% =====================================================================
%  LinkStudio - Theory and User Manual
%  Build:  latexmk -pdf linkstudio-manual.tex      (see README.md)
%  Data, code excerpts and reference tables come from tools/generate.py.
% =====================================================================
\documentclass[11pt,a4paper,twoside,openany]{book}
% =====================================================================
%  LinkStudio manual - preamble (packages, styles, macros)
% =====================================================================
\PassOptionsToPackage{obeyspaces,spaces}{url}
\usepackage[T1]{fontenc}
\usepackage[utf8]{inputenc}
\usepackage{lmodern}
\usepackage[a4paper,top=27mm,bottom=27mm,inner=28mm,outer=24mm,headheight=14pt]{geometry}
\usepackage{microtype}
\usepackage{amsmath,amssymb,amsthm,bm,mathtools}
\usepackage{siunitx}
\sisetup{per-mode=symbol,group-separator={\,},range-phrase=--}
\usepackage{booktabs,longtable,array,tabularx,multirow}
\usepackage{enumitem}
\setlist{itemsep=2pt,topsep=4pt}
\usepackage{xcolor}
\usepackage{graphicx}
\usepackage{caption,subcaption}
\captionsetup{font=small,labelfont=bf,format=hang,justification=raggedright,singlelinecheck=true}
\usepackage{fancyhdr}
\usepackage{titlesec}
\usepackage{tikz}
\usetikzlibrary{arrows.meta,positioning,calc,shapes.geometric,shapes.misc,fit,backgrounds,
  decorations.pathreplacing,decorations.markings,matrix,patterns}
\usepackage{pgfplots}
\pgfplotsset{compat=1.18}
\usepgfplotslibrary{groupplots,fillbetween}
\usepackage{listings}
\usepackage[most]{tcolorbox}
\usepackage[hidelinks]{hyperref}
\usepackage[nameinlink,capitalise,noabbrev]{cleveref}

% ------------------------------------------------------------- colours
\definecolor{ink}{HTML}{1F2933}
\definecolor{accent}{HTML}{1D4E89}
\definecolor{accent2}{HTML}{C8553D}
\definecolor{accent3}{HTML}{2A9D8F}
\definecolor{accent4}{HTML}{8E5DB5}
\definecolor{accent5}{HTML}{E9A23B}
\definecolor{soft}{HTML}{EEF3F9}
\definecolor{softwarm}{HTML}{FBF1EC}
\definecolor{softgreen}{HTML}{EAF6F4}
\definecolor{codebg}{HTML}{F7F8FA}
\definecolor{codeframe}{HTML}{D5DBE3}
\definecolor{kwcolor}{HTML}{1D4E89}
\definecolor{strcolor}{HTML}{A23E48}
\definecolor{commentcolor}{HTML}{6B7785}
% signal-kind colours (same as the GUI)
\definecolor{kbits}{HTML}{8E5DB5}
\definecolor{kllr}{HTML}{D6589F}
\definecolor{ksym}{HTML}{2A9D8F}
\definecolor{ksamp}{HTML}{1D6FB8}
\definecolor{kgrid}{HTML}{E9A23B}
\definecolor{kcsi}{HTML}{C8553D}

% ------------------------------------------------------------- headings
\titleformat{\chapter}[display]
  {\normalfont\sffamily\color{accent}}
  {\large\scshape Chapter \thechapter}{4pt}
  {\Huge\bfseries}[\vspace{2pt}{\color{accent}\titlerule[0.8pt]}]
\titlespacing*{\chapter}{0pt}{-10pt}{24pt}
\titleformat{\part}[display]
  {\normalfont\sffamily\centering\color{accent}}
  {\Large\scshape Part \thepart}{8pt}
  {\Huge\bfseries}
\titleformat{\section}{\normalfont\sffamily\Large\bfseries\color{ink}}{\color{accent}\thesection}{0.8em}{}
\titleformat{\subsection}{\normalfont\sffamily\large\bfseries\color{ink}}{\color{accent}\thesubsection}{0.8em}{}
\titleformat{\subsubsection}{\normalfont\sffamily\normalsize\bfseries\color{ink}}{}{0em}{}

\pagestyle{fancy}
\fancyhf{}
\fancyhead[LE]{\sffamily\small\color{commentcolor}\leftmark}
\fancyhead[RO]{\sffamily\small\color{commentcolor}\rightmark}
\fancyfoot[C]{\sffamily\small\thepage}
\renewcommand{\headrulewidth}{0.3pt}

% ------------------------------------------------------------- theorems
\theoremstyle{definition}
\newtheorem{definition}{Definition}[chapter]
\newtheorem{proposition}{Proposition}[chapter]
\newtheorem{example}{Example}[chapter]
\theoremstyle{remark}
\newtheorem*{remark}{Remark}

% ------------------------------------------------------------- code
\lstdefinestyle{py}{
  language=Python,
  basicstyle=\ttfamily\footnotesize\color{ink},
  keywordstyle=\color{kwcolor}\bfseries,
  stringstyle=\color{strcolor},
  commentstyle=\color{commentcolor}\itshape,
  numbers=left, numberstyle=\tiny\color{commentcolor}, numbersep=6pt,
  backgroundcolor=\color{codebg},
  frame=single, rulecolor=\color{codeframe}, framesep=3pt,
  xleftmargin=16pt, framexleftmargin=14pt,
  breaklines=true, breakatwhitespace=false, columns=fullflexible, keepspaces=true,
  showstringspaces=false, tabsize=4, upquote=true,
  captionpos=t, abovecaptionskip=2pt, belowcaptionskip=4pt,
  morekeywords={self,None,True,False},
  literate={π}{{$\pi$}}1 {β}{{$\beta$}}1 {—}{{--}}1
}
\lstdefinestyle{yaml}{
  basicstyle=\ttfamily\footnotesize\color{ink},
  backgroundcolor=\color{codebg}, frame=single, rulecolor=\color{codeframe},
  breaklines=true, columns=fullflexible, keepspaces=true, showstringspaces=false,
  commentstyle=\color{commentcolor}\itshape, morecomment=[l]{\#},
  captionpos=t
}
\lstdefinestyle{shell}{
  basicstyle=\ttfamily\footnotesize\color{ink},
  backgroundcolor=\color{codebg}, frame=single, rulecolor=\color{codeframe},
  columns=fullflexible, keepspaces=true, showstringspaces=false, captionpos=t
}
\lstset{style=py}
\renewcommand{\lstlistingname}{Code}
\renewcommand{\lstlistlistingname}{List of code excerpts}
\crefname{lstlisting}{Code}{Code}
\Crefname{lstlisting}{Code}{Code}

% Excerpts are extracted from the code base by tools/generate.py.
% \snipdef{name}{file}{first}{last} is written to generated/snippets.tex.
\makeatletter
\newcommand{\snipdef}[4]{%
  \expandafter\xdef\csname snipfile@#1\endcsname{\detokenize{#2}}%
  \expandafter\gdef\csname snipfirst@#1\endcsname{#3}%
  \expandafter\gdef\csname sniplast@#1\endcsname{#4}}
\newcommand{\snipsource}[1]{%
  \texttt{\csname snipfile@#1\endcsname}, lines
  \csname snipfirst@#1\endcsname--\csname sniplast@#1\endcsname}
% \codeexcerpt{name}{caption}
\newcommand{\codeexcerpt}[2]{%
  \edef\snip@first{\csname snipfirst@#1\endcsname}%
  \lstinputlisting[style=py,firstnumber=\snip@first,label={lst:#1},
    caption={#2\hfill{\normalfont\footnotesize\color{commentcolor}\snipsource{#1}}}]
    {generated/snippets/#1.py}}
\makeatother
% Generated by manual/tools/generate.py: \snipdef{name}{file}{first line}{last line}
\snipdef{block_reset}{linkstudio/core/block.py}{102}{114}
\snipdef{graph_toposort}{linkstudio/core/graph.py}{193}{215}
\snipdef{graph_execute}{linkstudio/core/graph.py}{240}{265}
\snipdef{stop_check}{linkstudio/core/simulation.py}{35}{47}
\snipdef{parallel_merge}{linkstudio/core/parallel.py}{140}{145}
\snipdef{sweep_point}{linkstudio/core/sweep.py}{140}{146}
\snipdef{awgn_noise_var}{linkstudio/phy/channel/awgn.py}{49}{72}
\snipdef{awgn_process}{linkstudio/phy/channel/awgn.py}{74}{82}
\snipdef{wilson}{linkstudio/phy/metrics.py}{13}{21}
\snipdef{demap_soft}{linkstudio/phy/dsp/constellations.py}{72}{86}
\snipdef{qam_build}{linkstudio/phy/dsp/constellations.py}{118}{125}
\snipdef{rrc_taps}{linkstudio/phy/dsp/filters.py}{19}{36}
\snipdef{mf_normalize}{linkstudio/phy/waveform/pulse_shaping.py}{96}{106}
\snipdef{pulse_shaper}{linkstudio/phy/waveform/pulse_shaping.py}{46}{53}
\snipdef{phase_noise}{linkstudio/phy/channel/impairments.py}{74}{79}
\snipdef{iq_imbalance}{linkstudio/phy/channel/impairments.py}{117}{124}
\snipdef{power_amplifier}{linkstudio/phy/channel/impairments.py}{150}{163}
\snipdef{quantizer}{linkstudio/phy/channel/impairments.py}{191}{198}
\snipdef{sos}{linkstudio/phy/channel/fading.py}{68}{86}
\snipdef{flat_fading}{linkstudio/phy/channel/fading.py}{115}{126}
\snipdef{linear_eq}{linkstudio/phy/receiver.py}{87}{108}
\snipdef{lms_loop}{linkstudio/phy/receiver.py}{147}{166}
\snipdef{loop_gains}{linkstudio/phy/sync.py}{24}{31}
\snipdef{cubic}{linkstudio/phy/sync.py}{36}{44}
\snipdef{om_estimate}{linkstudio/phy/sync.py}{117}{123}
\snipdef{gardner_loop}{linkstudio/phy/sync.py}{128}{136}
\snipdef{frame_sync}{linkstudio/phy/sync.py}{191}{212}
\snipdef{dd_pll}{linkstudio/phy/sync.py}{246}{252}
\snipdef{ofdm_mod}{linkstudio/phy/waveform/ofdm.py}{94}{108}
\snipdef{ofdm_sync}{linkstudio/phy/waveform/ofdm.py}{148}{159}
\snipdef{ls_slope}{linkstudio/phy/waveform/ofdm.py}{235}{251}
\snipdef{mimo_linear}{linkstudio/phy/mimo.py}{152}{161}
\snipdef{alamouti_dec}{linkstudio/phy/mimo.py}{196}{200}
\snipdef{crc_matrix}{linkstudio/phy/coding/crc.py}{30}{39}
\snipdef{lfsr}{linkstudio/phy/dsp/lfsr.py}{27}{40}
\snipdef{viterbi}{linkstudio/phy/coding/convolutional.py}{67}{86}
\snipdef{ldpc_encode}{linkstudio/phy/coding/ldpc.py}{71}{74}
\snipdef{ldpc_minsum}{linkstudio/phy/coding/ldpc.py}{92}{103}
\snipdef{polar_rel}{linkstudio/phy/coding/polar.py}{35}{39}
\snipdef{polar_encode}{linkstudio/phy/coding/polar.py}{44}{52}
\snipdef{polar_sc}{linkstudio/phy/coding/polar.py}{59}{83}
\snipdef{depuncture}{linkstudio/phy/coding/puncture.py}{55}{64}
\snipdef{cable}{linkstudio/phy/wireline.py}{296}{305}
\snipdef{xtalk}{linkstudio/phy/wireline.py}{357}{362}
\snipdef{echo_nlms}{linkstudio/phy/wireline.py}{479}{484}
\snipdef{dfe_loop}{linkstudio/phy/wireline.py}{565}{578}
\snipdef{waterfill}{linkstudio/phy/wireline.py}{616}{619}
\snipdef{dmt_hermitian}{linkstudio/phy/wireline.py}{686}{692}
\snipdef{jitter}{linkstudio/phy/wireline.py}{602}{606}
\snipdef{wiener}{linkstudio/phy/optical.py}{51}{53}
\snipdef{im_mod}{linkstudio/phy/optical.py}{172}{176}
\snipdef{iq_mzm}{linkstudio/phy/optical.py}{280}{289}
\snipdef{disp_op}{linkstudio/phy/optical.py}{313}{324}
\snipdef{beta2}{linkstudio/phy/optical.py}{298}{302}
\snipdef{ssfm}{linkstudio/phy/optical.py}{333}{373}
\snipdef{pmd_transfer}{linkstudio/phy/optical.py}{426}{436}
\snipdef{ase}{linkstudio/phy/optical.py}{521}{530}
\snipdef{add_ase}{linkstudio/phy/optical.py}{533}{535}
\snipdef{osnr}{linkstudio/phy/optical.py}{68}{74}
\snipdef{link_loop}{linkstudio/phy/optical.py}{687}{696}
\snipdef{photodiode}{linkstudio/phy/optical.py}{1005}{1010}
\snipdef{coherent_rx}{linkstudio/phy/optical.py}{1054}{1073}
\snipdef{cd_comp}{linkstudio/phy/optical.py}{1140}{1141}
\snipdef{fo_est}{linkstudio/phy/optical.py}{1180}{1187}
\snipdef{butterfly}{linkstudio/phy/optical.py}{1251}{1265}
\snipdef{bps}{linkstudio/phy/optical.py}{1306}{1319}
\snipdef{wdm_demux}{linkstudio/phy/optical.py}{794}{800}
\snipdef{fso}{linkstudio/phy/optical.py}{866}{869}
\snipdef{fso_params}{linkstudio/phy/optical.py}{814}{817}
\snipdef{rin}{linkstudio/phy/optical.py}{59}{61}
\snipdef{mmf_delay}{linkstudio/phy/multimode.py}{63}{66}
\snipdef{mmf_rms}{linkstudio/phy/multimode.py}{74}{80}
\snipdef{mmf_chromatic}{linkstudio/phy/multimode.py}{94}{96}
\snipdef{mmf_groups}{linkstudio/phy/multimode.py}{194}{208}
\snipdef{mmf_launch}{linkstudio/phy/multimode.py}{108}{125}
\snipdef{mmf_response}{linkstudio/phy/multimode.py}{237}{243}
\snipdef{mmf_default_alpha}{linkstudio/phy/multimode.py}{184}{190}
\snipdef{laser_rhs}{linkstudio/phy/lasers.py}{142}{149}
\snipdef{laser_steady}{linkstudio/phy/lasers.py}{105}{109}
\snipdef{laser_small_signal}{linkstudio/phy/lasers.py}{236}{243}
\snipdef{laser_overlap}{linkstudio/phy/lasers.py}{90}{99}
\snipdef{vcsel_thermal}{linkstudio/phy/lasers.py}{280}{292}
\snipdef{vcsel_vi}{linkstudio/phy/lasers.py}{267}{270}
\snipdef{laser_block}{linkstudio/phy/lasers.py}{337}{345}
\snipdef{mzm_field}{linkstudio/phy/optical.py}{128}{132}
\snipdef{mzm_block}{linkstudio/phy/optical.py}{229}{233}
\snipdef{raman}{linkstudio/phy/optical.py}{411}{415}
\snipdef{amp_sat}{linkstudio/phy/optical.py}{541}{543}
\snipdef{link_gn}{linkstudio/phy/optical.py}{709}{716}
\snipdef{gn_eta}{linkstudio/phy/theory.py}{233}{240}
\snipdef{osnr_link}{linkstudio/phy/theory.py}{202}{207}
\snipdef{pam_ser}{linkstudio/phy/theory.py}{148}{150}
\snipdef{coh_isnr}{linkstudio/phy/optical.py}{1067}{1072}
\snipdef{delay_kernel}{linkstudio/phy/channel/fading.py}{51}{60}
\snipdef{tr38901}{linkstudio/phy/channel/propagation.py}{35}{56}
\snipdef{dielectric}{linkstudio/phy/wireline.py}{268}{281}
\snipdef{pic}{linkstudio/phy/wireline.py}{248}{257}
\snipdef{fso_channel}{linkstudio/phy/optical.py}{915}{935}

\def\numblocks{90}

\def\numscenarios{42}

\def\dmtrateA{5.2}
\def\dmtrateB{1.3}
\def\dmtrateC{0.3}

\def\fiberLD{1.153}
\def\fiberbeta{21.68}

\def\solitonPzero{0.667}
\def\solitonL{4.61}

\newcommand{\mmfalphaopt}{1.9777}
\newcommand{\mmfalphaomfour}{1.9971}
\newcommand{\mmfalphaomthree}{2.0235}
\newcommand{\mmfdelta}{0.0093}
\newcommand{\mmfgroups}{18}
\newcommand{\mmfsigmamin}{12.2}


% ------------------------------------------------------------- boxes
\newtcolorbox{incode}[1][]{enhanced,breakable,colback=soft,colframe=accent,boxrule=0pt,
  leftrule=2.5pt,arc=0pt,outer arc=0pt,left=8pt,right=8pt,top=5pt,bottom=5pt,
  fonttitle=\sffamily\bfseries\small,coltitle=accent,colbacktitle=soft,
  title={In the code#1},attach title to upper={\quad},before skip=8pt,after skip=8pt}
\newtcolorbox{keypoint}[1][Key idea]{enhanced,breakable,colback=softgreen,colframe=accent3,boxrule=0pt,
  leftrule=2.5pt,arc=0pt,outer arc=0pt,left=8pt,right=8pt,top=5pt,bottom=5pt,
  fonttitle=\sffamily\bfseries\small,coltitle=accent3,colbacktitle=softgreen,
  title={#1},attach title to upper={\quad},before skip=8pt,after skip=8pt}
\newtcolorbox{pitfall}[1][Watch out]{enhanced,breakable,colback=softwarm,colframe=accent2,boxrule=0pt,
  leftrule=2.5pt,arc=0pt,outer arc=0pt,left=8pt,right=8pt,top=5pt,bottom=5pt,
  fonttitle=\sffamily\bfseries\small,coltitle=accent2,colbacktitle=softwarm,
  title={#1},attach title to upper={\quad},before skip=8pt,after skip=8pt}

% ------------------------------------------------------------- inline names
% Inline code names use the url machinery on detokenized text: special characters print
% verbatim, long names can break after / _ . = , and the commands are robust (usable in
% captions). Braces must be balanced.
\makeatletter
\DeclareUrlCommand\cs@tt{\urlstyle{tt}}
\g@addto@macro\UrlBreaks{\do\_\do\.\do\=\do\,}
\DeclareRobustCommand{\code}[1]{\expandafter\cs@tt\expandafter{\detokenize{#1}}}
\DeclareRobustCommand{\file}[1]{\expandafter\cs@tt\expandafter{\detokenize{#1}}}
\DeclareRobustCommand{\blk}[1]{\textcolor{accent}{\code{#1}}}          % block type
\DeclareRobustCommand{\prm}[1]{\textcolor{accent2!85!black}{\code{#1}}}% parameter
\DeclareRobustCommand{\meta}[1]{\textcolor{accent4}{\code{#1}}}        % metadata key
\makeatother
\newcommand{\kind}[1]{\textsf{\small #1}}

% ------------------------------------------------------------- maths
\newcommand{\E}{\mathbb{E}}
\newcommand{\R}{\mathbb{R}}
\newcommand{\C}{\mathbb{C}}
\newcommand{\CN}{\mathcal{CN}}
\newcommand{\N}{\mathcal{N}}
\newcommand{\Q}{\operatorname{Q}}
\newcommand{\erfc}{\operatorname{erfc}}
\newcommand{\sinc}{\operatorname{sinc}}
\newcommand{\diag}{\operatorname{diag}}
\newcommand{\sgn}{\operatorname{sgn}}
\newcommand{\Var}{\operatorname{Var}}
\newcommand{\F}{\mathcal{F}}
\newcommand{\EsNo}{E_s/N_0}
\newcommand{\EbNo}{E_b/N_0}
\newcommand{\dB}{\,\mathrm{dB}}
\newcommand{\T}{^{\mathsf{T}}}
\newcommand{\Her}{^{\mathsf{H}}}
\newcommand{\conj}{^{*}}
\newcommand{\vect}[1]{\bm{#1}}
\newcommand{\mat}[1]{\mathbf{#1}}
\DeclareMathOperator*{\argmin}{arg\,min}
\DeclareMathOperator*{\argmax}{arg\,max}

% ------------------------------------------------------------- TikZ styles
\tikzset{
  >={Latex[length=2mm]},
  blockbox/.style={draw=accent,fill=soft,rounded corners=2pt,minimum height=8mm,minimum width=16mm,
    align=center,font=\sffamily\footnotesize,inner sep=3pt},
  warmbox/.style={blockbox,draw=accent2,fill=softwarm},
  greenbox/.style={blockbox,draw=accent3,fill=softgreen},
  greybox/.style={blockbox,draw=black!40,fill=black!4},
  sumnode/.style={draw=accent,fill=white,circle,inner sep=0pt,minimum size=5mm,font=\footnotesize},
  mulnode/.style={draw=accent,fill=white,circle,inner sep=0pt,minimum size=5mm,font=\footnotesize,
    path picture={\draw[accent] (path picture bounding box.south east)--(path picture bounding box.north west)
    (path picture bounding box.south west)--(path picture bounding box.north east);}},
  sig/.style={->,thick,color=ink},
  lbl/.style={font=\scriptsize\sffamily,color=ink},
  note/.style={font=\scriptsize\itshape,color=commentcolor,align=center},
  layer/.style={draw=#1,fill=#1!7,rounded corners=4pt,inner sep=6pt},
  dot/.style={circle,fill=ink,inner sep=0pt,minimum size=3pt},
}

% ------------------------------------------------------------- pgfplots styles
\pgfplotscreateplotcyclelist{linkstudio}{
  {accent,thick},{accent2,thick},{accent3,thick},{accent4,thick},{accent5,thick},{black!60,thick}}
\pgfplotsset{
  every axis/.append style={font=\footnotesize,
    grid=both, major grid style={black!14}, minor grid style={black!6},
    legend style={font=\scriptsize,draw=black!20,fill=white,fill opacity=0.9,text opacity=1},
    tick label style={font=\scriptsize}, label style={font=\footnotesize},
    cycle list name=linkstudio, unbounded coords=jump, axis line style={black!60}},
  sim/.style={only marks,mark=*,mark size=1.6pt,#1},
  th/.style={no marks,thick,#1},
  simline/.style={thick,mark=*,mark size=1.5pt,#1},
  scatterpts/.style={only marks,mark=*,mark size=0.35pt,#1},
  berplot/.style={xlabel={$E_b/N_0$ [dB]},ylabel={BER},
    yticklabel style={font=\scriptsize}, legend pos=south west},
  constaxis/.style={width=4.4cm,height=4.4cm,axis equal image,xmin=-1.6,xmax=1.6,ymin=-1.6,ymax=1.6,
    xtick={-1,0,1},ytick={-1,0,1},grid=major,title style={font=\scriptsize\sffamily}},
}
\newcommand{\datafile}[1]{generated/data/#1.dat}

% ------------------------------------------------------------- misc
\setcounter{tocdepth}{1}
\setcounter{secnumdepth}{2}
\renewcommand{\arraystretch}{1.15}
\hypersetup{pdftitle={LinkStudio - Theory and User Manual},pdfauthor={LinkStudio project}}


\begin{document}
\frontmatter
\begin{titlepage}
  \centering
  \vspace*{30mm}
  {\sffamily\color{accent}\fontsize{40}{44}\selectfont\bfseries LinkStudio\par}
  \vspace{6mm}
  {\sffamily\Large\color{ink} A modular physical-layer communication simulator\par}
  \vspace{4mm}
  {\color{accent}\rule{0.6\textwidth}{0.8pt}\par}
  \vspace{6mm}
  {\sffamily\LARGE Theory and User Manual\par}
  \vspace{22mm}
  \begin{tikzpicture}[scale=0.9,transform shape]
    \node[blockbox,fill=kbits!12,draw=kbits] (a) {bits};
    \node[blockbox,right=6mm of a,fill=ksym!12,draw=ksym] (b) {symbols};
    \node[blockbox,right=6mm of b,fill=ksamp!12,draw=ksamp] (c) {waveform};
    \node[warmbox,right=6mm of c] (d) {channel};
    \node[blockbox,right=6mm of d,fill=ksamp!12,draw=ksamp] (e) {receiver};
    \node[blockbox,right=6mm of e,fill=kllr!12,draw=kllr] (f) {LLRs};
    \foreach \x/\y in {a/b,b/c,c/d,d/e,e/f} \draw[sig] (\x)--(\y);
  \end{tikzpicture}
  \vfill
  {\sffamily\small Wireless \textperiodcentered{} Wireline \textperiodcentered{} Optical \textperiodcentered{} Coding \textperiodcentered{} Synchronization \textperiodcentered{} MIMO\par}
  \vspace{4mm}
  {\sffamily\small Engine v0.4 \textperiodcentered{} \today\par}
\end{titlepage}

\cleardoublepage
\tableofcontents
\chapter{About this manual}

LinkStudio is a block-based simulator of the physical layer (PHY) of communication systems. A
simulation is a directed graph of blocks: sources, modulators, pulse shapers, channels,
synchronizers, equalizers, decoders and meters. The same engine covers radio links,
copper lines and optical fiber. It is driven from Python, from YAML scenario files, from the
command line or from a browser-based editor.

This manual has two goals and is organized accordingly.

\begin{description}[leftmargin=0pt,font=\sffamily\bfseries\color{accent}]
  \item[Part I -- Using and extending LinkStudio.] How the engine works (signals, blocks, graphs,
  Monte-Carlo runs, sweeps, scenarios) and how to drive it from the GUI, the CLI and the
  Python API. It ends with a developer guide for writing new blocks.
  \item[Part II -- Theory and implementation.] One chapter per subsystem. Every chapter first
  develops the theory with the notation used in the code, then shows \emph{where} and
  \emph{how} it is implemented, and finally validates the implementation against a closed-form
  result. Part II can be read as a compact textbook on digital transmission whose equations
  can all be checked in the code.
\end{description}

\section*{How to read the chapters of Part II}
Each theory chapter follows the same pattern:
\begin{enumerate}
  \item \textbf{Model.} The physical or mathematical model, derived with the assumptions stated
  explicitly.
  \item \textbf{Implementation.} A blue \emph{In the code} box names the blocks, files and
  metadata involved. Short excerpts show the lines that realize the key equations, for example
  the split-step solver of the nonlinear Schrödinger equation (\cref{lst:ssfm}). Excerpts are
  extracted automatically from the source, so their file names and line numbers match the code
  base this manual was built from.
  \item \textbf{Validation.} A figure comparing the simulator with theory. All measured curves
  in this manual (markers in the plots) are produced by running the real blocks with
  \file{manual/tools/generate.py}; solid lines are closed-form results.
\end{enumerate}

Green \emph{Key idea} boxes summarize a result; orange \emph{Watch out} boxes point at
common mistakes or at modelling limitations.

\section*{Typographic conventions}
\begin{center}
\begin{tabular}{@{}ll@{}}
\toprule
\blk{awgn} & a block type (as written in scenarios and in the GUI palette)\\
\prm{snr_db} & a block parameter\\
\meta{noise_var} & a metadata key travelling with a signal\\
\file{linkstudio/phy/optical.py} & a file, relative to \file{backend/}\\
\kind{SAMPLES} & a signal kind (\cref{sec:signal})\\
\bottomrule
\end{tabular}
\end{center}

\section*{Notation}
\begin{center}
\begin{tabular}{@{}ll@{\qquad}ll@{}}
\toprule
$x(t)$, $x[n]$ & complex envelope (continuous, sampled) & $f_s$ & sample rate\\
$T$, $R_s=1/T$ & symbol period, symbol rate & $N_{\mathrm{sps}}$ & samples per symbol\\
$E_s$, $E_b$ & energy per symbol, per information bit & $N_0$ & one-sided noise PSD\\
$M$, $k=\log_2 M$ & constellation size, bits per symbol & $R$ & code rate\\
$\CN(\mu,\sigma^2)$ & circular complex Gaussian & $\Q(x)$ & Gaussian tail function\\
$(\cdot)\conj$, $(\cdot)\T$, $(\cdot)\Her$ & conjugate, transpose, Hermitian & $\star$ & convolution\\
$\F\{\cdot\}$ & Fourier transform & $\E[\cdot]$ & expectation\\
\bottomrule
\end{tabular}
\end{center}
Fourier transforms follow the NumPy sign convention used everywhere in the code:
\begin{equation}
  X(\omega)=\int x(t)\,e^{-j\omega t}\,dt, \qquad x(t)=\frac{1}{2\pi}\int X(\omega)\,e^{j\omega t}\,d\omega .
  \label{eq:fourier-convention}
\end{equation}
This convention fixes the signs of the dispersion operator and of the nonlinear phase in the
fiber model (\cref{sec:nlse}). The discrete transforms are \code{np.fft.fft} and
\code{np.fft.ifft}; they are unitary (\code{norm="ortho"}) in OFDM and DMT.

\section*{Rebuilding this manual}
\begin{lstlisting}[style=shell]
python manual/tools/generate.py      # run the simulator, extract excerpts and tables
cd manual && latexmk -pdf linkstudio-manual.tex
\end{lstlisting}
Re-running the generator after a code change updates every figure, excerpt and reference
table, so the manual stays consistent with the code.


\mainmatter
\part{Using and extending LinkStudio}
\chapter{Overview and quick start}\label{ch:overview}

\section{What LinkStudio is}
LinkStudio simulates the physical layer of a communication link at the level of
\emph{complex-envelope samples}. A link is described as a graph of blocks. Every block is a
small, self-contained model (a QAM mapper, an 80-km fiber span, a Viterbi decoder, a BER
meter) with typed inputs and outputs and validated parameters. The engine executes the
graph on \emph{batches} of independent frames and accumulates metrics until a stop criterion
is met.

\Cref{fig:architecture} shows the layers of the code base. They depend on each other only
downwards:
\begin{itemize}
  \item the \textbf{core} knows nothing about communications;
  \item the \textbf{PHY library} knows nothing about the user interface;
  \item the \textbf{API} turns graphs and results into JSON for the browser front end.
\end{itemize}
The same scenario therefore runs identically from a Python script, the CLI or the GUI.

\begin{figure}[htbp]
\centering
\begin{tikzpicture}[node distance=3mm,
  lay/.style={draw=#1,fill=#1!6,rounded corners=3pt,minimum width=12.0cm,minimum height=13mm,
    align=left,font=\sffamily\footnotesize,inner sep=5pt,text width=11.6cm}]
  \node[lay=accent4] (ui) {\textbf{Front end} \hfill\file{frontend/src}\\
    React + TypeScript: graph editor (React Flow), inspector, Live / Simulate / Sweep panels, scopes};
  \node[lay=accent5,below=of ui] (api) {\textbf{API} \hfill\file{linkstudio/api}\\
    FastAPI: REST (blocks, scenarios, validate, layout, export, jobs), WebSocket live mode, probes};
  \node[lay=accent3,below=of api] (phy) {\textbf{PHY library} \hfill\file{linkstudio/phy}\\
    \numblocks{} registered blocks: sources, modulation, waveform, channel, sync, receiver, MIMO,
    coding, wireline, optical, metrics; closed-form theory};
  \node[lay=accent,below=of phy] (core) {\textbf{Core engine} \hfill\file{linkstudio/core}\\
    Signal, Block, Port, registry, Graph (validation, execution), Subsystem, Simulation,
    Sweep, Scenario (YAML)};
  \node[lay=gray,below=of core] (np) {\textbf{Numerics}\\ NumPy, SciPy (FFT, special
    functions, sparse matrices), Pydantic (parameter validation), PyYAML};
  \node[font=\sffamily\scriptsize,color=commentcolor,right=2mm of ui.east,anchor=west,rotate=90,
    yshift=-3mm] at ($(ui.east)!0.5!(np.east)$) {depends on $\downarrow$};
  \draw[->,thick,commentcolor] ($(ui.east)+(0.25,0)$) -- ($(np.east)+(0.25,0)$);
  \node[greybox,left=4mm of api,minimum width=18mm] (cli) {CLI\\\file{linkstudio.cli}};
  \node[greybox,left=4mm of phy,minimum width=18mm] (scripts) {scripts\\\file{scripts/}};
  \node[greybox,left=4mm of core,minimum width=18mm] (tests) {tests\\\file{tests/}};
  \draw[->,commentcolor] (cli.east) -- (api.west);
  \draw[->,commentcolor] (scripts.east) -- (phy.west);
  \draw[->,commentcolor] (tests.east) -- (core.west);
\end{tikzpicture}
\caption{Layers of LinkStudio. Paths are relative to the repository; Python packages live in
\file{backend/}.}
\label{fig:architecture}
\end{figure}

\subsection{Capabilities at a glance}
\begin{center}
\small
\begin{tabularx}{\textwidth}{@{}lX@{}}
\toprule
\textbf{Domain} & \textbf{Models} \\
\midrule
Modulation & PSK, square QAM, PAM, custom constellations with Gray labels; hard, exact and max-log soft demapping\\
Waveforms & RRC/RC/Gaussian/rectangular pulses, matched filter, rate changes, OFDM with pilots and CP, DMT with water-filling\\
Channels & AWGN (SNR, $E_s/N_0$, $E_b/N_0$), Rayleigh/Rician flat fading, Jakes Doppler, 3GPP EPA/EVA/ETU, MIMO with antenna correlation\\
RF front end & carrier offset, phase noise, IQ imbalance, PA models (Rapp, Saleh, clipping), quantizer, timing offset\\
Synchronization & Schmidl--Cox frame sync, Gardner timing loop with Oerder--Meyr start, decision-directed PLL, CP-based OFDM sync\\
Receivers & one-tap, ZF/MMSE FIR, LMS/NLMS, FFE+DFE, CTLE, echo canceller, ZF/MMSE/ML MIMO detection, Alamouti\\
Coding & CRC, scrambler, interleavers, repetition, Hamming, convolutional + Viterbi, puncturing, QC-LDPC (BP), polar (SC)\\
Wireline & NRZ, Manchester, AMI, MLT-3, 4B/5B, 64B/66B, skin-effect cables, ANSI NEXT/FEXT, impulse noise, jitter, ADSL/VDSL DMT\\
Optical & laser/EML with chirp and RIN, IQ Mach--Zehnder, NLSE/Manakov fiber with PMD, EDFA with ASE, WDM, PIN/APD, coherent receiver and DSP, FSO turbulence\\
Metrics & BER with confidence interval, SER, FER/BLER, EVM/MER, power/PAPR, OSNR, iteration counts\\
\bottomrule
\end{tabularx}
\end{center}

\section{Installation}
LinkStudio needs Python 3.11 or newer. The front end is already built in
\file{frontend/dist}; Node.js is only needed to modify it.
\begin{lstlisting}[style=shell]
cd backend
pip install -e ".[ui,viz,dev]"    # engine + web server + plotting + tests
python -m pytest                   # the test-suite (about 170 tests)
python scripts/main.py --quick     # runs every demo and compares with theory
\end{lstlisting}

\section{Three ways to run a simulation}\label{sec:three-ways}

\paragraph{1. The graphical editor.}
\code{python -m linkstudio serve} opens \url{http://127.0.0.1:8000}. Pick a scenario from the
browser (\cref{ch:interfaces}), press \emph{Start live} and move the
\prm{snr_db} slider of the \blk{awgn} block: constellation, spectrum and BER update while
the simulation runs.

\paragraph{2. A YAML scenario from the command line.}
Scenario files describe the graph and how to run it (\cref{sec:scenarios}):
\begin{lstlisting}[style=shell]
python -m linkstudio run scenarios/tutorials/01_bpsk_awgn.yaml
python -m linkstudio blocks              # list all block types
python -m linkstudio describe fiber      # parameters and ports of one block
\end{lstlisting}

\paragraph{3. The Python API.}
The snippet below builds the graph of \cref{fig:first-graph} and measures the BER of QPSK at
$E_b/N_0=6$\,dB. The result is about $2.4\times10^{-3}$, the value of
$\Q(\sqrt{2E_b/N_0})$.
\begin{lstlisting}[style=py,numbers=none]
import linkstudio as cs

g = cs.Graph("qpsk")
g.create("src", "random_bits", n_bits=2000)
g.create("map", "mapper", scheme="psk", order=4)
g.create("awgn", "awgn", snr_db=6.0)                  # Eb/N0 by default
g.create("demap", "demapper", method="hard")
g.create("ber", "ber_meter")
g.chain("src", "map", "awgn", "demap", "ber.rx")     # default ports
g.connect("src", "ber.tx")                            # explicit edge

result = cs.Simulation(g, batch_size=64, seed=1).run(
    stop=cs.StopCriteria.errors("ber.bit_errors", 200))
print(result.summary())
\end{lstlisting}

\begin{figure}[htbp]
\centering
\begin{tikzpicture}[node distance=7mm]
  \node[blockbox] (src) {\blk{src}\\random\_bits};
  \node[blockbox,right=of src] (map) {\blk{map}\\mapper};
  \node[warmbox,right=of map] (awgn) {\blk{awgn}\\awgn};
  \node[blockbox,right=of awgn] (dem) {\blk{demap}\\demapper};
  \node[greenbox,right=of dem] (ber) {\blk{ber}\\ber\_meter};
  \draw[sig,kbits] (src) -- node[above,lbl]{BITS} (map);
  \draw[sig,ksym] (map) -- node[above,lbl]{SYMBOLS} (awgn);
  \draw[sig,ksym] (awgn) -- (dem);
  \draw[sig,kbits] (dem) -- node[above,lbl]{rx} (ber);
  \draw[sig,kbits] (src.south) -- ++(0,-6mm) -| node[pos=0.25,below,lbl]{tx} (ber.south);
\end{tikzpicture}
\caption{The graph built by the quick-start snippet. Edge colours follow the signal kind.}
\label{fig:first-graph}
\end{figure}

\section{Repository map}
\begin{center}
\small
\begin{tabularx}{\textwidth}{@{}lX@{}}
\toprule
\file{backend/linkstudio/core/} & engine: \file{signal.py}, \file{block.py}, \file{registry.py}, \file{graph.py}, \file{composite.py}, \file{simulation.py}, \file{sweep.py}, \file{scenario.py}\\
\file{backend/linkstudio/phy/} & block library; \file{dsp/} holds shared numerical helpers (filters, constellations, LFSR, analysis), \file{theory.py} closed-form references\\
\file{backend/linkstudio/api/} & web server: \file{app.py}, \file{editor.py} (GUI $\leftrightarrow$ graph), \file{jobs.py}, \file{probes.py}\\
\file{backend/scenarios/} & ready-made scenarios grouped by category (\cref{app:scenarios})\\
\file{backend/scripts/} & \file{main.py} and one demo per subsystem, each checking theory\\
\file{backend/tests/} & pytest suite (engine, DSP, physics, API, every scenario)\\
\file{frontend/} & React application\\
\file{docs/} & short Markdown notes (architecture, engine, GUI)\\
\file{manual/} & this manual and its generator\\
\bottomrule
\end{tabularx}
\end{center}

\chapter{The simulation engine}\label{ch:engine}

The engine (\file{linkstudio/core}) is about one thousand lines of Python and contains no
communication theory. It provides five abstractions:
\begin{itemize}
  \item \emph{signals} that flow along edges;
  \item \emph{blocks} that transform them;
  \item \emph{graphs} that connect blocks;
  \item \emph{simulations} that run graphs batch after batch;
  \item \emph{sweeps} that repeat simulations over a parameter grid.
\end{itemize}
This chapter describes each of them and the rules they enforce.

% ------------------------------------------------------------------
\section{Signals}\label{sec:signal}
A \code{Signal} (\file{core/signal.py}) is an immutable record with four fields:
\begin{description}[font=\ttfamily\color{accent}]
  \item[data] a NumPy array. Axis 0 is always the \emph{batch} (independent frames simulated in
  parallel) and the last axis is time, bit index or subcarrier index. Multi-dimensional kinds
  add axes in between, as in \cref{fig:signal-shapes}.
  \item[kind] the semantic type of the data, a flag of the \code{Kind} enumeration
  (\cref{tab:kinds}).
  \item[fs] the sample or symbol rate in Hz, when it is meaningful.
  \item[meta] a dictionary of metadata that travels downstream with the data (\cref{sec:metadata}).
\end{description}

\begin{figure}[htbp]
\centering
\begin{tikzpicture}[x=1cm,y=1cm]
  % 2D
  \begin{scope}
    \foreach \r in {0,...,3} \foreach \c in {0,...,7}
      \draw[fill=ksym!10,draw=ksym!60] (\c*0.4,-\r*0.4) rectangle ++(0.4,0.4);
    \draw[->] (-0.2,0.6) -- node[above,lbl]{time (last axis)} (3.4,0.6);
    \draw[->] (-0.25,0.4) -- node[left,lbl,align=right]{batch\\(axis 0)} (-0.25,-1.2);
    \node[lbl,below] at (1.6,-1.35) {\code{(B, n)}: bits, symbols, samples};
  \end{scope}
  % 3D streams
  \begin{scope}[xshift=4.9cm]
    \foreach \p in {1,0} {
      \foreach \r in {0,...,3} \foreach \c in {0,...,6}
        \draw[fill=ksamp!10,draw=ksamp!60] (\c*0.4+\p*0.25,-\r*0.4+\p*0.25) rectangle ++(0.4,0.4);}
    \node[lbl] at (3.45,0.85) {stream / polarization};
    \node[lbl,below] at (1.6,-1.35) {\code{(B, P, n)}: MIMO, dual pol.};
  \end{scope}
  % grid
  \begin{scope}[xshift=9.6cm]
    \foreach \p in {2,1,0} {
      \foreach \r in {0,...,3} \foreach \c in {0,...,5}
        \draw[fill=kgrid!12,draw=kgrid!80] (\c*0.4+\p*0.22,-\r*0.4+\p*0.22) rectangle ++(0.4,0.4);}
    \node[lbl,align=center] at (1.5,1.15) {OFDM symbols $\times$ subcarriers};
    \node[lbl,below] at (1.4,-1.35) {\code{(B, S, N_FFT)}: resource grid};
  \end{scope}
\end{tikzpicture}
\caption{Array conventions. Every block may assume that axis 0 is the batch and the last axis is
time (or frequency for grids).}
\label{fig:signal-shapes}
\end{figure}

\begin{table}[htbp]
\centering\small
\caption{Signal kinds. Kinds are bit flags, so a port can accept several of them
(\code{Kind.BITS | Kind.LLR}). The colours are those used by the GUI.}
\label{tab:kinds}
\begin{tabular}{@{}lllp{6.8cm}@{}}
\toprule
Kind & & dtype & Meaning\\
\midrule
\kind{BITS}    & \textcolor{kbits}{\rule{8pt}{8pt}} & \code{uint8} & hard bits in $\{0,1\}$\\
\kind{LLR}     & \textcolor{kllr}{\rule{8pt}{8pt}}  & \code{float64} & log-likelihood ratios, $L=\ln P(b{=}0)/P(b{=}1)$ (positive means 0)\\
\kind{SYMBOLS} & \textcolor{ksym}{\rule{8pt}{8pt}}  & \code{complex128} & one sample per symbol (constellation domain)\\
\kind{SAMPLES} & \textcolor{ksamp}{\rule{8pt}{8pt}} & \code{complex128}/\code{float64} & oversampled waveform; optical fields in $\sqrt{\mathrm{W}}$\\
\kind{GRID}    & \textcolor{kgrid}{\rule{8pt}{8pt}} & \code{complex128} & OFDM resource grid (frequency domain)\\
\kind{CSI}     & \textcolor{kcsi}{\rule{8pt}{8pt}}  & \code{complex128} & channel state information; the form is given by \meta{csi_type}\\
\kind{FLAGS}   & \textcolor{black!50}{\rule{8pt}{8pt}} & \code{bool} & one flag per frame (e.g. CRC check)\\
\bottomrule
\end{tabular}
\end{table}

Signals are never modified in place. A block derives a new signal with
\code{x.evolve(new_data, key=value)}. This keeps the input's kind, rate and metadata and
merges the new keys, so information added upstream is preserved unless a block deliberately
overrides it.

% ------------------------------------------------------------------
\section{Blocks}\label{sec:blocks}
A block (\file{core/block.py}) declares three things:
\begin{itemize}
  \item a \code{Params} model, a Pydantic class whose fields carry types, defaults, bounds
  and descriptions;
  \item named, typed \code{inputs} and \code{outputs} (\code{Port} objects);
  \item a \code{process(inputs, ctx)} method mapping one batch of input signals to output
  signals.
\end{itemize}
Measurement blocks also override \code{metrics()}.

Ports carry the accepted (or produced) kinds, a \code{required} flag for inputs and, for
outputs, an optional \code{same_as} reference: \code{Port.like("in")} produces the same kind
as input \code{in}. The GUI uses this information to colour handles and to refuse invalid
connections, and the graph validator uses it to propagate kinds (\cref{sec:validation}).

\subsection{Life cycle}
\begin{figure}[htbp]
\centering
\begin{tikzpicture}[node distance=6mm and 8mm]
  \node[blockbox] (init) {\code{__init__}\\validate Params};
  \node[blockbox,right=of init] (setup) {\code{setup()}\\derived state};
  \node[blockbox,right=of setup] (reset) {\code{reset(seed, path)}\\RNG + \code{on_reset()}};
  \node[greenbox,right=of reset] (proc) {\code{process()}\\one batch};
  \node[blockbox,below=9mm of proc] (met) {\code{metrics()}};
  \node[warmbox,below=9mm of setup] (setp) {\code{set_param()}};
  \draw[sig] (init) -- (setup);
  \draw[sig] (setup) -- (reset);
  \draw[sig] (reset) -- (proc);
  \draw[sig] (proc.north east) to[out=60,in=0,looseness=4] node[right,lbl]{next batch} (proc.east);
  \draw[sig] (proc) -- node[right,lbl]{any time} (met);
  \draw[sig] (setp) -- node[left,lbl]{rebuild} (setup);
\end{tikzpicture}
\caption{Block life cycle. \code{setup()} rebuilds derived state (filter taps, code matrices,
trellises) after construction and after every parameter change. \code{on_reset()} clears
accumulators before a run.}
\label{fig:lifecycle}
\end{figure}

\subsection{Reproducible randomness}
Every block owns a private NumPy generator. On reset it is re-seeded from the master seed of
the run \emph{and} the block's hierarchical path:
\codeexcerpt{block_reset}{Seeding the random generator of a block.}
Because the CRC-32 of the path enters the seed sequence, two \blk{awgn} blocks named
differently draw independent noise. The same graph with the same seed produces bit-identical
results, whatever the execution order of unrelated blocks and whether the run starts from the
GUI, the CLI or a script.

% ------------------------------------------------------------------
\section{The registry}
\code{@register_block("type_name", category=...)} adds a class to a global registry
(\file{core/registry.py}). Registered blocks can be created by name (\code{create_block},
\code{Graph.create}) and are described by \code{block_catalogue()}. The catalogue contains
the documentation, the ports and the JSON Schema of the parameters. The GUI builds its palette
and its parameter forms from this catalogue, and \cref{app:blocks} of this manual is generated
from it.

% ------------------------------------------------------------------
\section{Graphs}\label{sec:graphs}
A \code{Graph} holds named blocks and edges. Edges are stored as a map
$\text{destination input}\mapsto\text{source output}$, so every input has at most one driver,
while an output may fan out to any number of inputs. Ports are addressed as
\code{"block.port"}; omitting the port selects the block's first declared port, which makes
\code{g.chain("src", "map", "awgn")} a compact way to wire a pipeline. Parameters are
addressed as \code{"block.param"}, and \code{"/"} descends into subsystems
(\code{"rx/demap.method"}).

\subsection{Validation}\label{sec:validation}
\code{Graph.validate()} enforces three rules and returns the resolved kind of every output:
\begin{enumerate}
  \item every \emph{required} input is driven, either by an edge or, inside a subsystem, by an
  external input;
  \item the graph is acyclic;
  \item along every edge the produced kind intersects the accepted kinds.
\end{enumerate}
Kinds are resolved in topological order. An output declared with \code{same_as} takes the
kind actually delivered to the referenced input, intersected with the kinds that input accepts.
An \blk{interleaver} fed with LLRs therefore produces \kind{LLR}, and the validator can reject
an LLR output connected to a block that only accepts symbols, before anything is executed.

The topological order is computed with Kahn's algorithm:
\codeexcerpt{graph_toposort}{Topological sort of the block graph.}
Ties are broken by insertion order, so the execution order is deterministic. Feedback loops
(PLLs, adaptive filters, decision feedback) live \emph{inside} blocks, where they can run at the
sample rate. The graph itself carries one batch at a time from sources to sinks.

\begin{figure}[htbp]
\centering
\begin{tikzpicture}[node distance=6mm and 6mm]
  \node[blockbox] (src) {src\\\scriptsize\color{accent2}1};
  \node[blockbox,right=of src] (enc) {enc\\\scriptsize\color{accent2}2};
  \node[blockbox,right=of enc] (map) {map\\\scriptsize\color{accent2}3};
  \node[warmbox,right=of map] (ch) {ch\\\scriptsize\color{accent2}4};
  \node[blockbox,right=of ch] (eq) {eq\\\scriptsize\color{accent2}5};
  \node[blockbox,right=of eq] (dec) {dec\\\scriptsize\color{accent2}6};
  \node[greenbox,right=of dec] (ber) {ber\\\scriptsize\color{accent2}7};
  \draw[sig] (src)--(enc); \draw[sig] (enc)--(map); \draw[sig] (map)--(ch);
  \draw[sig] (ch)--node[above,lbl]{out}(eq); \draw[sig] (eq)--(dec); \draw[sig] (dec)--node[above,lbl]{rx}(ber);
  \draw[sig,kcsi] (ch.south) -- ++(0,-4mm) -| node[pos=0.25,below,lbl]{csi} (eq.south);
  \draw[sig,kbits] (src.north) -- ++(0,5mm) -| node[pos=0.25,above,lbl]{tx} (ber.north);
\end{tikzpicture}
\caption{A graph with fan-out (\code{src.out} drives \code{enc.in} and \code{ber.tx}) and a
side channel (\code{ch.csi}). Red numbers are the execution order.}
\end{figure}

\subsection{Execution}
\code{Graph.execute(ctx)} walks the blocks in order, gathers their inputs from the signals
already produced, checks the delivered kinds at run time and calls \code{process}:
\codeexcerpt{graph_execute}{The core of the executor.}
All produced signals are returned keyed by \code{"block.port"}; internal signals of
subsystems appear as \code{"sub/block.port"}. These keys are what the GUI probes and what
\code{Simulation.run(capture=...)} returns. Exceptions raised inside a block are re-raised as
\code{BlockError} with the block name, so the user sees which block failed and why.

% ------------------------------------------------------------------
\section{Subsystems}\label{sec:subsystems}
A \code{Subsystem} (\file{core/composite.py}) wraps a whole graph as a single block.
Its ports map to chosen inner ports, and its output kinds are those resolved by validating the
inner graph. A transmitter or a receiver can thus be built once and reused in many scenarios;
the same mechanism will let a future network layer treat a whole PHY link as one node.
Parameters and metrics of inner blocks are reached with the \code{"/"} separator:
\code{rx/dec.max_iter}, \code{rx/dec/avg_iterations}.

\begin{figure}[htbp]
\centering
\begin{tikzpicture}[node distance=5mm and 7mm]
  \node[blockbox] (src) {src};
  \node[draw=accent4,dashed,rounded corners,inner sep=7pt,right=12mm of src,
        label={[lbl,accent4]above:subsystem \code{tx} (file \file{coded_qam_tx.yaml})}] (tx) {
    \tikz[node distance=5mm]{\node[blockbox,minimum width=10mm] (e) {enc}; \node[blockbox,minimum width=10mm,right=of e] (m) {map}; \draw[sig] (e)--(m);}};
  \node[warmbox,right=12mm of tx] (ch) {awgn};
  \node[draw=accent4,dashed,rounded corners,inner sep=7pt,right=12mm of ch,
        label={[lbl,accent4]above:subsystem \code{rx}}] (rx) {
    \tikz[node distance=5mm]{\node[blockbox,minimum width=10mm] (d) {demap}; \node[blockbox,minimum width=10mm,right=of d] (c) {dec}; \draw[sig] (d)--(c);}};
  \draw[sig] (src) -- node[above,lbl]{bits} (tx);
  \draw[sig] (tx) -- node[above,lbl]{symbols} (ch);
  \draw[sig] (ch) -- (rx);
\end{tikzpicture}
\caption{Subsystems nest graphs; from outside they behave as ordinary blocks.}
\end{figure}

% ------------------------------------------------------------------
\section{Simulations and stop criteria}\label{sec:simulation}
A \code{Simulation} (\file{core/simulation.py}) runs a graph batch after batch. One batch
contains \code{batch_size} independent frames, which NumPy processes as a single vectorized
operation, so a larger batch amortizes the Python overhead. Meter blocks accumulate counts
across batches. The run stops at the first satisfied condition of a \code{StopCriteria}:
\codeexcerpt{stop_check}{Stop criteria are checked after every batch.}
The classic rule for error-rate estimation is \code{StopCriteria.errors("ber.bit_errors", 100)}:
stop once 100 errors have been counted. \Cref{sec:mc} shows why this gives a relative
accuracy of about $\pm20\,\%$ independent of the BER itself. \code{run} also accepts an
\code{on_batch} callback, which the web server uses to report progress and to cancel jobs.

% ------------------------------------------------------------------
\section{Sweeps}\label{sec:sweeps}
A \code{Sweep} (\file{core/sweep.py}) runs a fresh simulation at every point of a parameter
grid. In \code{grid} mode the points are the Cartesian product of the value lists; in
\code{zip} mode the lists are paired element-wise. It returns a table with one record per
point, which can be filtered, printed or written to CSV.
\codeexcerpt{sweep_point}{Each sweep point sets the parameters and runs a new simulation.}

\begin{keypoint}[Common random numbers]
By default every point reuses the same seed (\code{same_seed=True}). Two points then see the
same bits and the same noise realizations, scaled differently. If $X$ and $Y$ are the estimates
at two neighbouring points,
\[
  \Var(X-Y)=\Var X+\Var Y-2\operatorname{Cov}(X,Y),
\]
and the positive covariance created by shared randomness makes the \emph{differences} between
points much less noisy. The curves look smooth and comparisons between configurations are
fair. Use different seeds only when the points must be statistically independent.
\end{keypoint}

% ------------------------------------------------------------------
\section{Metadata propagation}\label{sec:metadata}
Physical quantities that a downstream block needs but cannot infer from the samples travel in
\code{Signal.meta}. Examples:
\begin{itemize}
  \item the constellation used by the mapper;
  \item the oversampling factor;
  \item the nominal signal power;
  \item the noise variance after equalization;
  \item the accumulated chromatic dispersion of a fiber link.
\end{itemize}
Each block copies the metadata of its input and adds or overrides what it changes.
\Cref{tab:meta-flow} follows the most important keys through a pulse-shaped link. The
mechanism is what lets the \blk{awgn} block set $E_b/N_0$ correctly whatever comes before it,
and lets a soft demapper compute exact LLRs without being told the SNR.
\Cref{app:metadata} lists all keys.

\begin{table}[htbp]
\centering\small
\caption{Metadata along a pulse-shaped QPSK link with $N_{\mathrm{sps}}$ samples per symbol
and $E_s/N_0=\gamma$. The waveform power is per sample; symbols have unit energy.}
\label{tab:meta-flow}
\begin{tabular}{@{}lcccc@{}}
\toprule
after block & \meta{power} & \meta{sps} & \meta{noise_var} & also sets\\
\midrule
\blk{mapper}         & 1 & 1 & -- & \meta{constellation}, \meta{bits_per_symbol}\\
\blk{pulse_shaper}   & $1/N_{\mathrm{sps}}$ & $N_{\mathrm{sps}}$ & -- & \meta{pulse_taps}, \meta{delay}\\
\blk{awgn} (\prm{esn0}) & $1/N_{\mathrm{sps}}$ & $N_{\mathrm{sps}}$ & $N_0=1/\gamma$ & \\
\blk{matched_filter} & 1 & 1 & $1/\gamma$ & \meta{n_symbols}\\
\blk{demapper}       & \multicolumn{4}{l}{uses \meta{constellation} and \meta{noise_var} for exact LLRs}\\
\bottomrule
\end{tabular}
\end{table}

\begin{pitfall}
Metadata describes the \emph{nominal} signal. A block that changes the signal in a way it
cannot describe (a CTLE, an unknown channel) sets the affected keys to \code{None}, and
downstream blocks fall back to measuring. When you write a block, propagate or invalidate
every key your processing changes.
\end{pitfall}

% ------------------------------------------------------------------
\section{Scenarios}\label{sec:scenarios}
A scenario (\file{core/scenario.py}) is a graph plus its run configuration, stored as YAML:
\begin{lstlisting}[style=yaml]
name: QPSK over AWGN
category: Tutorials            # used by the GUI scenario browser
description: ...
seed: 1
batch_size: 64
blocks:
  src:   {type: random_bits, params: {n_bits: 2000}}
  map:   {type: mapper, params: {scheme: psk, order: 4}}
  awgn:  {type: awgn, params: {snr_db: 6}}
  demap: {type: demapper, params: {method: hard}}
  ber:   {type: ber_meter}
  rx:    {type: subsystem, file: subsystems/coded_qam_rx.yaml}   # optional
chain: [src, map, awgn, demap, ber.rx]      # default-port connections
edges:
  - src.out -> ber.tx                        # explicit connections
simulation: {stop: {metric: ber.bit_errors, threshold: 200, max_batches: 500}}
sweep: {params: {awgn.snr_db: [0, 2, 4, 6, 8]}}   # optional: run a sweep instead
probes: [{key: awgn.out, view: constellation}]   # GUI scopes
layout: {src: [0, 0], map: [230, 0]}             # GUI positions
\end{lstlisting}
Files saved from the GUI contain only non-default parameters, so they stay short and readable
and remain usable from the CLI. \Cref{app:scenarios} lists the scenarios shipped with LinkStudio.

\chapter{Graphical editor, API and command line}\label{ch:interfaces}

\section{The graphical editor}
\code{python -m linkstudio serve} starts the web server and opens the editor in the browser.
\Cref{fig:gui-layout} shows its layout.

\begin{figure}[htbp]
\centering
\begin{tikzpicture}[x=1cm,y=1cm,font=\sffamily\scriptsize]
  \draw[black!40,fill=black!3,rounded corners=2pt] (0,0) rectangle (14,8.2);
  % top bar
  \draw[accent,fill=soft] (0.1,7.5) rectangle (13.9,8.1);
  \node[anchor=west] at (0.2,7.8) {\textbf{scenario name} \quad New \quad Save \quad Save as \quad Undo/Redo \quad Layout \quad ? \quad Export .py \quad seed};
  \node[anchor=east,accent3] at (13.8,7.8) {\textbf{valid}};
  % palette
  \draw[accent4,fill=accent4!6] (0.1,2.9) rectangle (2.6,7.4);
  \node[align=center] at (1.35,6.95) {\textbf{Blocks}\\search};
  \foreach \y/\t in {6.2/sources,5.75/encoders,5.3/mappers,4.85/modulators,4.4/media,3.95/decoders,3.5/measurements}
    \node[draw=black!25,fill=white,minimum width=2.1cm,minimum height=3.5mm] at (1.35,\y) {\t};
  % canvas
  \draw[black!30,fill=white] (2.7,2.9) rectangle (10.6,7.4);
  \node at (6.65,7.15) {\textbf{Canvas}};
  \foreach \x/\t/\c in {3.4/src/kbits,4.9/map/ksym,6.4/awgn/ksym,7.9/demap/kllr,9.6/ber/gray}
    \node[draw=\c,fill=\c!10,rounded corners=1pt,minimum width=1.1cm,minimum height=6mm] (n\t) at (\x,5.2) {\t};
  \draw[->,kbits,thick] (nsrc)--(nmap); \draw[->,ksym,thick] (nmap)--(nawgn);
  \draw[->,ksym,thick] (nawgn)--(ndemap); \draw[->,kllr,thick] (ndemap)--(nber);
  \draw[->,kbits,thick] (nsrc.south) -- ++(0,-0.6) -| (nber.south);
  \draw[black!30,fill=black!5] (9.2,3.0) rectangle (10.5,3.8); \node at (9.85,3.4) {minimap};
  % inspector
  \draw[accent2,fill=accent2!5] (10.7,2.9) rectangle (13.9,7.4);
  \node[align=left,anchor=north west] at (10.8,7.3) {\textbf{Inspector}\\[2pt] name\\ parameters (forms,\\ sliders, \texttt{auto})\\ frequency response\\[3pt] Flip \;\; Duplicate\\ Delete};
  % dock
  \draw[accent3,fill=accent3!6] (0.1,0.1) rectangle (13.9,2.8);
  \node[anchor=north west] at (0.2,2.75) {\textbf{Live} \quad Simulate \quad Sweep};
  \draw[black!30,fill=white] (0.3,0.3) rectangle (3.0,2.2); \node at (1.65,1.25) {metrics + CI};
  \foreach \x/\t in {3.2/constellation,5.9/spectrum,8.6/eye diagram,11.3/BER curve}
    { \draw[black!30,fill=white] (\x,0.3) rectangle ++(2.5,1.9); \node at (\x+1.25,1.25) {\t}; }
\end{tikzpicture}
\caption{Layout of the editor: palette, canvas, inspector and the run dock with its three
modes.}
\label{fig:gui-layout}
\end{figure}

\paragraph{Palette.}
The \numblocks{} blocks are grouped by functional class in signal-chain order: sources, encoders, bit
processing, mappers, multiplexers, framing, modulators, filters and resamplers, amplifiers and
attenuators, transmission media, impairments, then the receiver classes (demodulators,
synchronization, equalizers, detectors and estimators, demultiplexers, demappers, decoders) and
the measurements. Complementary classes (encoders and decoders, mappers and demappers, \dots)
share a colour. Only the large classes are split into subclasses (code families, media); the
others are flat lists. Every block draws its function on the canvas: codewords with parity, the
shift register and trellis of convolutional codes, the Tanner graph, the interleaver permutation
drawn from its seed (or the block matrix numbered in write order), the scrambler LFSR with its
real taps and seed, pulse shapes, constellations, and the physical blocks (adders, mixers,
antennas, fibers, lasers) as drawings without a frame.

\paragraph{Editing.}
\begin{itemize}
  \item \textbf{Adding blocks:} drag a block from the palette or double-click it.
  \item \textbf{Connecting:} drag from an output handle to an input handle. Handle colours are
  the signal kinds of \cref{tab:kinds}; incompatible connections are refused, and dashed
  handles are optional inputs.
  \item \textbf{Selection:} click a block; Shift or Ctrl + click adds blocks, Shift + drag on
  the empty canvas selects the blocks in a box, Ctrl+A selects all and Esc clears the
  selection. With several blocks selected a bar over the canvas offers \emph{Align row} (the
  links between them become straight), \emph{Align column}, \emph{Distribute} (equal gaps),
  \emph{Flip}, \emph{Duplicate} and \emph{Delete}.
  \item \textbf{Clipboard:} Ctrl+C, Ctrl+X and Ctrl+V copy, cut and paste blocks with the
  connections between them; the copy is placed at the mouse and pasting again steps the
  copies. Ctrl+D duplicates, Delete deletes, the arrow keys move the selection (Shift + arrow
  by a larger step).
  \item \textbf{Direction:} a block faces right (inputs on the left, outputs on the right) or
  left. F, or \emph{Flip} in the inspector, turns the selected blocks around. A block dragged
  to the other side of the blocks it is linked to turns around by itself, so a chain can run
  right to left or turn back on itself without moving any port. The block drawings are
  mirrored with the block; their text stays readable and plots against time or frequency keep
  their axes.
  \item \textbf{Moving ports:} Ctrl + drag a handle to move it anywhere on the block border,
  e.g.\ for feedback paths. Port places and directions are saved in the scenario
  (\code{port_layout}).
  \item \textbf{Alignment:} while a block is dragged it snaps so that its connected ports line
  up with the ports at the other end (straight links), or else so that its edges or centre line
  up with a neighbouring block; a guide line shows the alignment.
  \item \textbf{Automatic layout} (Ctrl+L) arranges the blocks along the signal flow
  (\cref{fig:auto-layout}). Each block goes in the column given by its longest distance from a
  source, and sources that feed late blocks are pulled forward. The longest path (the main
  chain) stays on one straight line. Branches such as training sequences, channel estimates
  and meters go in lanes under it, ordered by barycentre sweeps to reduce crossings. A long
  chain wraps into rows sized from the shape of the canvas and run alternately left to right
  and right to left (a horseshoe); the blocks of the reversed rows face left. Scopes and meters
  that tap a link from their bottom side sit above the gap after the block they tap.
  Unconnected blocks go in a last row.
  \item \textbf{Links} are drawn with orthogonal segments routed around the blocks (A$^\ast$
  search on a sparse grid of block edges and channels, with a penalty per bend), so they never
  pass behind a block. All links are routed together. Where links of different nets would run
  on top of each other along a channel, their shared segments are moved apart so that each
  link stays visible.
  \item \textbf{Inspector:} parameter forms are generated from each block's schema. Bounded and
  dB-valued parameters get sliders, and \code{auto} leaves a value to upstream metadata. For
  linear blocks the inspector also plots the \emph{frequency response} (magnitude, phase or group
  delay): the model transfer function of cables, CTLE, lasers and VCSELs, photodiodes, single-
  and multimode fibers and pulse shapers, and an estimate $\hat H=S_{xy}/S_{xx}$ of one
  realization for random multipath channels.
  \item \textbf{Undo:} Ctrl+Z and Ctrl+Y undo and redo any change: blocks, connections,
  parameters, positions, port places and directions, scopes.
  \item \textbf{Shortcuts:} ? (or the ? button) lists them all. Besides those above: Ctrl+S
  saves, Ctrl+Shift+S saves as, Ctrl+O opens the scenario browser, Ctrl+B and Ctrl+I show or hide
  the block palette and the inspector.
  Saving keeps the run settings (\code{simulation}, \code{sweep}) and the ports of a subsystem
  file, so a scenario saved from the GUI runs the same from the command line.
\end{itemize}

\begin{figure}[htbp]
\centering
\begin{tikzpicture}[x=1cm,y=1cm,font=\sffamily\scriptsize,
    b/.style={draw=accent,fill=soft,rounded corners=1pt,minimum width=1.25cm,minimum height=6mm},
    l/.style={->,thick,black!60}]
  \foreach \i/\t in {0/src,1/enc,2/map,3/mod,4/tx}
    \node[b] (r\i) at (1.9*\i,2) {\t};
  \foreach \i/\t in {4/fiber,3/rx,2/eq,1/demap,0/dec}
    \node[b] (q\i) at (1.9*\i,0) {\t};
  \node[b,draw=gray,fill=gray!8] (ber) at (-1.9,0) {ber};
  \node[b,draw=gray,fill=gray!8,minimum width=1cm] (sc) at (2.85,2.9) {scope};
  \foreach \i/\j in {0/1,1/2,2/3,3/4} \draw[l] (r\i) -- (r\j);
  \draw[l] (r4.east) -- ++(0.35,0) |- (q4.east);
  \foreach \i/\j in {4/3,3/2,2/1,1/0} \draw[l] (q\i) -- (q\j);
  \draw[l] (q0) -- (ber);
  \draw[l] (r0.south) -- ++(0,-0.4) -| (ber.north);
  \fill[black!60] (2.85,2) circle (1.3pt);
  \draw[thick,black!60] (2.85,2) -- (sc.south);
  \node[anchor=west,commentcolor] at (8.9,2) {row 1: faces right};
  \node[anchor=west,commentcolor] at (8.9,0) {row 2: faces left};
\end{tikzpicture}
\caption{Automatic layout of a long chain: it wraps into a second row that runs right to left
with its blocks facing left. A scope tapping a link from its bottom side sits above it.}
\label{fig:auto-layout}
\end{figure}

The run buttons are disabled while the graph is invalid, and the validator's message is shown
next to them.

\paragraph{Modes.}
\begin{description}[font=\sffamily\bfseries]
  \item[Live] simulates one frame per step (batch size 1), 1--30 steps per second, and streams
  every probe; the metrics accumulate across steps. Parameter changes are applied between steps
  and the meters restart, so the effect of a slider is visible immediately. Structural edits
  restart the session. The rate only sets how often the scopes refresh, not the results.
  \item[Simulate] runs a Monte-Carlo job in the background, either for a fixed number of
  batches or until a metric reaches a threshold. \emph{frames/batch} sets the batch size (saved in
  the scenario) and \emph{cores} the number of worker processes (\cref{sec:parallel}). It shows the
  final metrics (BER with its 95\,\% confidence interval) and the probes of the last batch.
  \item[Sweep] runs one or two nested parameter sweeps and plots any metric, with an optional
  theoretical curve; with several cores the points run in parallel. Curves can be kept and
  overlaid, and the results exported as CSV or PNG.
\end{description}

\paragraph{The Scope block.}
Signals are displayed by the \blk{scope} block (class \emph{measurements}), an instrument with a
screen. Its inputs enter from the bottom side: drop the scope from the palette onto a link, or
drag one of its free inputs onto a link, to tap the output that feeds it. The tap is drawn as a
junction dot just after that output. A scope has up to four inputs; the selector in its corner
chooses the displayed one, and the bottom-right corner resizes it. The plot selector under its
name (parameter \prm{view}) lists the views that suit the displayed signal:
\begin{center}\small
\begin{tabular}{@{}ll@{}}
\toprule
Kind & Views (the first is the \code{auto} choice)\\
\midrule
\kind{SAMPLES} & spectrum (Welch), constellation, eye diagram, time\\
\kind{SYMBOLS} & constellation, time, spectrum\\
\kind{GRID} & constellation, resource-grid magnitude\\
\kind{CSI} & impulse / frequency response\\
\kind{LLR} & histogram, time\\
\kind{BITS} & time\\
\bottomrule
\end{tabular}
\end{center}
A view that does not suit the signal falls back to the first one of its kind. After a
\emph{Simulate} run every view is already computed, so switching needs no new run; in
\emph{Live} the next frame shows the new view, and the measurements are not restarted. Older
scenarios with one scope type per view (\code{scope_constellation}, \dots) are converted to
\blk{scope} with the matching view when they are loaded.

Axes carry physical units. The server sends values in base SI units together with their unit,
taken from the \meta{unit} metadata: $\sqrt{\mathrm{W}}$ for optical fields, A
for photocurrents, none for normalized baseband signals. The editor picks the SI prefix
(ps, GHz, mA, \dots) so that tick labels stay short. Spectra are absolute power spectral
densities (dBm/Hz for optical fields, one-sided for real signals). The eye diagram shows time
in seconds over two unit intervals and its vertical axis is centred on the mean of the traces
and spans all of them.

\section{The HTTP and WebSocket API}
The editor talks to the FastAPI application in \file{linkstudio/api/app.py}. Everything it does
can also be scripted.

\begin{center}\small
\begin{tabularx}{\textwidth}{@{}llX@{}}
\toprule
Method & Path & Purpose\\
\midrule
GET  & \code{/api/health} & server version and number of CPU cores\\
GET  & \code{/api/blocks} & block catalogue: documentation, ports, parameter JSON Schema\\
GET  & \code{/api/scenarios} & scenario list with name, category, description, block count\\
GET/PUT & \code{/api/scenarios/{path}} & load / save a scenario in editor format\\
POST & \code{/api/validate} & validate an editor graph; returns errors and port kinds\\
POST & \code{/api/layout} & server-side layout (the editor lays out locally)\\
POST & \code{/api/export/python} & standalone Python script equivalent to the graph\\
POST & \code{/api/response} & frequency response of one block (magnitude, phase, group delay)\\
POST & \code{/api/jobs} & start a simulate or sweep job\\
GET  & \code{/api/jobs/{id}} & job status, progress, partial and final results\\
POST & \code{/api/jobs/{id}/cancel} & cancel a job\\
WS   & \code{/api/live} & interactive live mode\\
\bottomrule
\end{tabularx}
\end{center}

Jobs run in a thread pool (\file{api/jobs.py}); the \code{workers} entry of the job
configuration sends the batches or sweep points to that many worker processes. Their results
contain JSON-safe metrics (NaN
becomes \code{null}), the probe data of the last batch and, for sweeps of a single parameter
whose value matches a known modulation, theoretical BER curves. A failed job reports the
message \code{failed} with the details in a separate \code{detail} field.

\begin{figure}[htbp]
\centering
\begin{tikzpicture}[x=1cm,y=-0.62cm,font=\sffamily\scriptsize]
  \node[blockbox] (c) at (0,0) {browser};
  \node[blockbox] (s) at (7,0) {server (\code{/api/live})};
  \draw[black!40] (0,0.6) -- (0,10.5); \draw[black!40] (7,0.6) -- (7,10.5);
  \draw[->] (0,1.5) -- node[above]{\code{start} \{graph, rate\}} (7,1.5);
  \draw[<-] (0,2.5) -- node[above]{\code{started}} (7,2.5);
  \draw[<-,accent3] (0,3.6) -- node[above]{\code{frame} \{batch, metrics, probes, step\_ms\}} (7,3.6);
  \draw[<-,accent3] (0,4.6) -- node[above]{\code{frame}} (7,4.6);
  \draw[->,accent2] (0,5.6) -- node[above]{\code{set_param} \{path, value\}} (7,5.6);
  \node[anchor=west,note] at (7.1,6.3) {applied before the next batch;\\meters restart};
  \draw[<-,accent3] (0,7.0) -- node[above]{\code{frame}} (7,7.0);
  \draw[->] (0,8.0) -- node[above]{\code{probes} / \code{rate} / \code{reset}} (7,8.0);
  \draw[->] (0,9.2) -- node[above]{\code{stop}} (7,9.2);
  \draw[<-] (0,10.2) -- node[above]{\code{stopped}} (7,10.2);
\end{tikzpicture}
\caption{Message flow of the live mode. The server steps the simulation in a worker thread and
paces frames to the requested rate.}
\end{figure}

\section{Command line}
\begin{center}\small
\begin{tabularx}{\textwidth}{@{}lX@{}}
\toprule
\code{python -m linkstudio blocks [--category C]} & list block types\\
\code{python -m linkstudio describe TYPE} & ports, parameters and documentation of a block\\
\code{python -m linkstudio run FILE.yaml [--csv OUT]} & run a scenario (simulation or sweep)\\
\code{python -m linkstudio serve [--host H] [--port P] [--no-browser]} & start the editor\\
\code{python scripts/main.py [--quick]} & run all demos with theory checks\\
\bottomrule
\end{tabularx}
\end{center}
The demo scripts in \file{scripts/} (one per subsystem) are short, readable examples of the
Python API. Each prints simulated values next to theory and returns a pass/fail status.

\chapter{Developer guide: writing blocks}\label{ch:extending}

Adding a model to LinkStudio means writing one block class. The engine, the CLI, the scenario
loader and the GUI pick it up automatically once it is registered. This chapter builds a
complete example and lists the conventions a block must respect to behave like the built-in
ones.

\section{A complete example: free-space path loss}
Suppose we want a block that attenuates a radio signal by the free-space loss between two
isotropic antennas at distance $d$ and carrier frequency $f_c$ (Friis):
\begin{equation}
  L_{\mathrm{fs}} = \left(\frac{4\pi d f_c}{c}\right)^{2},\qquad
  L_{\mathrm{fs}}\big|_{\mathrm{dB}} = 20\log_{10} d + 20\log_{10} f_c - 147.55 .
\end{equation}
The amplitude is scaled by $L_{\mathrm{fs}}^{-1/2}$. The nominal power metadata and any noise
already present must be scaled by $L_{\mathrm{fs}}^{-1}$, so that a downstream \blk{awgn} block
in \prm{nominal} mode still computes the SNR from the right reference.

\begin{lstlisting}[style=py,caption={A new block, e.g. in \file{linkstudio/phy/channel/pathloss.py}.},label=lst:example-block]
import numpy as np
from pydantic import Field

from linkstudio.core import Block, BlockParams, Kind, Port, RunContext, Signal, register_block

C_LIGHT = 299_792_458.0


@register_block("free_space_loss", category="medium", subcategory="Wireless", label="Free-Space Loss")
class FreeSpaceLoss(Block):
    """Friis free-space path loss between isotropic antennas."""

    class Params(BlockParams):                       # validated, serializable, shown in the GUI
        distance_m: float = Field(100.0, gt=0, description="Link distance.")
        carrier_hz: float = Field(2.4e9, gt=0, description="Carrier frequency.")

    inputs = {"in": Port(Kind.SYMBOLS | Kind.SAMPLES)}
    outputs = {"out": Port.like("in")}               # same kind as the input

    def setup(self) -> None:                         # derived state, rebuilt on parameter change
        p = self.params
        self.loss = (4 * np.pi * p.distance_m * p.carrier_hz / C_LIGHT) ** 2

    def on_reset(self) -> None:                      # runtime state, cleared before a run
        self._frames = 0

    def process(self, inputs: dict[str, Signal], ctx: RunContext) -> dict[str, Signal]:
        x = inputs["in"]
        meta = {k: x.meta[k] / self.loss for k in ("power", "noise_var") if x.meta.get(k) is not None}
        self._frames += x.batch_size
        return {"out": x.evolve(x.data / np.sqrt(self.loss), **meta)}

    def metrics(self) -> dict[str, float]:
        return {"loss_db": 10 * np.log10(self.loss), "frames": self._frames}
\end{lstlisting}

To register it, import the module in the package \code{__init__.py}
(\file{linkstudio/phy/channel/__init__.py}). From then on:
\begin{itemize}
  \item \code{g.create("pl", "free_space_loss", distance_m=500)} works from Python;
  \item \code{{type: free_space_loss, params: {distance_m: 500}}} works in YAML;
  \item the block appears in the GUI palette under \emph{Transmission media} $\to$ \emph{Wireless} as \emph{Free-Space Loss}, with a form for its two
  parameters;
  \item \code{pl.distance_m} can be swept;
  \item \code{pl.loss_db} is reported among the metrics.
\end{itemize}

\section{Conventions}
\begin{enumerate}[label=\textbf{C\arabic*.}]
  \item \textbf{Shapes.} Accept any leading batch shape. Operate on the last axis and use
  \code{...} indexing, so the block works for \code{(B, n)} as well as \code{(B, P, n)} arrays.
  \item \textbf{Vectorize over the batch.} Loops over samples are acceptable for adaptive
  algorithms, but each iteration should process all frames at once (see the LMS loop in
  \cref{lst:lms_loop}).
  \item \textbf{Kinds.} Declare the narrowest accepted kinds; use \code{Port.like} when the
  output kind follows the input.
  \item \textbf{Metadata.} Propagate with \code{evolve}. Update the keys your processing
  changes (\meta{power}, \meta{noise_var}, \meta{sps}, \meta{delay}, \dots) and set them to
  \code{None} when you cannot describe the result.
  \item \textbf{Randomness.} Draw only from \code{self.rng}; never use NumPy's global
  generator. This keeps runs reproducible (\cref{sec:blocks}).
  \item \textbf{Real signals stay real.} Wireline and IM/DD chains are real-valued; do not
  return complex arrays for real inputs.
  \item \textbf{Errors.} Raise \code{BlockError} with a message telling the user how to fix the
  problem (\emph{``Timing recovery needs an oversampled input (sps >= 2)''}).
  \item \textbf{Metrics accumulate.} Sum counts in \code{process}, compute ratios in
  \code{metrics}, clear them in \code{on_reset}.
\end{enumerate}

\section{Testing a block against theory}
The test-suite follows one rule: a physical model is validated against a closed-form result,
not against itself. \file{tests/helpers.py} provides \code{uncoded_link} and \code{ber_at}. A
typical test reads:
\begin{lstlisting}[style=py,numbers=none]
def test_awgn_ber_matches_theory():
    g = uncoded_link("qam", 16, snr_db=10.0, n_bits=4 * 1200)
    assert np.isclose(ber_at(g, errors=300), theory.ber_qam(16, 10.0), rtol=0.2)
\end{lstlisting}
With 300 errors the relative standard deviation of the estimate is about $1/\sqrt{300}\approx
6\,\%$, so a 20\,\% tolerance makes false alarms negligible but still detects a 1\,dB error. For
our example the natural check is the metric itself: at \SI{100}{m} and \SI{2.4}{GHz},
$L_{\mathrm{fs}}=80.05$\,dB. \Cref{ch:validation} lists the physics checks of the whole library.

\section{Adding a scenario}
Save a graph from the GUI, or write a YAML file under \file{backend/scenarios/<category>/}
with \code{name}, \code{category} and \code{description}. The test
\file{tests/test_scenarios.py} loads and runs every scenario file, so a broken scenario is
caught automatically.


\part{Theory and implementation}
\chapter[Baseband signal model, noise and Monte-Carlo estimation]{Baseband signal model, noise and\\ Monte-Carlo estimation}\label{ch:baseband}

Every chapter of Part II relies on three conventions:
\begin{itemize}
  \item a signal is represented by its sampled complex envelope;
  \item noise is circular complex Gaussian with a variance tied to $N_0$;
  \item performance is estimated by counting errors.
\end{itemize}
This chapter makes the three precise.

% ------------------------------------------------------------------
\section{The complex envelope}
A real band-pass signal centred on the carrier $f_c$ can always be written as
\begin{equation}
  s(t)=\Re\!\left\{x(t)\,e^{j2\pi f_c t}\right\}=x_I(t)\cos 2\pi f_c t-x_Q(t)\sin 2\pi f_c t,
  \qquad x(t)=x_I(t)+j\,x_Q(t),
  \label{eq:envelope}
\end{equation}
where $x(t)$ is the \emph{complex envelope}. If $s$ occupies $|f\pm f_c|<B/2$, then $x$ is a
low-pass signal of bandwidth $B/2$. It can be sampled at $f_s\ge B$, whatever the carrier
frequency. Linear time-invariant band-pass systems act on $x$ through their low-pass
equivalent impulse response. A carrier frequency error $\Delta f$ becomes a rotation
$x(t)e^{j2\pi\Delta f t}$, and an oscillator phase error $\phi(t)$ becomes $x(t)e^{j\phi(t)}$.

\begin{incode}
All \kind{SYMBOLS} and \kind{SAMPLES} signals are sampled complex envelopes
$x[n]=x(n/f_s)$ with \code{Signal.fs}$=f_s$. Two exceptions are worth noting:
\begin{itemize}
  \item wireline and IM/DD electrical signals are \emph{real} baseband signals ($x_Q\equiv0$),
  and the blocks keep their arrays real;
  \item optical fields are envelopes around the optical carrier, scaled in $\sqrt{\mathrm{W}}$
  (\cref{ch:optical}).
\end{itemize}
Powers are \emph{envelope} powers, $P=\E|x[n]|^2$. The factor $\tfrac12$ relating envelope and
band-pass powers cancels in every SNR, so it is never applied.
\end{incode}

% ------------------------------------------------------------------
\section{Additive white Gaussian noise}\label{sec:awgn}
\subsection{From $N_0$ to a sample variance}
Thermal noise at the receiver input is modelled as white with two-sided band-pass power
spectral density $N_0/2$. Its complex envelope $w(t)$ is circular complex Gaussian with
two-sided PSD $N_0$ (in envelope units). Observed through any filter of noise bandwidth $f_s$
and sampled at $f_s$, it gives independent samples
\begin{equation}
  w[n]\sim\CN(0,\sigma^2),\qquad \sigma^2=N_0 f_s,
  \qquad \E\big[(\Re w)^2\big]=\E\big[(\Im w)^2\big]=\sigma^2/2 .
  \label{eq:sample-variance}
\end{equation}
The code works in units where the sample period is the time unit, so $N_0$ and $\sigma^2$
coincide numerically. The metadata key \meta{noise_var} is $\sigma^2$, the noise variance per
complex sample (or per symbol after matched filtering).

\subsection{SNR definitions}
Let $P$ be the signal power per sample and $N_{\mathrm{sps}}$ the number of samples per symbol.
The energy per symbol is $E_s=P\,N_{\mathrm{sps}}T_s$ and the noise PSD is
$N_0=\sigma^2 T_s$, so
\begin{equation}
  \frac{E_s}{N_0}=\frac{P\,N_{\mathrm{sps}}}{\sigma^2}.
  \label{eq:esn0}
\end{equation}
A symbol carries $k=\log_2 M$ coded bits. A code of rate $R$ and $N_s$ spatially multiplexed
streams turn these into $k R N_s$ information bits per channel use, so
\begin{equation}
  \frac{E_b}{N_0}=\frac{E_s}{k R N_s N_0}
  \quad\Longrightarrow\quad
  \sigma^2=\frac{P\,N_{\mathrm{sps}}}{(E_b/N_0)\,k R N_s}.
  \label{eq:ebn0}
\end{equation}
The \blk{awgn} block implements the three definitions of \cref{tab:snr-modes}.
$N_{\mathrm{sps}}$, $k$, $R$ and $N_s$ are read from the metadata set by the pulse shaper (or
OFDM modulator), the mapper, the encoders and the stream demultiplexer. \Cref{eq:ebn0} is
therefore correct for any chain, which is why the coded and MIMO curves of later chapters can
be compared with uncoded theory on the same $E_b/N_0$ axis.

\begin{table}[htbp]
\centering\small
\caption{SNR modes of the \blk{awgn} block.}
\label{tab:snr-modes}
\begin{tabular}{@{}lll@{}}
\toprule
\prm{mode} & meaning of \prm{snr_db} & noise variance $\sigma^2$\\
\midrule
\code{snr}  & signal-to-noise ratio per sample & $P/\mathrm{SNR}$\\
\code{esn0} & $E_s/N_0$ & $P N_{\mathrm{sps}}/(E_s/N_0)$\\
\code{ebn0} & $E_b/N_0$ (default) & $P N_{\mathrm{sps}}/\big((E_b/N_0)\,k R N_s\big)$\\
\bottomrule
\end{tabular}
\end{table}

The reference power $P$ is, in order of priority:
\begin{enumerate}
  \item the parameter \prm{signal_power};
  \item the nominal power \meta{power} carried by the metadata (\prm{power_reference}\code{=nominal},
  the textbook definition of \emph{average} SNR);
  \item the power measured over the batch.
\end{enumerate}
The nominal reference matters for fading: the average SNR is defined by the transmitted power,
not by the instantaneous power after a deep fade.

\codeexcerpt{awgn_noise_var}{Noise variance from the SNR definition (\file{awgn.py}).}
\codeexcerpt{awgn_process}{Adding the noise. Real inputs get $\sigma^2/2$ per real dimension.}

\begin{remark}
For a real signal only the in-phase noise component matters, with variance $\sigma^2/2=N_0/2$.
A real antipodal (NRZ, 2-PAM) signal then has the same error probability as BPSK, as it should.
This is tested in \file{tests/test_wireline_optical.py}.
\end{remark}

% ------------------------------------------------------------------
\section{Monte-Carlo estimation of error rates}\label{sec:mc}
\subsection{The estimator and its accuracy}
Let $N$ bits be observed and $K$ of them be in error. If errors are independent with
probability $p$, then $K\sim\mathrm{Binomial}(N,p)$ and the natural estimator
$\hat p=K/N$ is unbiased with
\begin{equation}
  \Var\hat p=\frac{p(1-p)}{N},\qquad
  \frac{\sqrt{\Var\hat p}}{p}=\sqrt{\frac{1-p}{Np}}\approx\frac{1}{\sqrt{\E K}} .
\end{equation}
The relative accuracy depends only on the number of \emph{errors}, not on the number of bits.
Stopping a simulation after a fixed number of errors (\cref{sec:simulation}) therefore gives the
same relative accuracy at every SNR: 100 errors give about $\pm10\,\%$ standard deviation, and
$\pm20\,\%$ at 95\,\% confidence.

\subsection{The Wilson score interval}
The normal approximation $\hat p\pm z\sqrt{\hat p(1-\hat p)/N}$ fails when $K$ is small: with
$K=0$ it gives a zero-width interval. The Wilson interval~\cite{wilson1927} instead inverts the test
$|\hat p-p|\le z\sqrt{p(1-p)/N}$ for the unknown $p$. Squaring gives the quadratic
$(1+z^2/N)p^2-(2\hat p+z^2/N)p+\hat p^2\le0$, whose roots are
\begin{equation}
  p_{\pm}=\frac{\hat p+\dfrac{z^2}{2N}\pm z\sqrt{\dfrac{\hat p(1-\hat p)}{N}+\dfrac{z^2}{4N^2}}}{1+\dfrac{z^2}{N}},
  \qquad z=1.96\ \text{for 95\,\%}.
  \label{eq:wilson}
\end{equation}
The interval is never empty, stays inside $[0,1]$ and is accurate even for a handful of errors.
\blk{ber_meter} reports it as \code{ber_ci_low} and \code{ber_ci_high}, and the GUI shows it next
to the BER.
\codeexcerpt{wilson}{The Wilson interval in \file{metrics.py}.}

\begin{figure}[htbp]
\centering
\begin{tikzpicture}
\begin{semilogxaxis}[width=10cm,height=5.2cm,xlabel={number of observed errors $K$},
  ylabel={$p_\pm/\hat p$},ymin=0,ymax=3,xmin=1,xmax=2000,legend pos=north east]
  \addplot[name path=hi,accent,thick] table[x=errors,y=hi]{\datafile{wilson}};
  \addplot[name path=lo,accent,thick] table[x=errors,y=lo]{\datafile{wilson}};
  \addplot[accent!15] fill between[of=hi and lo];
  \addplot[black!50,dashed,domain=1:2000,samples=2] {1};
  \legend{95\,\% Wilson bounds}
\end{semilogxaxis}
\end{tikzpicture}
\caption{Width of the 95\,\% confidence interval relative to the estimate, as a function of the
number of observed errors (computed with \code{wilson_interval} at $p=10^{-4}$). Ten errors leave
a factor of about two of uncertainty; one hundred errors bring it to $\pm20\,\%$.}
\label{fig:wilson}
\end{figure}

\begin{pitfall}[Correlated errors]
The binomial model assumes independent errors. After a decoder, or in block fading, errors come
in bursts within a frame. The interval is then optimistic, because the effective number of
independent observations is closer to the number of erroneous \emph{frames}. For coded systems,
judge the accuracy by \code{frame_errors} as well.
\end{pitfall}

% ------------------------------------------------------------------
\section{Closed-form references}
\file{linkstudio/phy/theory.py} collects the closed-form results used in tests, demos, the GUI
and this manual. All functions accept arrays and SNRs in dB.
\begin{center}\small
\begin{tabularx}{\textwidth}{@{}lX@{}}
\toprule
\code{ber_bpsk}, \code{ber_pam}, \code{ber_qam} & exact Gray-coded BER over AWGN (\cref{sec:ber-awgn})\\
\code{ser_psk}, \code{ber_psk} & Craig's integral for M-PSK\\
\code{ber_bpsk_rayleigh}, \code{ber_bpsk_mrc}, \code{ber_rayleigh} & fading and diversity (\cref{sec:rayleigh-ber})\\
\code{ber_hamming_hard} & post-decoding BER of Hamming codes\\
\code{ook_q_factor}, \code{q_to_ber}, \code{osnr_to_esn0_db} & optical receivers (\cref{ch:optical})\\
\code{gaussian_pulse_broadening} & dispersion of a Gaussian pulse\\
\code{shannon_capacity}, \code{ebn0_min_db} & information-theoretic limits\\
\bottomrule
\end{tabularx}
\end{center}

\chapter{Digital modulation and demodulation}\label{ch:modulation}

% ------------------------------------------------------------------
\section{Constellations}
A memoryless modulator maps each group of $k$ bits $\vect b=(b_0,\dots,b_{k-1})$ to one of
$M=2^k$ complex points, $s=\mathcal{S}(\vect b)$. The label $\ell=\sum_i b_i2^{k-1-i}$ is
read MSB first and indexes the table of points. LinkStudio builds four families, all normalized to
unit average energy $\E|s|^2=1$:
\begin{align}
  \text{M-PAM:}\quad & s_\ell = \frac{2i-(M-1)}{\sqrt{(M^2-1)/3}},\quad \ell=g(i),\\
  \text{M-QAM:}\quad & s_\ell = \frac{a_{\ell_I}+j\,a_{\ell_Q}}{\sqrt{2(M-1)/3}},\quad
     \ell=(\ell_I,\ell_Q),\ a_\cdot\ \text{Gray }\sqrt M\text{-PAM levels},\\
  \text{M-PSK:}\quad & s_\ell = e^{\,j(2\pi i/M+\phi_0)},\quad \ell=g(i),\quad
     \phi_0=\pi/4\ \text{for QPSK, else }0,
\end{align}
where $g(i)=i\oplus\lfloor i/2\rfloor$ is the binary-reflected Gray code. Adjacent points
therefore differ in exactly one bit. The denominators are the average energies of the
unnormalized PAM and QAM grids, $\frac1M\sum_i(2i-M+1)^2=(M^2-1)/3$ per dimension. Custom
constellations are normalized the same way.

\begin{figure}[htbp]
\centering
\begin{tikzpicture}[scale=0.85]
  \draw[->,black!50] (-4.4,0)--(4.4,0) node[right,lbl]{I};
  \draw[->,black!50] (0,-4.4)--(0,4.4) node[above,lbl]{Q};
  \foreach \x/\bx in {-3/00,-1/01,1/11,3/10}
    \foreach \y/\by in {-3/00,-1/01,1/11,3/10} {
      \fill[accent] (\x,\y) circle (2.5pt);
      \node[font=\scriptsize\ttfamily,above=1pt] at (\x,\y) {\textcolor{accent2}{\bx}\textcolor{accent3}{\by}};
    }
  \foreach \x/\bx in {-3/00,-1/01,1/11,3/10}
    \node[font=\scriptsize\ttfamily,accent2,below] at (\x,-4.0) {\bx};
  \node[lbl,accent2,below] at (0,-4.45) {I bits $(b_0b_1)$};
  \node[lbl,accent3,rotate=90] at (-4.7,0) {Q bits $(b_2b_3)$};
  \draw[dashed,black!25] (-2,-4) -- (-2,4) (0,-4)--(0,4) (2,-4)--(2,4) (-4,-2)--(4,-2) (-4,2)--(4,2);
\end{tikzpicture}
\caption{16-QAM with the labels produced by \code{qam(16)}. The first half of the label selects
the in-phase level and the second half the quadrature level, each with a Gray 4-PAM labelling.
Dashed lines are the decision boundaries.}
\label{fig:qam16}
\end{figure}

\codeexcerpt{qam_build}{Square QAM as the product of two Gray PAMs (\file{dsp/constellations.py}).}

\begin{incode}
\blk{mapper} (\file{phy/modulation.py}) builds the constellation in \code{setup()} and maps
frames of $n k$ bits to $n$ symbols. It sets \meta{constellation}, \meta{bits_per_symbol},
\meta{power}$=1$, \meta{sps}$=1$ and the symbol rate \prm{symbol_rate} as \code{fs}.
\blk{demapper} takes the constellation from the metadata unless one is given explicitly.
\end{incode}

% ------------------------------------------------------------------
\section{Optimal detection}
For the observation $y=s+w$ with $w\sim\CN(0,\sigma^2)$ and equiprobable symbols, the
maximum-likelihood decision is the nearest point:
\begin{equation}
  \hat s=\argmax_{s\in\mathcal{S}} p(y\,|\,s)=\argmax_{s}\frac{1}{\pi\sigma^2}e^{-|y-s|^2/\sigma^2}
  =\argmin_{s\in\mathcal{S}}|y-s|^2 .
\end{equation}
\code{Constellation.nearest} evaluates all $M$ distances in chunks of $2^{16}$ symbols, which
bounds memory. \prm{method}\code{=hard} returns the label bits of the nearest point.

\subsection{Bit log-likelihood ratios}\label{sec:llr}
Soft-input decoders need, for every bit, the a-posteriori log-likelihood ratio. With
$\mathcal{S}_i^{(b)}$ the subset of points whose $i$-th label bit equals $b$, and equiprobable
bits,
\begin{equation}
  L(b_i)=\ln\frac{P(b_i=0\,|\,y)}{P(b_i=1\,|\,y)}
  =\ln\!\sum_{s\in\mathcal{S}_i^{(0)}}\!e^{-|y-s|^2/\sigma^2}
   -\ln\!\sum_{s\in\mathcal{S}_i^{(1)}}\!e^{-|y-s|^2/\sigma^2}.
  \label{eq:llr-exact}
\end{equation}
Keeping only the dominant term of each sum gives the max-log approximation
\begin{equation}
  L(b_i)\approx\frac{1}{\sigma^2}\Big(\min_{s\in\mathcal{S}_i^{(1)}}|y-s|^2-\min_{s\in\mathcal{S}_i^{(0)}}|y-s|^2\Big),
  \label{eq:llr-maxlog}
\end{equation}
which is exact for BPSK and loses only a few tenths of a dB with good codes. A positive LLR
means bit $0$ throughout LinkStudio. The variance $\sigma^2$ may differ from symbol to symbol, for
example after zero-forcing equalization of a fading channel; \meta{noise_var} may therefore be
an array broadcastable to the symbols.
\codeexcerpt{demap_soft}{Exact (log-sum-exp) and max-log LLRs.}

\begin{keypoint}
The LLR scale matters. Equalizers update \meta{noise_var} with their noise enhancement
($\sigma^2/|h|^2$ for a one-tap ZF equalizer, $\sigma^2\|\vect w\|^2$ for an FIR equalizer,
the post-detection variance for MIMO). The decoder downstream then receives LLRs with the
correct reliability without any manual tuning.
\end{keypoint}

% ------------------------------------------------------------------
\section{Error probability over AWGN}\label{sec:ber-awgn}
With $\Q(x)=\frac12\erfc(x/\sqrt2)$ the Gaussian tail, the classic results implemented in
\file{theory.py} are:
\begin{itemize}
  \item \textbf{BPSK and Gray QPSK}: $P_b=\Q\big(\sqrt{2E_b/N_0}\big)$. QPSK is two orthogonal BPSK
  channels sharing the symbol energy.
  \item \textbf{Gray M-PAM}, exact (Cho and Yoon~\cite{cho2002}): with $L=M$, $m=\log_2L$ and
  $\gamma_b=E_b/N_0$,
  \begin{equation}
    P_b=\frac1m\sum_{q=1}^{m}\frac1L\!\!\sum_{i=0}^{(1-2^{-q})L-1}\!\!
    (-1)^{\lfloor i2^{q-1}/L\rfloor}\Big(2^{q-1}-\Big\lfloor\frac{i2^{q-1}}{L}+\frac12\Big\rfloor\Big)
    \erfc\!\Big((2i+1)\sqrt{\tfrac{3m\gamma_b}{L^2-1}}\Big).
  \end{equation}
  \item \textbf{Square Gray M-QAM}: the in-phase and quadrature components are independent
  $\sqrt M$-PAM signals carrying half the bits each, so $P_b^{\text{QAM}}(M)=P_b^{\text{PAM}}(\sqrt M)$
  at the same $E_b/N_0$.
  \item \textbf{M-PSK}: the symbol error rate is Craig's integral~\cite{craig1991}
  \begin{equation}
    P_s=\frac1\pi\int_0^{(M-1)\pi/M}\exp\!\Big(-\frac{E_s}{N_0}\frac{\sin^2(\pi/M)}{\sin^2\theta}\Big)d\theta,
    \qquad P_b\approx P_s/k ,
  \end{equation}
  where the Gray approximation $P_b\approx P_s/k$ is tight at moderate SNR.
\end{itemize}

\begin{figure}[htbp]
\centering
\begin{tikzpicture}
\begin{semilogyaxis}[berplot,width=12cm,height=8cm,xmin=-2,xmax=16,ymin=1e-6,ymax=0.5,
  legend columns=2,legend pos=south west,legend style={font=\scriptsize}]
  \addplot[th=accent] table[x=ebn0,y=bpsk]{\datafile{ber_awgn_theory}};
  \addplot[sim=accent] table[x=ebn0,y=bpsk]{\datafile{ber_awgn_sim}};
  \addplot[th=accent3] table[x=ebn0,y=psk8]{\datafile{ber_awgn_theory}};
  \addplot[sim={accent3,mark=square*}] table[x=ebn0,y=psk8]{\datafile{ber_awgn_sim}};
  \addplot[th=accent2] table[x=ebn0,y=qam16]{\datafile{ber_awgn_theory}};
  \addplot[sim={accent2,mark=triangle*}] table[x=ebn0,y=qam16]{\datafile{ber_awgn_sim}};
  \addplot[th=accent4] table[x=ebn0,y=qam64]{\datafile{ber_awgn_theory}};
  \addplot[sim={accent4,mark=diamond*}] table[x=ebn0,y=qam64]{\datafile{ber_awgn_sim}};
  \draw[black!60,dashed] (axis cs:-1.59,1e-6) -- (axis cs:-1.59,0.5)
     node[pos=0.08,right,font=\scriptsize,black!70]{Shannon limit $-1.59$\,dB};
  \legend{BPSK/QPSK theory,simulated,8-PSK theory,simulated,16-QAM theory,simulated,64-QAM theory,simulated}
\end{semilogyaxis}
\end{tikzpicture}
\caption{Bit error rate over AWGN: simulated (markers, chain
\blk{random_bits}$\to$\blk{mapper}$\to$\blk{awgn}$\to$\blk{demapper}$\to$\blk{ber_meter}, up to
200 errors per point) and exact theory (lines). Points with fewer than 10 errors are omitted.}
\label{fig:ber-awgn}
\end{figure}

\paragraph{The ultimate limit.}
For a spectral efficiency $\eta$ bit/s/Hz, the AWGN capacity $C=\log_2(1+\mathrm{SNR})$ with
$\mathrm{SNR}=\eta E_b/N_0$ requires $E_b/N_0\ge(2^\eta-1)/\eta$, which tends to
$\ln2=-1.59$\,dB as $\eta\to0$ (\code{ebn0_min_db}). The gap between this limit and the uncoded
curves of \cref{fig:ber-awgn} is what channel coding recovers (\cref{ch:coding}).

\chapter{Pulse shaping and matched filtering}\label{ch:pulse}

% ------------------------------------------------------------------
\section{Linear modulation}
A linearly modulated waveform is a train of pulses weighted by the symbols,
\begin{equation}
  x(t)=\sum_m a_m\,g(t-mT).
\end{equation}
For i.i.d. zero-mean symbols with $\E|a_m|^2=E_a$, the power spectral density is
\begin{equation}
  S_x(f)=\frac{E_a}{T}\,|G(f)|^2 ,
  \label{eq:linear-psd}
\end{equation}
so the spectrum is shaped entirely by the pulse. The receiver filters with $h(t)$ and samples
the output at $t=mT+\tau_0$. With $p=g\star h$ the end-to-end pulse,
\begin{equation}
  y(mT+\tau_0)=a_m\,p(\tau_0)+\underbrace{\sum_{l\neq m}a_l\,p\big((m-l)T+\tau_0\big)}_{\text{intersymbol interference}}+\,(w\star h)(mT+\tau_0).
\end{equation}

\section{The Nyquist criterion}
The ISI vanishes for every symbol sequence if and only if $p(nT)=\delta[n]$ (taking
$\tau_0=0$). By the Poisson summation formula this is equivalent to the frequency-domain
condition
\begin{equation}
  \sum_{k=-\infty}^{\infty}P\!\left(f-\frac kT\right)=T\qquad\forall f .
  \label{eq:nyquist}
\end{equation}
The narrowest solution is the ideal low-pass of bandwidth $1/(2T)$, the sinc pulse, which is
not realizable and decays only as $1/t$. The \emph{raised-cosine} (RC) family trades bandwidth
for decay. Its spectrum has a cosine roll-off of relative width $\beta$:
\begin{equation}
  P_{\mathrm{RC}}(f)=
  \begin{cases}
    T, & |f|\le\frac{1-\beta}{2T},\\[2pt]
    \frac T2\Big[1+\cos\Big(\frac{\pi T}{\beta}\big(|f|-\frac{1-\beta}{2T}\big)\Big)\Big], & \frac{1-\beta}{2T}<|f|\le\frac{1+\beta}{2T},\\[2pt]
    0, & \text{otherwise},
  \end{cases}
\end{equation}
\begin{equation}
  p_{\mathrm{RC}}(t)=\sinc\frac tT\,\frac{\cos(\pi\beta t/T)}{1-(2\beta t/T)^2}.
\end{equation}
The occupied bandwidth is $(1+\beta)R_s$ and the pulse decays as $1/t^3$.

\section{The matched filter}
Among all receive filters, the one maximizing the SNR at the sampling instant follows from the
Cauchy--Schwarz inequality. For white noise of PSD $N_0$,
\begin{equation}
  \mathrm{SNR}=\frac{\big|\int G(f)H(f)\,df\big|^2}{N_0\int|H(f)|^2df}
  \le\frac{\int|G(f)|^2df}{N_0}=\frac{E_g}{N_0},
\end{equation}
with equality for $H(f)=G\conj(f)$, that is $h(t)=g\conj(-t)$. To have both a matched filter
and zero ISI, the RC response is split equally between transmitter and receiver. Each side
uses the \emph{root-raised-cosine} (RRC) pulse, $|G_{\mathrm{RRC}}(f)|^2=P_{\mathrm{RC}}(f)$,
whose impulse response is
\begin{equation}
  g_{\mathrm{RRC}}(t)=\frac{\sin\big(\pi\frac tT(1-\beta)\big)+4\beta\frac tT\cos\big(\pi\frac tT(1+\beta)\big)}
  {\pi\frac tT\big(1-(4\beta\frac tT)^2\big)},
\end{equation}
with the finite limits $1-\beta+4\beta/\pi$ at $t=0$ and
$\frac{\beta}{\sqrt2}\big[(1+\frac2\pi)\sin\frac{\pi}{4\beta}+(1-\frac2\pi)\cos\frac{\pi}{4\beta}\big]$ at
$t=\pm T/(4\beta)$. These special points are handled explicitly in the code:
\codeexcerpt{rrc_taps}{RRC taps with their singular points (\file{dsp/filters.py}).}

\begin{figure}[htbp]
\centering
\begin{tikzpicture}
\begin{groupplot}[group style={group size=2 by 1,horizontal sep=16mm},width=6.9cm,height=5cm]
\nextgroupplot[xlabel={$t/T$},ylabel={normalized amplitude},xmin=-4,xmax=4,ymin=-0.3,ymax=1.1,
  xtick={-4,...,4},legend pos=north east,title={pulses, $\beta=0.25$}]
  \addplot[th=accent] table[x=t,y=rc]{\datafile{pulses_time}};
  \addplot[th=accent2,dashed] table[x=t,y=rrc]{\datafile{pulses_time}};
  \legend{RC (Nyquist),RRC}
\nextgroupplot[xlabel={$fT$},ylabel={$|P_{\mathrm{RC}}(f)|/P(0)$},xmin=0,xmax=1.1,ymin=0,ymax=1.1,
  legend pos=north east,title={RC spectrum}]
  \addplot[th=black!60] table[x=f,y=b0]{\datafile{pulses_freq}};
  \addplot[th=accent] table[x=f,y=b25]{\datafile{pulses_freq}};
  \addplot[th=accent3] table[x=f,y=b50]{\datafile{pulses_freq}};
  \addplot[th=accent2] table[x=f,y=b100]{\datafile{pulses_freq}};
  \legend{$\beta=0$,$\beta=0.25$,$\beta=0.5$,$\beta=1$}
\end{groupplot}
\end{tikzpicture}
\caption{Left: the RC pulse crosses zero at every non-zero multiple of $T$, while the RRC pulse
alone does not. Right: spectra of the truncated RC taps produced by \code{rc_taps} (span
32 symbols). All curves pass through $0.5$ at $fT=0.5$, as required by \cref{eq:nyquist}.}
\label{fig:pulses}
\end{figure}

% ------------------------------------------------------------------
\section{Discrete-time implementation}
\begin{incode}
\blk{pulse_shaper} upsamples by $N_{\mathrm{sps}}$ (zero stuffing) and convolves with
unit-energy taps spanning \prm{span} symbols. It records the taps in \meta{pulse_taps}, the
group delay $(L-1)/2$ in \meta{delay}, and the power per sample $1/N_{\mathrm{sps}}$ (unit
symbol energy spread over $N_{\mathrm{sps}}$ samples).

\blk{matched_filter} convolves with $g\conj[-n]$ (or another pulse), samples at the known
delay and rescales.
\end{incode}

\codeexcerpt{pulse_shaper}{The pulse shaper and its metadata.}

With unit-energy taps the peak of the end-to-end pulse is
$p[0]=\sum_n|g[n]|^2=1$, and the noise gain of the matched filter is $\sum_n|h[n]|^2=1$. The
block nevertheless normalizes by the actual peak and scales \meta{noise_var} by
$\sum|h|^2/|p[0]|^2$, so that mismatched receive filters (an RC receive filter, a rectangular
integrate-and-dump) also yield symbols at nominal amplitude with the correct noise variance:
\codeexcerpt{mf_normalize}{Normalization of the matched filter output.}

Following \cref{eq:esn0}, the AWGN block sets $\sigma^2=P N_{\mathrm{sps}}/(E_s/N_0)=1/(E_s/N_0)$
on the waveform. After the matched filter the symbols have unit energy and noise variance
$\sigma^2$, so the symbol SNR equals $E_s/N_0$, the matched-filter bound. This is verified in
\file{tests/test_physics.py::test_pulse_shaped_link_matches_theory}: the BER of an
RRC-shaped link equals the symbol-rate theory.

\begin{pitfall}[Truncation]
Taps are truncated to \prm{span} symbols. Short spans raise the spectral side lobes and leave
a small residual ISI that grows as $\beta\to0$. For dense WDM or very small roll-offs, use
spans of 32 symbols or more (the coherent optical scenarios use 32).
\end{pitfall}

\section{Eye diagrams}
Overlaying traces of two symbol periods shows the ISI at a glance (\cref{fig:eye}). After the
RRC transmit filter alone the eye is partially closed, because an RRC pulse is not a Nyquist
pulse. After the matched filter all traces pass through $\pm1$ at the sampling instant.

\begin{figure}[htbp]
\centering
\begin{tikzpicture}
\begin{groupplot}[group style={group size=2 by 1,horizontal sep=14mm},width=6.9cm,height=5cm,
  xmin=-1,xmax=1,ymin=-1.5,ymax=1.5,xlabel={$t/T$},xtick={-1,-0.5,0,0.5,1}]
\nextgroupplot[title={after the RRC transmit filter},ylabel={amplitude}]
  \addplot[accent,very thin,opacity=0.5] table[x=x,y=y]{\datafile{eye_tx}};
\nextgroupplot[title={after the matched filter}]
  \addplot[accent3,very thin,opacity=0.5] table[x=x,y=y]{\datafile{eye_mf}};
  \draw[accent2,dashed] (axis cs:0,-1.5) -- (axis cs:0,1.5);
\end{groupplot}
\end{tikzpicture}
\caption{Eye diagrams of binary PAM with RRC pulses ($\beta=0.35$, 16 samples/symbol, 30\,dB
SNR), cut from the outputs of \blk{pulse_shaper} and \blk{matched_filter} (\prm{decimate}\code{=false}).
The dashed line is the sampling instant used by the matched filter.}
\label{fig:eye}
\end{figure}

\chapter{Radio front-end impairments}\label{ch:impairments}

Real transmitters and receivers distort the signal before any propagation effect. This chapter
models the impairments implemented in \file{phy/channel/impairments.py}. \Cref{fig:impair-const}
shows their signatures on a 16-QAM constellation.

\begin{figure}[htbp]
\centering
\begin{tikzpicture}
\begin{groupplot}[group style={group size=4 by 1,horizontal sep=4mm,y descriptions at=edge left},
  constaxis,width=4.1cm,height=4.1cm]
\nextgroupplot[title={AWGN only}]
  \addplot[scatterpts=accent] table[x=i,y=q]{\datafile{const_awgn}};
\nextgroupplot[title={phase noise}]
  \addplot[scatterpts=accent3] table[x=i,y=q]{\datafile{const_pn}};
\nextgroupplot[title={IQ imbalance}]
  \addplot[scatterpts=accent4] table[x=i,y=q]{\datafile{const_iq}};
\nextgroupplot[title={Saleh PA, 4\,dB IBO}]
  \addplot[scatterpts=accent2] table[x=i,y=q]{\datafile{const_pa}};
\end{groupplot}
\end{tikzpicture}
\caption{16-QAM at $E_s/N_0=26$\,dB through \blk{phase_noise} (150\,Hz linewidth at
1\,MBd), \blk{iq_imbalance} (1.5\,dB, $8^\circ$) and \blk{power_amplifier} (Saleh). Phase
noise smears points along circles; IQ imbalance turns the square grid into a parallelogram;
the PA compresses and rotates the outer points.}
\label{fig:impair-const}
\end{figure}

% ------------------------------------------------------------------
\section{Carrier frequency and phase offsets}
A local oscillator at $f_c+\Delta f$ with phase $\phi_0$ turns the received envelope into
\begin{equation}
  y[n]=x[n]\,e^{\,j(2\pi\Delta f\,n/f_s+\phi_0)} .
\end{equation}
\blk{carrier_offset} applies this rotation (\prm{cfo_hz}, \prm{phase_deg}, or a uniformly random
phase per frame). Synchronization (\cref{ch:sync}) estimates and removes it.

% ------------------------------------------------------------------
\section{Oscillator phase noise}\label{sec:phase-noise}
A free-running oscillator accumulates phase errors. The standard model, and the one used for
lasers in \cref{ch:optical}, is a Wiener process: the instantaneous frequency deviation is
white, so the phase performs a random walk. In discrete time,
\begin{equation}
  \phi[n]=\phi[n-1]+\varepsilon[n],\qquad \varepsilon[n]\sim\N\!\Big(0,\ \frac{2\pi\Delta\nu}{f_s}\Big).
  \label{eq:wiener}
\end{equation}
\begin{proposition}
The field $e^{j\phi(t)}$ of a Wiener phase with increment variance $2\pi\Delta\nu|\tau|$ has a
Lorentzian spectrum of full width at half maximum $\Delta\nu$.
\end{proposition}
\begin{proof}
$\phi(t+\tau)-\phi(t)\sim\N(0,2\pi\Delta\nu|\tau|)$, and $\E e^{jX}=e^{-\Var X/2}$ for a Gaussian
$X$, so the autocorrelation is $R(\tau)=e^{-\pi\Delta\nu|\tau|}$. Its Fourier transform is
\[
  S(f)=\frac{\Delta\nu/(2\pi)}{f^2+(\Delta\nu/2)^2},
\]
which falls to half its peak at $f=\pm\Delta\nu/2$.
\end{proof}
\codeexcerpt{phase_noise}{Wiener phase noise (\blk{phase_noise}).}
The phase variance grows linearly with time. What matters for a receiver is therefore the
product $\Delta\nu\,T$ of linewidth and symbol duration: carrier recovery must track the phase
faster than it wanders.

% ------------------------------------------------------------------
\section{IQ imbalance}
In a quadrature down-converter the Q branch may have a gain $g\ne1$ and a phase error $\varphi$
relative to the I branch. Writing the two branch outputs and recombining them gives the
widely-linear model
\begin{equation}
  y=K_1x+K_2x\conj+d,\qquad K_1=\frac{1+g\,e^{-j\varphi}}{2},\qquad K_2=\frac{1-g\,e^{j\varphi}}{2},
  \label{eq:iq}
\end{equation}
where $d$ is a DC offset. The conjugate term creates a mirror image of the signal. Its
strength is measured by the image rejection ratio $\mathrm{IRR}=|K_1|^2/|K_2|^2$; for $g=1.5$\,dB
and $\varphi=8^\circ$ it is about $19$\,dB. On a constellation it shears the square grid
(\cref{fig:impair-const}).
\codeexcerpt{iq_imbalance}{IQ imbalance and DC offset (\blk{iq_imbalance}).}

% ------------------------------------------------------------------
\section{Power amplifier nonlinearity}
A memoryless band-pass nonlinearity acts on the envelope $x=re^{j\theta}$ as
\begin{equation}
  y=A(r)\,e^{\,j(\theta+\Phi(r))},
\end{equation}
where $A(r)$ is the AM/AM curve and $\Phi(r)$ the AM/PM curve. The input is first scaled so that
its mean power lies \prm{input_backoff_db} below the saturation level $A_{\mathrm{sat}}=1$
(input back-off, IBO). Three models are implemented:
\begin{align}
  \text{Rapp (solid state):}\quad & A(r)=\frac{r}{\big(1+r^{2p}\big)^{1/(2p)}},\qquad\Phi=0,\\
  \text{Saleh (TWTA):}\quad & A(r)=\frac{r}{1+\beta_a r^2},\qquad \Phi(r)=\frac{\alpha_\phi r^2}{1+\beta_\phi r^2},\\
  \text{ideal clipper:}\quad & A(r)=\min(r,1),\qquad\Phi=0.
\end{align}
The Saleh AM/AM curve is divided by $\alpha_a$ so that all models have unit small-signal gain;
the default constants are those of Saleh's original TWT fit~\cite{saleh1981}. The smoothness
$p$ of the Rapp model interpolates between a soft compression ($p=1$) and the ideal clipper
($p\to\infty$).

\begin{figure}[htbp]
\centering
\begin{tikzpicture}
\begin{groupplot}[group style={group size=2 by 1,horizontal sep=18mm},width=6.8cm,height=5.2cm,
  xlabel={input amplitude $r$},xmin=0,xmax=2]
\nextgroupplot[ylabel={output amplitude $A(r)$},ymin=0,ymax=1.2,legend pos=south east,title={AM/AM}]
  \addplot[th=black!40,dashed,domain=0:1.2] {x};
  \addplot[th=accent] table[x=r,y=rapp1]{\datafile{pa_curves}};
  \addplot[th=accent3] table[x=r,y=rapp2]{\datafile{pa_curves}};
  \addplot[th=accent4] table[x=r,y=rapp5]{\datafile{pa_curves}};
  \addplot[th=black!70] table[x=r,y=clip]{\datafile{pa_curves}};
  \addplot[th=accent2] table[x=r,y=saleh]{\datafile{pa_curves}};
  \legend{linear,Rapp $p{=}1$,Rapp $p{=}2$,Rapp $p{=}5$,clipper,Saleh}
\nextgroupplot[ylabel={phase shift $\Phi(r)$ [deg]},ymin=0,ymax=25,title={AM/PM (Saleh)}]
  \addplot[th=accent2] table[x=r,y=saleh_pm_deg]{\datafile{pa_curves}};
\end{groupplot}
\end{tikzpicture}
\caption{PA characteristics measured by driving the \blk{power_amplifier} block with an amplitude
ramp (back-off chosen so that the block's input scaling is unity).}
\label{fig:pa}
\end{figure}
\codeexcerpt{power_amplifier}{Back-off scaling and the AM/AM, AM/PM models.}

\begin{pitfall}
The nonlinearity generates spectral regrowth (out-of-band emission) and an in-band error that
is \emph{not} Gaussian. For OFDM, whose envelope is nearly Gaussian (\cref{sec:papr}), an IBO
of several dB is needed. Compare the \blk{power_meter} PAPR with the back-off.
\end{pitfall}

% ------------------------------------------------------------------
\section{Analog-to-digital conversion}
\blk{quantizer} models a uniform mid-rise ADC with $b$ bits on each of I and Q. The full scale
$A$ is either given or set to \prm{loading_factor}$\times$ the RMS value $\sigma_x$ of each
component. The step is $\Delta=2A/2^b$. For a signal that does not clip, the quantization error is
approximately uniform with variance $\Delta^2/12$, giving
\begin{equation}
  \mathrm{SQNR}=\frac{\sigma_x^2}{\Delta^2/12}=\frac{3\cdot2^{2b}}{(A/\sigma_x)^2}
  \ \Rightarrow\ \mathrm{SQNR}_{\mathrm{dB}}\approx6.02\,b+4.77-20\log_{10}\frac{A}{\sigma_x}.
\end{equation}
With the default loading factor of 4 (clipping probability about $6\times10^{-5}$ for a Gaussian
component), an 8-bit ADC gives about 41\,dB.
\codeexcerpt{quantizer}{Uniform mid-rise quantization with clipping.}

% ------------------------------------------------------------------
\section{Timing offset}
\blk{timing_offset} delays the signal by $d=\lfloor d\rfloor+\mu$ samples. The integer part is a
shift. The fractional part uses a Blackman-windowed sinc interpolator of 33 taps,
$h[n]=\sinc(n-16-\mu)\,w[n]/\sum_n\sinc(n-16-\mu)\,w[n]$, which is accurate as long as the signal
is oversampled. The frame length is preserved, so a delay pushes the tail out of the frame; the
\blk{preamble} block appends guard symbols for this reason.

\chapter{Fading channels and equalization}\label{ch:fading}

% ------------------------------------------------------------------
\section{The multipath channel}
In a radio channel the transmitted signal reaches the receiver along many paths with
different delays $\tau_l$, complex gains and Doppler shifts. The low-pass equivalent channel is
a linear time-varying filter,
\begin{equation}
  y(t)=\sum_l a_l(t)\,x(t-\tau_l)+w(t),\qquad h(\tau;t)=\sum_l a_l(t)\,\delta(\tau-\tau_l).
\end{equation}
Two scales characterize it:
\begin{itemize}
  \item the \emph{delay spread} $\sigma_\tau$, whose inverse is the coherence bandwidth
  $B_c\approx1/(5\sigma_\tau)$;
  \item the maximum Doppler shift $f_D=v f_c/c$, whose inverse sets the coherence time
  $T_c\approx0.423/f_D$.
\end{itemize}
A signal of bandwidth $B\ll B_c$ sees \emph{flat} fading (a single complex gain). If $B\gtrsim
B_c$ the fading is frequency selective and causes ISI.

% ------------------------------------------------------------------
\section{Rayleigh and Rician fading}
When a path is the superposition of many unresolved scatterers with random phases, the
central limit theorem makes its gain circular complex Gaussian, $h\sim\CN(0,1)$. The envelope
$|h|$ is then Rayleigh distributed and the power $|h|^2$ exponential. A line-of-sight
component of relative power $K$ (the Rician factor) adds a deterministic term:
\begin{equation}
  h=\sqrt{\frac{K}{K+1}}\,e^{j\theta_0}+\sqrt{\frac{1}{K+1}}\,h_{\mathrm{Ray}},\qquad \E|h|^2=1 .
\end{equation}

\subsection{Time correlation: the Clarke--Jakes model}
For isotropic scattering in the horizontal plane, a path arriving at angle $\alpha$ is shifted
by $f_D\cos\alpha$. Averaging over a uniform $\alpha$ gives~\cite{clarke1968} the autocorrelation and the U-shaped
Doppler spectrum
\begin{equation}
  R_h(\tau)=\E\big[h(t+\tau)h\conj(t)\big]=J_0(2\pi f_D\tau),\qquad
  S_h(f)=\frac{1}{\pi f_D\sqrt{1-(f/f_D)^2}},\ |f|<f_D .
  \label{eq:jakes}
\end{equation}
LinkStudio generates it with a random-phase sum of sinusoids:
\begin{equation}
  h[n]=\frac{1}{\sqrt{N}}\sum_{i=1}^{N}e^{\,j(2\pi f_D\cos\alpha_i\,n/f_s+\phi_i)},\qquad
  \alpha_i,\phi_i\sim\mathcal U[0,2\pi).
\end{equation}
Each term has autocorrelation $\E_\alpha e^{j2\pi f_D\cos\alpha\,\tau}=J_0(2\pi f_D\tau)$, so
the sum reproduces \cref{eq:jakes} exactly in expectation. Its first-order statistics approach
Gaussian as $N$ grows (16 sinusoids by default).
\codeexcerpt{sos}{Sum-of-sinusoids Rayleigh process (\file{channel/fading.py}).}

\begin{figure}[htbp]
\centering
\begin{tikzpicture}
\begin{groupplot}[group style={group size=2 by 1,horizontal sep=17mm},width=6.9cm,height=5cm]
\nextgroupplot[xlabel={$f_D\tau$},ylabel={$R_h(\tau)$},xmin=0,xmax=3,ymin=-0.5,ymax=1.05,
  legend pos=north east,title={autocorrelation}]
  \addplot[th=accent] table[x=lag,y=theory]{\datafile{jakes_acf}};
  \addplot[sim=accent2,mark size=1.1pt] table[x=lag,y=sim]{\datafile{jakes_acf}};
  \legend{$J_0(2\pi f_D\tau)$,simulated}
\nextgroupplot[xlabel={$f_D t$},ylabel={$|h|^2$ [dB]},xmin=0,xmax=4,ymin=-30,ymax=10,
  title={fading envelope}]
  \addplot[accent,thin] table[x=t,y=env_db]{\datafile{jakes_envelope}};
\end{groupplot}
\end{tikzpicture}
\caption{Left: autocorrelation of \code{sum_of_sinusoids} (400 realizations, $f_D/f_s=0.01$)
against \cref{eq:jakes}. Right: a realization; fades deeper than 20\,dB occur about once per
wavelength travelled.}
\label{fig:jakes}
\end{figure}

\begin{incode}
\blk{flat_fading} multiplies the input by $h$ and publishes the true gains on its \code{csi}
output (\meta{csi_type}\code{="gain"}). Three time-variation modes are available:
\begin{itemize}
  \item \prm{block}: one coefficient per frame (quasi-static);
  \item \prm{fast}: independent per sample (ideal interleaving);
  \item \prm{jakes}: correlated with Doppler \prm{doppler_hz}; needs \code{fs}.
\end{itemize}
\end{incode}
\codeexcerpt{flat_fading}{Fading generation and the Rician line-of-sight term.}

% ------------------------------------------------------------------
\section{Tapped delay lines and 3GPP profiles}
\blk{multipath} implements the discrete model
$y[n]=\sum_{l}h_l[n]\,x[n-d_l]$ with integer sample delays $d_l=\mathrm{round}(\tau_lf_s)$.
Taps falling on the same sample have their powers added, and the profile is normalized to unit
total power. Taps are Rayleigh (or static with \prm{fading}\code{=false}), constant per frame or
each with an independent Jakes process. The 3GPP extended models of TS~36.104~\cite{3gpp36104} are built in:

\begin{center}\small
\begin{tabular}{@{}lll@{}}
\toprule
Model & delays [ns] & relative powers [dB]\\
\midrule
EPA & 0, 30, 70, 90, 110, 190, 410 & 0, $-1$, $-2$, $-3$, $-8$, $-17.2$, $-20.8$\\
EVA & 0, 30, 150, 310, 370, 710, 1090, 1730, 2510 & 0, $-1.5$, $-1.4$, $-3.6$, $-0.6$, $-9.1$, $-7$, $-12$, $-16.9$\\
ETU & 0, 50, 120, 200, 230, 500, 1600, 2300, 5000 & $-1$, $-1$, $-1$, 0, 0, 0, $-3$, $-5$, $-7$\\
\bottomrule
\end{tabular}
\end{center}
Their RMS delay spreads are 43, 357 and 991\,ns. The channel impulse response is published
on \code{csi}: \meta{csi_type}\code{="cir"} with shape $(B,L)$ for block fading, or
\code{"cir_tv"} with shape $(B,L,n)$ for time-varying taps.

% ------------------------------------------------------------------
\section{Error probability with fading}\label{sec:rayleigh-ber}
With coherent detection and perfect channel knowledge, the instantaneous SNR of BPSK in flat
fading is $\gamma=|h|^2E_b/N_0$, exponentially distributed with mean $\bar\gamma$. Averaging
the AWGN error probability over it gives
\begin{equation}
  \bar P_b=\int_0^\infty\Q\big(\sqrt{2\gamma}\big)\,\frac{e^{-\gamma/\bar\gamma}}{\bar\gamma}\,d\gamma
  =\frac12\left(1-\sqrt{\frac{\bar\gamma}{1+\bar\gamma}}\right)\ \xrightarrow{\ \bar\gamma\to\infty\ }\ \frac{1}{4\bar\gamma}.
  \label{eq:rayleigh-ber}
\end{equation}
The error rate now falls only as $1/\mathrm{SNR}$ instead of exponentially: 10\,dB more SNR
buys one decade of BER. Diversity, the subject of \cref{sec:diversity}, restores a steeper
slope. \code{ber_rayleigh} performs the same average numerically for any AWGN curve.
\Cref{fig:diversity} compares \cref{eq:rayleigh-ber} with the chain
\blk{flat_fading}$\to$\blk{awgn}$\to$\blk{one_tap_equalizer}.

% ------------------------------------------------------------------
\section{Equalization}\label{sec:equalization}
\subsection{One-tap equalization}
For flat fading, $y=hs+w$, and dividing by the known gain gives $\hat s=s+w/h$. The noise
variance becomes $\sigma^2/|h|^2$ symbol by symbol, and \blk{one_tap_equalizer} writes exactly
this array to \meta{noise_var}, so the soft demapper weights faded symbols correctly. The
\prm{phase} variant only derotates, which is sufficient for PSK.

\subsection{Linear FIR equalizers designed from the channel}
For a frequency-selective channel with impulse response $\vect h=(h_0,\dots,h_{L-1})$, a
linear equalizer $\vect w$ of $N$ taps produces the combined response $\mat H\vect w$, where
$\mat H$ is the $(N+L-1)\times N$ convolution matrix of $\vect h$. Zero forcing asks the combined
response to approximate a unit impulse $\vect e_d$ at decision delay $d$. The MMSE design adds
the noise power $\sigma^2$ as a regularizer:
\begin{equation}
  \vect w_d=\argmin_{\vect w}\ \|\mat H\vect w-\vect e_d\|^2+\lambda\|\vect w\|^2
  =(\mat H\Her\mat H+\lambda\mat I)^{-1}\mat H\Her\vect e_d,
  \qquad\lambda=\begin{cases}\sigma^2&\text{MMSE}\\ \approx0&\text{ZF.}\end{cases}
\end{equation}
All delays are solved at once, since the columns of $(\mat H\Her\mat H+\lambda\mat
I)^{-1}\mat H\Her$ are the solutions for every $d$. The delay with the largest useful gain
$(\mat H\vect w_d)_d$ is selected. The output is then scaled by $1/(\mat H\vect w_d)_d$ to remove
the MMSE bias, and the noise variance becomes $\sigma^2\|\vect w\|^2$.
\codeexcerpt{linear_eq}{Least-squares design of ZF/MMSE equalizers with delay selection (\blk{linear_equalizer}).}

\subsection{Adaptive equalization: LMS}
When the channel is unknown, the equalizer $z[n]=\vect w\Her\vect u[n]$, with
$\vect u[n]=(y[n],y[n-1],\dots)$, can minimize $\E|d[n]-z[n]|^2$ by stochastic gradient
descent. The instantaneous gradient with respect to $\vect w\conj$ is $-\vect u\,e\conj$ with
$e=d-z$, giving the LMS update and its normalized version:
\begin{equation}
  \vect w\leftarrow\vect w+\mu\,\vect u[n]\,e\conj[n],\qquad
  \mu_{\mathrm{NLMS}}=\frac{\tilde\mu}{\|\vect u[n]\|^2+\epsilon},\quad 0<\tilde\mu<2 .
\end{equation}
The reference $d[n]$ is the known training sequence for the first \prm{n_train} symbols and the
nearest constellation point afterwards (decision-directed mode). The residual MSE measured after
training is written to \meta{noise_var}.
\codeexcerpt{lms_loop}{The (N)LMS loop runs over time but processes all frames of the batch at once.}

\begin{keypoint}[ZF versus MMSE]
Zero forcing inverts spectral nulls and amplifies the noise there. MMSE balances ISI against
noise and tends to ZF at high SNR. The decision-feedback equalizer of \cref{sec:dfe} removes
post-cursor ISI without noise enhancement, at the cost of error propagation.
\end{keypoint}

\chapter{Synchronization}\label{ch:sync}

A receiver does not know when symbols start, where the frame begins, or the exact carrier
frequency and phase. \Cref{fig:sync-chain} shows the synchronized single-carrier receiver of
\file{phy/sync.py}, which recovers these quantities in three stages.

\begin{figure}[htbp]
\centering
\begin{tikzpicture}[node distance=4mm and 5mm]
  \node[blockbox] (map) {mapper};
  \node[blockbox,right=of map] (pre) {preamble};
  \node[blockbox,right=of pre] (ps) {pulse\\shaper};
  \node[warmbox,right=of ps,minimum width=26mm] (ch) {delay $\tau$, CFO $\Delta f$,\\phase noise, AWGN};
  \node[blockbox,right=of ch] (mf) {matched\\filter};
  \node[greenbox,below=9mm of mf] (tr) {timing\\recovery};
  \node[greenbox,left=of tr] (fs) {frame\\sync};
  \node[greenbox,left=of fs] (cr) {carrier\\recovery};
  \node[blockbox,left=of cr] (dm) {demapper};
  \draw[sig] (map)--(pre); \draw[sig] (pre)--(ps); \draw[sig] (ps)--(ch); \draw[sig] (ch)--(mf);
  \draw[sig] (mf)--node[right,lbl]{$N_{\mathrm{sps}}\ge2$}(tr); \draw[sig] (tr)--node[above,lbl]{1 sps}(fs);
  \draw[sig] (fs)--(cr); \draw[sig] (cr)--(dm);
  \node[note,below=1mm of tr] {fractional $\tau$};
  \node[note,below=1mm of fs] {integer offset,\\coarse $\Delta f$, phase};
  \node[note,below=1mm of cr] {residual phase};
\end{tikzpicture}
\caption{The synchronized receiver of the \emph{Synchronized QPSK receiver} scenario.}
\label{fig:sync-chain}
\end{figure}

\begin{figure}[htbp]
\centering
\begin{tikzpicture}
\begin{groupplot}[group style={group size=4 by 1,horizontal sep=4mm,y descriptions at=edge left},
  constaxis,width=4.1cm,height=4.1cm]
\nextgroupplot[title={nominal sampling}]
  \addplot[scatterpts=black!50] table[x=i,y=q]{\datafile{sync_raw}};
\nextgroupplot[title={after timing rec.}]
  \addplot[scatterpts=accent] table[x=i,y=q]{\datafile{sync_tr}};
\nextgroupplot[title={after frame sync}]
  \addplot[scatterpts=accent4] table[x=i,y=q]{\datafile{sync_fs}};
\nextgroupplot[title={after the PLL}]
  \addplot[scatterpts=accent3] table[x=i,y=q]{\datafile{sync_cr}};
\end{groupplot}
\end{tikzpicture}
\caption{QPSK through a delay of 9.37 samples, a 3\,kHz CFO at 1\,MBd, 10\,Hz phase noise and
12\,dB $E_b/N_0$. Panels: sampled at the nominal instants (ISI and rotation), after Gardner
timing recovery (rotating ring), after preamble-based CFO and phase correction (residual
spread), after the decision-directed PLL.}
\label{fig:sync-const}
\end{figure}

% ------------------------------------------------------------------
\section{The preamble}
\blk{preamble} prepends $L$ known QPSK symbols made of two identical halves of length
$h=L/2$, the Schmidl--Cox structure~\cite{schmidl1997}, and appends \prm{guard} zero symbols.
The repetition makes the timing metric insensitive to the CFO and yields a frequency estimate
(\cref{sec:frame-sync}). The guard gives the receiver slack for unknown delays.

% ------------------------------------------------------------------
\section{Symbol timing recovery}
After the matched filter the waveform is sampled at $N_{\mathrm{sps}}$ samples per symbol, but
the optimum instants $mT+\tau$ fall between samples. Timing recovery has three parts:
\begin{itemize}
  \item an \emph{interpolator} that evaluates the waveform at arbitrary instants;
  \item a \emph{timing error detector} (TED);
  \item a \emph{loop filter} that turns detector outputs into corrections.
\end{itemize}

\subsection{Interpolation}
A cubic Lagrange (Farrow) interpolator evaluates $x(i+\mu)$ for $0\le\mu<1$ from four
neighbours:
\begin{equation}
  x(i+\mu)=c_0+c_1\mu+c_2\mu^2+c_3\mu^3,\quad
  \begin{aligned}
  c_0&=x_0, & c_1&=-\tfrac13x_{-1}-\tfrac12x_0+x_1-\tfrac16x_2,\\
  c_2&=\tfrac12x_{-1}-x_0+\tfrac12x_1, & c_3&=-\tfrac16x_{-1}+\tfrac12x_0-\tfrac12x_1+\tfrac16x_2 .
  \end{aligned}
\end{equation}
\codeexcerpt{cubic}{Farrow-structure cubic interpolation (\code{cubic_interpolate}).}

\subsection{Feedforward start: the Oerder--Meyr estimator}
A linearly modulated signal is cyclostationary. The mean of $|x(t)|^2$ is periodic in $T$, and
its first Fourier coefficient has a phase proportional to the timing offset. Over $K$ symbols,
\begin{equation}
  \hat\varepsilon=-\frac{1}{2\pi}\arg\sum_{k=0}^{KN_{\mathrm{sps}}-1}|x[k]|^2e^{-j2\pi k/N_{\mathrm{sps}}}
  \qquad(\text{in symbols, modulo }1),
\end{equation}
which needs $N_{\mathrm{sps}}\ge3$~\cite{oerder1988}. The block uses it to initialize the
loop, which removes the long acquisition (and the risk of hang-up) of a loop started from an
arbitrary phase.
\codeexcerpt{om_estimate}{Oerder--Meyr initial estimate.}

\subsection{The Gardner detector}
Gardner's TED~\cite{gardner1986} uses two samples per symbol, the current strobe and the
midpoint between strobes:
\begin{equation}
  e[m]=\Re\Big\{x\conj\big((m-\tfrac12)T+\hat\tau\big)\,\big[x(mT+\hat\tau)-x((m-1)T+\hat\tau)\big]\Big\}.
\end{equation}
It is independent of the carrier phase and works before carrier recovery. Its mean response
to a timing error, the \emph{S-curve}, is zero at the correct instant and has slope $K_d$
there (\cref{fig:scurve}). With the convention of the code, $e>0$ means the strobes are late.

\begin{figure}[htbp]
\centering
\begin{tikzpicture}
\begin{axis}[width=9.5cm,height=5.4cm,xlabel={timing error $\tau/T$},ylabel={$\E\,e$},
  xmin=-0.5,xmax=0.5,legend pos=north west]
  \addplot[th=accent] table[x=tau,y=b20]{\datafile{gardner_scurve}};
  \addplot[th=accent3] table[x=tau,y=b35]{\datafile{gardner_scurve}};
  \addplot[th=accent2] table[x=tau,y=b60]{\datafile{gardner_scurve}};
  \legend{$\beta=0.2$,$\beta=0.35$,$\beta=0.6$}
\end{axis}
\end{tikzpicture}
\caption{Measured S-curves of the Gardner detector for QPSK with RRC pulses (8 samples per symbol,
matched-filtered output of the real blocks, interpolated with the same cubic interpolator). The
slope at the origin grows with the roll-off; for $\beta=0.35$ it is close to 1, the default
\prm{detector_gain}.}
\label{fig:scurve}
\end{figure}

\subsection{The second-order loop}\label{sec:loop}
\begin{figure}[htbp]
\centering
\begin{tikzpicture}[node distance=5mm and 9mm]
  \node[blockbox] (ted) {detector\\$K_d$};
  \node[blockbox,right=of ted] (kp) {$K_p$};
  \node[blockbox,below=of kp] (ki) {$K_i\dfrac{1}{1-z^{-1}}$};
  \node[sumnode,right=of kp] (s) {$+$};
  \node[blockbox,right=of s] (nco) {NCO\\$\dfrac{1}{1-z^{-1}}$};
  \node[blockbox,right=of nco] (int) {interpolator};
  \draw[sig] (ted) -- node[above,lbl]{$e$} (kp);
  \draw[sig] (ted.south) |- (ki);
  \draw[sig] (kp) -- (s); \draw[sig] (ki) -| (s);
  \draw[sig] (s) -- (nco); \draw[sig] (nco) -- node[above,lbl]{$\hat\tau$} (int);
  \draw[sig] (int.north) -- ++(0,6mm) -| node[pos=0.25,above,lbl]{strobes} (ted.north);
\end{tikzpicture}
\caption{Proportional-plus-integral loop used for timing (and, with a phase detector, for the
PLL).}
\end{figure}
A proportional-plus-integral loop filter makes the loop second order. It tracks a constant
offset with zero error and a frequency (clock drift) with zero steady-state error. The gains
follow from the normalized noise bandwidth $B_nT$ and the damping $\zeta$~\cite{mengali1997}:
\begin{equation}
  \theta=\frac{B_nT}{\zeta+\frac{1}{4\zeta}},\qquad
  K_p=\frac{4\zeta\theta}{(1+2\zeta\theta+\theta^2)K_d},\qquad
  K_i=\frac{4\theta^2}{(1+2\zeta\theta+\theta^2)K_d} .
\end{equation}
\codeexcerpt{loop_gains}{Second-order loop gains.}
\codeexcerpt{gardner_loop}{The timing loop: interpolate, detect, filter, advance.}
A small $B_nT$ averages noise better but tracks more slowly; $0.01$ (the default) is a common
compromise.

% ------------------------------------------------------------------
\section{Frame synchronization and coarse CFO}\label{sec:frame-sync}
At symbol rate, with a residual CFO of $\epsilon$ cycles per symbol, the two halves of the
preamble differ only by the phase $2\pi\epsilon h$. The Schmidl--Cox metric is evaluated at
every candidate start $d$:
\begin{equation}
  P(d)=\sum_{k=0}^{h-1}r\conj[d+k]\,r[d+k+h],\qquad
  M(d)=\frac{|P(d)|^2}{\big(\sum_k|r[d+k+h]|^2\big)^2},\qquad
  \hat\epsilon=\frac{\arg P(\hat d)}{2\pi h}.
\end{equation}
The estimate is unambiguous for $|\epsilon|<1/(2h)$. Because $M(d)$ has a plateau, its maximum is
used only for the CFO estimate. After derotation, the frame start is located by
cross-correlation with the known preamble,
$\hat d=\argmax_d\big|\sum_kr[d+k]\,p\conj[k]\big|$. The CFO is then refined at the exact
position, and the common phase is taken from the correlation peak. Finally the payload symbols
are extracted.
\codeexcerpt{frame_sync}{Schmidl--Cox CFO, correlation timing and refinement (\blk{frame_sync}).}

% ------------------------------------------------------------------
\section{Carrier phase tracking}
The residual phase (CFO residue and phase noise) is tracked by a decision-directed PLL. For
$y=s\,e^{j\Delta\phi}+w$ and a correct decision $\hat s$,
\begin{equation}
  e=\frac{\Im\{y\,\hat s\conj\}}{|\hat s|^2}=\sin\Delta\phi+\text{noise}\approx\Delta\phi ,
\end{equation}
so the detector gain is 1. The same second-order loop as above updates the phase estimate
every symbol. Like every decision-directed scheme it has an $M$-fold phase ambiguity, which the
preamble phase estimate resolves beforehand.
\codeexcerpt{dd_pll}{Decision-directed PLL (\blk{carrier_recovery}).}

\begin{pitfall}
The PLL bandwidth must exceed the phase-noise linewidth relative to the symbol rate. It must
also be small enough that decision errors at low SNR do not make it slip cycles. A cycle slip
rotates all subsequent decisions by $2\pi/M$ and appears as a burst of errors until the next
preamble. Differential encoding or pilot symbols are the classic remedies.
\end{pitfall}

\chapter{OFDM}\label{ch:ofdm}

Orthogonal frequency-division multiplexing turns a frequency-selective channel into many
parallel flat channels, each equalized by a single complex division. This chapter derives the
property, describes the subcarrier layout and the pilot-based receiver of
\file{phy/waveform/ofdm.py}, and examines the two weaknesses of OFDM: synchronization and
peak-to-average power ratio.

% ------------------------------------------------------------------
\section{Cyclic prefix and circular convolution}
An OFDM symbol carries $N$ frequency-domain values $X_k$ (the subcarrier symbols). The
transmitter computes the unitary inverse DFT and prepends the last $N_{\mathrm{cp}}$ samples
(the \emph{cyclic prefix}):
\begin{equation}
  x[n]=\frac{1}{\sqrt N}\sum_{k=0}^{N-1}X_k\,e^{j2\pi kn/N},\qquad n=-N_{\mathrm{cp}},\dots,N-1 .
\end{equation}
\begin{proposition}
If the channel impulse response has length $L\le N_{\mathrm{cp}}+1$, then after removing the
prefix and applying the unitary DFT,
\begin{equation}
  Y_k=H_kX_k+W_k,\qquad H_k=\sum_{l=0}^{L-1}h_l\,e^{-j2\pi kl/N},\qquad W_k\sim\CN(0,\sigma^2)\ \text{i.i.d.}
  \label{eq:ofdm-model}
\end{equation}
\end{proposition}
\begin{proof}
Thanks to the prefix, the $N$ retained received samples equal the \emph{circular} convolution of
$x$ with $h$: $\vect y=\mat C\vect x+\vect w$, with $\mat C$ circulant. Every circulant matrix is
diagonalized by the DFT, $\mat C=\mat F\Her\diag(H_0,\dots,H_{N-1})\mat F$ with $\mat F$ the
unitary DFT matrix. Hence $\mat F\vect y=\diag(H)\mat F\vect x+\mat F\vect w$, and
$\mat F\vect w$ is white with the same variance because $\mat F$ is unitary.
\end{proof}

\begin{figure}[htbp]
\centering
\begin{tikzpicture}[x=1cm,y=1cm]
  \foreach \s in {0,1,2} {
    \begin{scope}[xshift=\s*4.3cm]
      \draw[fill=kgrid!35,draw=kgrid] (0,0) rectangle (0.9,0.8);
      \draw[fill=kgrid!10,draw=kgrid] (0.9,0) rectangle (4.1,0.8);
      \fill[kgrid!35] (3.2,0) rectangle (4.1,0.8);
      \node[lbl] at (0.45,0.4) {CP};
      \node[lbl] at (2.05,0.4) {$N$ samples};
      \draw[->,kgrid!80!black,thin] (3.65,0.85) to[bend right=30] (0.45,0.85);
    \end{scope}}
  \draw[decorate,decoration={brace,mirror,amplitude=4pt}] (0,-0.1) -- node[below=4pt,lbl]{$N+N_{\mathrm{cp}}$} (4.1,-0.1);
  \draw[accent2,thick] (5.2,-0.55) -- (8.4,-0.55) node[right,lbl]{FFT window};
  \draw[accent2] (5.2,-0.45)--(5.2,-0.65) (8.4,-0.45)--(8.4,-0.65);
\end{tikzpicture}
\caption{OFDM symbols with cyclic prefix. Any FFT window starting inside the prefix (after the
channel transient) sees a cyclically shifted symbol, which only multiplies $X_k$ by a linear phase.}
\label{fig:ofdm-cp}
\end{figure}

\begin{incode}
\blk{ofdm_modulator} builds an \code{OFDMLayout}, maps data and pilots onto the grid, applies
\code{np.fft.ifft(..., norm="ortho")} and adds the prefix. The unitary transform keeps the
noise variance per subcarrier equal to the noise variance per time sample.

It sets \meta{sps}$=N/N_{\mathrm{act}}$ and
\meta{power}$=(N_{\mathrm{data}}P_d+N_{\mathrm{pil}})/N$. By \cref{eq:esn0} the AWGN block in
$E_s/N_0$ mode then produces $\sigma^2=1/(E_s/N_0)$, the per-subcarrier SNR. Cyclic-prefix and
pilot overheads are deliberately not counted in $E_s$.
\end{incode}
\codeexcerpt{ofdm_mod}{Resource mapping, unitary IFFT and cyclic prefix.}

% ------------------------------------------------------------------
\section{Subcarrier layout}
$N_{\mathrm{act}}$ subcarriers are used around DC. DC itself is left empty, as in Wi-Fi and LTE,
and the remaining edge subcarriers act as guard bands. Every \prm{pilot_spacing}-th active
subcarrier carries a known pilot. The last active subcarrier is always a pilot, so that
interpolation never extrapolates. Pilot values are a fixed pseudo-random QPSK sequence;
identical pilots would add up coherently into a periodic time-domain peak.

\begin{figure}[htbp]
\centering
\begin{tikzpicture}[x=0.36cm,y=1cm]
  % active k = -13..-1, 1..13; pilots where the active index is a multiple of 6, plus the last one
  \foreach \k in {-16,-15,-14,14,15,16} \draw[black!25] (\k,0)--(\k,0.25);
  \draw[black!60] (0,0)--(0,0.25); \node[lbl,below] at (0,0) {DC};
  \foreach \k in {-12,-11,-10,-9,-8,-6,-5,-4,-3,-2,1,2,3,4,5,7,8,9,10,11}
    \fill[accent3!70] (\k-0.3,0) rectangle (\k+0.3,0.65);
  \foreach \k in {-13,-7,-1,6,12,13}
    \fill[accent2] (\k-0.3,0) rectangle (\k+0.3,0.95);
  \draw[->] (-17,0) -- (17.5,0) node[right,lbl]{$k$};
  \fill[accent2] (-15,1.45) rectangle (-14.4,1.7); \node[lbl,right] at (-14.2,1.57) {pilot};
  \fill[accent3!70] (-9,1.45) rectangle (-8.4,1.7); \node[lbl,right] at (-8.2,1.57) {data};
  \draw[black!40] (-3,1.45)--(-3,1.7); \node[lbl,right] at (-2.8,1.57) {unused (guard)};
\end{tikzpicture}
\caption{Layout with $N_{\mathrm{act}}=26$ active subcarriers and a pilot every 6 (the first and the
last active subcarriers are pilots).}
\label{fig:ofdm-layout}
\end{figure}

% ------------------------------------------------------------------
\section{Channel estimation and equalization}
\subsection{Least squares on pilots}
At pilot subcarriers \cref{eq:ofdm-model} gives the LS estimate $\hat H_p=Y_p/X_p$, with error
variance $\sigma^2$ since pilots have unit modulus. The channel at data subcarriers is
obtained by linear interpolation, implemented as a fixed matrix from pilot to active
subcarriers. Optionally the estimates are averaged over the OFDM symbols of a frame, which
lowers the estimation noise when the channel is static.

\subsection{Removing the timing phase slope}
A timing offset of $\delta$ samples inside the prefix multiplies $H_k$ by $e^{-j2\pi k\delta/N}$.
This fast phase rotation across subcarriers is poorly approximated by linear interpolation of
the real and imaginary parts. The estimator therefore:
\begin{enumerate}
  \item estimates the dominant slope from products of adjacent pilots,
  $\hat s=\arg\sum_p\hat H_{p+1}\hat H_p\conj\,/\,\Delta_p$;
  \item removes it;
  \item interpolates the now slowly varying response;
  \item restores the slope.
\end{enumerate}
\codeexcerpt{ls_slope}{LS estimation with phase-slope removal (\blk{ofdm_channel_estimator}).}

\begin{figure}[htbp]
\centering
\begin{tikzpicture}
\begin{groupplot}[group style={group size=1 by 2,vertical sep=6mm,x descriptions at=edge bottom},
  width=12cm,height=4.4cm,xmin=-51,xmax=51,xlabel={subcarrier index $k$}]
\nextgroupplot[ylabel={$|H_k|$ [dB]},legend pos=south west,legend columns=3]
  \addplot[th=black!70] table[x=k,y=true_db]{\datafile{ofdm_chest}};
  \addplot[th=accent,dashed] table[x=k,y=ls_db]{\datafile{ofdm_chest}};
  \addplot[sim=accent2,mark size=1.4pt] table[x=k,y=db]{\datafile{ofdm_pilots}};
  \legend{true channel,LS + interpolation,pilot estimates}
\nextgroupplot[ylabel={$\arg H_k$ [rad]}]
  \addplot[th=black!70] table[x=k,y=true_ph]{\datafile{ofdm_chest}};
  \addplot[th=accent,dashed] table[x=k,y=ls_ph]{\datafile{ofdm_chest}};
\end{groupplot}
\end{tikzpicture}
\caption{One OFDM symbol over the EVA channel ($N=128$, 100 active subcarriers, pilot every 6,
$15$\,kHz spacing, $E_s/N_0=15$\,dB): true response from the channel's \code{csi} (\prm{perfect}
estimator) and LS estimate.}
\label{fig:ofdm-chest}
\end{figure}

\subsection{One-tap equalization}
\blk{ofdm_equalizer} extracts the data subcarriers and divides by the estimate,
$\hat X_k=Y_k/\hat H_k$. It writes the per-subcarrier noise variance $\sigma^2/|\hat H_k|^2$ to
\meta{noise_var}. With coding and interleaving across subcarriers, the soft demapper then gives
low weight to faded subcarriers. This is how OFDM obtains frequency diversity (see the Wi-Fi and
LTE scenarios).

% ------------------------------------------------------------------
\section{Blind timing and frequency synchronization}\label{sec:ofdm-sync}
The cyclic prefix is a copy of the end of each symbol, so the received signal is correlated
with itself $N$ samples later over the prefix. Van de Beek \emph{et al.}~\cite{vandebeek1997}
use this redundancy. With a CFO of $\varepsilon$ subcarrier spacings, $r[k+N]\approx r[k]e^{j2\pi\varepsilon}$
inside the prefix, so the correlation
\begin{equation}
  \gamma(\theta)=\sum_{m}\sum_{k=\theta+mL}^{\theta+mL+N_{\mathrm{cp}}-1}r[k]\,r\conj[k+N],\qquad L=N+N_{\mathrm{cp}},
\end{equation}
peaks in magnitude at the symbol start, and its phase gives the fractional CFO:
\begin{equation}
  \hat\theta=\argmax_\theta|\gamma(\theta)|,\qquad \hat\varepsilon=-\frac{1}{2\pi}\arg\gamma(\hat\theta),\quad|\varepsilon|<\tfrac12 .
\end{equation}
The block accumulates $\gamma$ over all OFDM symbols of the frame with cumulative sums, which
makes the search over $\theta$ linear in the frame length. It then starts the FFT window
\prm{backoff} samples early, inside the prefix, to stay clear of the channel's delay spread.
The resulting linear phase is absorbed by the channel estimator. Compared with the full ML
metric of van de Beek, the energy correction term is omitted; it matters only at low SNR.
\codeexcerpt{ofdm_sync}{CP correlation, timing and fractional CFO (\blk{ofdm_sync}).}

\begin{pitfall}
An integer CFO (a multiple of the subcarrier spacing) shifts the whole grid and is
invisible to the CP correlation. It must be resolved with a preamble or pilots, which this
block does not do.
\end{pitfall}

% ------------------------------------------------------------------
\section{Peak-to-average power ratio}\label{sec:papr}
Each time sample is a sum of $N_{\mathrm{act}}$ independent terms. By the central limit theorem
it is nearly $\CN(0,P)$, so $|x|^2/P$ is approximately exponential. Treating the $N$ Nyquist-rate
samples of a symbol as independent gives the classic approximation
\begin{equation}
  P\big(\mathrm{PAPR}>\gamma\big)\approx1-\big(1-e^{-\gamma}\big)^{N_{\mathrm{act}}} .
\end{equation}

\begin{figure}[htbp]
\centering
\begin{tikzpicture}
\begin{semilogyaxis}[width=10cm,height=6cm,xlabel={$\gamma$ [dB]},ylabel={$P(\mathrm{PAPR}>\gamma)$},
  xmin=4,xmax=13,ymin=1e-3,ymax=1,legend pos=south west]
  \addplot[sim=accent] table[x=papr,y=sim64]{\datafile{papr_ccdf}};
  \addplot[th=accent] table[x=papr,y=th64]{\datafile{papr_ccdf}};
  \addplot[sim={accent3,mark=square*}] table[x=papr,y=sim256]{\datafile{papr_ccdf}};
  \addplot[th=accent3] table[x=papr,y=th256]{\datafile{papr_ccdf}};
  \addplot[sim={accent2,mark=triangle*}] table[x=papr,y=sim1024]{\datafile{papr_ccdf}};
  \addplot[th=accent2] table[x=papr,y=th1024]{\datafile{papr_ccdf}};
  \legend{$N{=}64$ sim.,approx.,$N{=}256$ sim.,approx.,$N{=}1024$ sim.,approx.}
\end{semilogyaxis}
\end{tikzpicture}
\caption{CCDF of the per-symbol PAPR of QPSK-OFDM at the output of \blk{ofdm_modulator} (no
oversampling, no prefix) and the approximation above. A PAPR of 10--12\,dB forces the
power-amplifier back-off discussed in \cref{ch:impairments}.}
\label{fig:papr}
\end{figure}

\begin{remark}
The modulator produces Nyquist-rate samples, which miss part of the analog peaks; with
oversampling the PAPR is typically about 0.5\,dB higher.
\end{remark}

\chapter{Multiple antennas}\label{ch:mimo}

% ------------------------------------------------------------------
\section{The flat-fading MIMO channel}
With $N_t$ transmit and $N_r$ receive antennas and flat fading, one channel use is
\begin{equation}
  \vect y=\mat H\vect x+\vect w,\qquad \mat H=\frac{1}{\sqrt{N_t}}\mat G,\quad G_{ij}\sim\CN(0,1),\quad \vect w\sim\CN(\vect0,\sigma^2\mat I).
\end{equation}
The normalization by $\sqrt{N_t}$ keeps the total transmitted power independent of the number of
antennas. Every receive antenna then sees unit average signal power, and the SNR set in the
\blk{awgn} block is the SNR \emph{per receive antenna}. Antenna correlation follows the
Kronecker model with exponential correlation matrices,
\begin{equation}
  \mat H=\frac{1}{\sqrt{N_t}}\,\mat R_r^{1/2}\mat G\,\mat R_t^{1/2},\qquad [\mat R]_{ij}=\rho^{|i-j|}.
\end{equation}

\begin{incode}
Multi-stream signals have shape $(B,N_t,n)$.
\begin{itemize}
  \item \blk{stream_demux} distributes symbols round-robin over the streams and sets
  \meta{spatial_rate}$=N_t$, so that $E_b/N_0$ accounts for the $N_t$ symbols sent per channel
  use (\cref{eq:ebn0}).
  \item \blk{mimo_channel} draws $\mat H$ per frame (\prm{block}) or per symbol (\prm{fast}) and
  publishes it on \code{csi} (\meta{csi_type}\code{="mimo"}).
  \item \blk{mimo_detector} and \blk{alamouti_decoder} use it, and \blk{stream_mux} restores a
  single stream.
\end{itemize}
\end{incode}

% ------------------------------------------------------------------
\section{Detection for spatial multiplexing}
\subsection{Zero forcing}
$\mat W_{\mathrm{ZF}}=(\mat H\Her\mat H)^{-1}\mat H\Her$ separates the streams exactly,
$\hat{\vect x}=\vect x+\mat W\vect w$. Stream $k$ then has noise variance
$\sigma^2[(\mat H\Her\mat H)^{-1}]_{kk}$, which becomes large when $\mat H$ is ill-conditioned. For
$N_t=N_r$ each stream has only one degree of diversity.

\subsection{MMSE and its bias}
The linear MMSE estimator is $\mat W=(\mat H\Her\mat H+\sigma^2\mat I)^{-1}\mat H\Her$ (unit-power
symbols). Its output is biased: $\E[\hat x_k|x_k]=b_kx_k$ with $b_k=(\mat W\mat H)_{kk}<1$.
\begin{proposition}
For the MMSE filter, the residual interference plus noise of stream $k$ has variance
$b_k(1-b_k)$. After unbiasing ($\hat x_k/b_k$), the noise variance is therefore $(1-b_k)/b_k$,
and the post-detection SINR is $b_k/(1-b_k)$.
\end{proposition}
\begin{proof}
For the MMSE filter, the error covariance is
$\mat E=\E[(\hat{\vect x}-\vect x)(\hat{\vect x}-\vect x)\Her]=\mat I-\mat W\mat H$, so the MSE of
stream $k$ is $1-b_k$. Write $\hat x_k=b_kx_k+\nu_k$ with $\nu_k$ uncorrelated with $x_k$ (an
orthogonality property of the MMSE estimator). Then
$\mathrm{MSE}_k=(1-b_k)^2+\Var\nu_k=1-b_k$, which gives $\Var\nu_k=b_k(1-b_k)$.
\end{proof}
The block outputs the unbiased symbols and writes $(1-b_k)/b_k$ to \meta{noise_var}, so the
soft demapper receives correctly scaled LLRs.
\codeexcerpt{mimo_linear}{ZF and unbiased MMSE detection with post-detection noise variances.}

\subsection{Maximum likelihood}
$\hat{\vect x}=\argmin_{\vect x\in\mathcal S^{N_t}}\|\vect y-\mat H\vect x\|^2$ is optimal and
achieves diversity $N_r$ per stream. The block searches all $M^{N_t}$ candidates exhaustively
(limited to 4096) and outputs hard decisions.

\begin{figure}[htbp]
\centering
\begin{tikzpicture}
\begin{semilogyaxis}[berplot,width=10cm,height=6.5cm,xmin=0,xmax=20,ymin=1e-5,ymax=0.5]
  \addplot[simline=accent] table[x=ebn0,y=zf]{\datafile{mimo_smux}};
  \addplot[simline={accent3,mark=square*}] table[x=ebn0,y=mmse]{\datafile{mimo_smux}};
  \addplot[simline={accent2,mark=triangle*}] table[x=ebn0,y=ml]{\datafile{mimo_smux}};
  \legend{ZF,MMSE,ML}
\end{semilogyaxis}
\end{tikzpicture}
\caption{$2\times2$ spatial multiplexing of QPSK over fast Rayleigh fading (simulated).
ZF and MMSE have diversity 1, the same slope as single-antenna Rayleigh fading, and MMSE gains
about 2\,dB. ML reaches diversity 2.}
\label{fig:smux}
\end{figure}

% ------------------------------------------------------------------
\section{Diversity}\label{sec:diversity}
\subsection{Receive diversity: maximal-ratio combining}
With one transmit antenna, ZF and MMSE reduce to maximal-ratio combining,
$\hat x=\vect h\Her\vect y/\|\vect h\|^2$, whose SNR is the sum of the branch SNRs. For BPSK with
$L$ i.i.d. Rayleigh branches of mean SNR $\bar\gamma$ each~\cite{proakis2008},
\begin{equation}
  \bar P_b=p^L\sum_{l=0}^{L-1}\binom{L-1+l}{l}(1-p)^l,\qquad
  p=\frac{1-\mu}{2},\quad \mu=\sqrt{\frac{\bar\gamma}{1+\bar\gamma}},
  \label{eq:mrc}
\end{equation}
and asymptotically $\bar P_b\propto\bar\gamma^{-L}$: the slope of the curve is the diversity order.

\subsection{Transmit diversity: the Alamouti code}
Two symbols are sent over two antennas in two symbol periods~\cite{alamouti1998}:
\begin{equation}
  \begin{array}{c|cc}
    & \text{antenna 1} & \text{antenna 2}\\ \hline
    t_1 & s_1 & s_2\\
    t_2 & -s_2\conj & s_1\conj
  \end{array}
  \qquad\Longrightarrow\qquad
  \begin{aligned}
    \tilde s_1&=h_1\conj r_1+h_2r_2\conj=(|h_1|^2+|h_2|^2)\,s_1+\tilde w_1,\\
    \tilde s_2&=h_2\conj r_1-h_1r_2\conj=(|h_1|^2+|h_2|^2)\,s_2+\tilde w_2 .
  \end{aligned}
\end{equation}
The code is orthogonal, so the linear combiner decouples the symbols and gives full diversity
2 without channel knowledge at the transmitter. The decoder sums over any number of receive
antennas. Because the transmit power is split between the antennas, the SNR is $3$\,dB below
two-branch MRC: $\bar P_b^{\text{Alamouti}}(\bar\gamma)=\bar P_b^{\text{MRC},L=2}(\bar\gamma/2)$.
\codeexcerpt{alamouti_dec}{Alamouti combining over $N_r$ receive antennas.}

\begin{figure}[htbp]
\centering
\begin{tikzpicture}
\begin{semilogyaxis}[berplot,width=11cm,height=7cm,xmin=0,xmax=22,ymin=1e-5,ymax=0.5,legend columns=2]
  \addplot[th=black!50,dashed] table[x=ebn0,y=awgn]{\datafile{diversity_theory}};
  \addlegendentry{AWGN}
  \addplot[th=accent] table[x=ebn0,y=siso]{\datafile{diversity_theory}};
  \addlegendentry{Rayleigh, theory}
  \addplot[sim=accent] table[x=ebn0,y=siso]{\datafile{diversity_sim}};
  \addlegendentry{SISO simulated}
  \addplot[th=accent3] table[x=ebn0,y=mrc2]{\datafile{diversity_theory}};
  \addlegendentry{MRC $L{=}2$, theory}
  \addplot[sim={accent3,mark=square*}] table[x=ebn0,y=mrc2]{\datafile{diversity_sim}};
  \addlegendentry{$1{\times}2$ simulated}
  \addplot[th=accent2] table[x=ebn0,y=alamouti]{\datafile{diversity_theory}};
  \addlegendentry{Alamouti, theory}
  \addplot[sim={accent2,mark=triangle*}] table[x=ebn0,y=alamouti]{\datafile{diversity_sim}};
  \addlegendentry{$2{\times}1$ simulated}
\end{semilogyaxis}
\end{tikzpicture}
\caption{BPSK over Rayleigh fading: single antenna (\cref{eq:rayleigh-ber}), $1\times2$ MRC and
$2\times1$ Alamouti (\cref{eq:mrc}). The $E_b/N_0$ axis is per receive antenna. Diversity 2
doubles the slope; Alamouti runs parallel to MRC, 3\,dB to the right.}
\label{fig:diversity}
\end{figure}

\chapter{Channel coding}\label{ch:coding}

Channel coding adds structured redundancy so that the receiver can detect and correct errors.
\Cref{fig:coding-chain} shows the coding blocks of \file{phy/coding/}. They can be combined
freely. Every encoder multiplies the \meta{code_rate} metadata by its own rate, including tail
bits and CRC overhead, so the \blk{awgn} block always refers $E_b/N_0$ to the
\emph{information} bits (\cref{eq:ebn0}).

\begin{figure}[htbp]
\centering
\begin{tikzpicture}[node distance=3mm and 3.5mm,every node/.style={font=\sffamily\scriptsize}]
  \node[blockbox,minimum width=12mm] (s) {source};
  \node[greenbox,right=of s,minimum width=12mm] (crc) {CRC};
  \node[greenbox,right=of crc,minimum width=12mm] (scr) {scrambler};
  \node[greenbox,right=of scr,minimum width=12mm] (enc) {FEC\\encoder};
  \node[greenbox,right=of enc,minimum width=12mm] (pun) {puncturer};
  \node[greenbox,right=of pun,minimum width=12mm] (il) {interleaver};
  \node[blockbox,right=of il,minimum width=12mm] (map) {mapper};
  \node[warmbox,below=7mm of map,minimum width=12mm] (ch) {channel};
  \node[blockbox,left=of ch,minimum width=12mm] (dm) {soft\\demapper};
  \node[greenbox,left=of dm,minimum width=12mm] (dil) {de-inter-\\leaver};
  \node[greenbox,left=of dil,minimum width=12mm] (dpun) {depunc-\\turer};
  \node[greenbox,left=of dpun,minimum width=12mm] (dec) {FEC\\decoder};
  \node[greenbox,left=of dec,minimum width=12mm] (dscr) {descr.};
  \node[greenbox,left=of dscr,minimum width=12mm] (dcrc) {CRC\\check};
  \draw[sig,kbits] (s)--(crc); \draw[sig,kbits] (crc)--(scr); \draw[sig,kbits] (scr)--(enc);
  \draw[sig,kbits] (enc)--(pun); \draw[sig,kbits] (pun)--(il); \draw[sig,kbits] (il)--(map);
  \draw[sig] (map)--(ch); \draw[sig] (ch)--(dm); \draw[sig,kllr] (dm)--(dil); \draw[sig,kllr] (dil)--(dpun);
  \draw[sig,kllr] (dpun)--(dec); \draw[sig,kbits] (dec)--(dscr); \draw[sig,kbits] (dscr)--(dcrc);
\end{tikzpicture}
\caption{A complete coded chain. Purple edges carry bits, pink edges LLRs. The scrambler and the
interleaver are their own inverses (\prm{inverse}\code{=true} for the interleaver).}
\label{fig:coding-chain}
\end{figure}

% ------------------------------------------------------------------
\section{Cyclic redundancy check}
A CRC with generator polynomial $g(x)$ of degree $r$ appends the remainder
$p(x)=x^r u(x)\bmod g(x)$ to the message $u(x)$; the receiver recomputes it and compares.
Because polynomial division over GF(2) is linear, the parity is a linear function of the
message, $\vect p=\vect u\mat P\bmod 2$. Row $i$ of $\mat P$ is the remainder of the monomial
$x^{k-1-i+r}$. Computing all of them once turns CRC encoding and checking of a whole batch
into a single integer matrix product.
\codeexcerpt{crc_matrix}{The CRC parity matrix, built once per frame length and cached.}
A CRC of degree $r$ detects all error bursts of length $\le r$. It misses a random error
pattern with probability $\approx2^{-r}$. \blk{crc_decoder} reports \code{bler}, the block error
rate detected by the CRC, and outputs a validity flag per frame. CRC-8, CRC-16-CCITT, LTE
CRC-24A and IEEE 802.3 CRC-32 are available.

% ------------------------------------------------------------------
\section{LFSR sequences and scrambling}
A linear feedback shift register produces $s[k]=\bigoplus_{t\in\mathcal T}s[k-t]$. When the
feedback polynomial $1+\sum_{t\in\mathcal T}x^t$ is primitive of degree $m$, the sequence is an
$m$-sequence of period $2^m-1$ with nearly ideal autocorrelation. The \blk{prbs} source
(PRBS-7 to PRBS-31, ITU-T O.150) and the additive \blk{scrambler} (default $x^7+x^4+1$, IEEE
802.11) are built on it. Since every output bit depends only on bits at least
$\min\mathcal T$ positions back, chunks of that many bits can be computed with vector operations:
\codeexcerpt{lfsr}{Vectorized LFSR generation (\file{dsp/lfsr.py}).}
Scrambling is a XOR with the sequence and is its own inverse. On LLRs, descrambling flips the
sign of the affected values. The self-synchronizing $x^{58}+x^{39}+1$ scrambler of 64B/66B is
described in \cref{sec:line-codes}.

% ------------------------------------------------------------------
\section{Interleaving}
Fading and decoders both produce errors in bursts, while most codes are designed for
independent errors. An interleaver permutes the coded bits, and the de-interleaver applies the
inverse permutation, so a burst on the channel is spread over many codewords. The
\prm{block} interleaver writes row by row into an $R\times(n/R)$ matrix and reads column by
column. The \prm{random} interleaver applies a seeded pseudo-random permutation.

% ------------------------------------------------------------------
\section{Block codes: repetition and Hamming}
The repetition code sends every bit $n$ times. Its soft decoder simply adds the $n$ LLRs. The
Hamming code $(2^m-1,2^m-1-m)$ has as parity-check matrix $\mat H$ all non-zero $m$-bit columns,
arranged so that the code is systematic. A single error in position $j$ produces the syndrome
$\vect s=\mat H\vect e\T=\vect h_j$, which identifies the position; a lookup table maps each
syndrome to the bit to flip. Decoding errors occur with $\ge2$ channel errors; the approximation
\begin{equation}
  P_b\approx\frac1n\sum_{j=2}^{n}(j+1)\binom nj p^j(1-p)^{n-j},\qquad p=\Q\Big(\sqrt{2R\,E_b/N_0}\Big),
\end{equation}
counts the $j$ channel errors plus at most one wrong correction (\code{ber_hamming_hard}).

% ------------------------------------------------------------------
\section{Convolutional codes and the Viterbi algorithm}\label{sec:conv}
\subsection{Encoder and trellis}
A rate-$1/n$ feed-forward convolutional encoder with constraint length $K$ computes $n$ output
bits per input bit,
\begin{equation}
  c_j[t]=\bigoplus_{i=0}^{K-1}g_{j,i}\,u[t-i],\qquad j=1,\dots,n,
\end{equation}
with generators given in octal: $(171,133)_8$ is the industry-standard $K=7$ code of IEEE
802.11, DVB and many others, with free distance $d_{\mathrm{free}}=10$. The encoder state is
the last $K-1$ input bits, giving $2^{K-1}$ states. Appending $K-1$ zeros (\prm{terminate})
returns the trellis to state 0.

\begin{figure}[htbp]
\centering
\begin{tikzpicture}[x=2.1cm,y=0.95cm]
  \foreach \t in {0,...,4} {
    \foreach \s/\lab in {0/00,1/01,2/10,3/11} \node[dot] (n\t\s) at (\t,-\s) {};
    \node[lbl,above] at (\t,0.25) {$t=\t$};}
  \foreach \s/\lab in {0/00,1/01,2/10,3/11} \node[lbl,left] at (-0.1,-\s) {\lab};
  \foreach \t [evaluate=\t as \u using int(\t+1)] in {0,...,3} {
    % input 0 (solid), input 1 (dashed); state = (u[t-1], u[t-2]) MSB = most recent
    \draw[accent,thin] (n\t0)--(n\u0); \draw[accent2,thin,dashed] (n\t0)--(n\u2);
    \draw[accent,thin] (n\t1)--(n\u0); \draw[accent2,thin,dashed] (n\t1)--(n\u2);
    \draw[accent,thin] (n\t2)--(n\u1); \draw[accent2,thin,dashed] (n\t2)--(n\u3);
    \draw[accent,thin] (n\t3)--(n\u1); \draw[accent2,thin,dashed] (n\t3)--(n\u3);}
  \draw[accent3,line width=1.6pt] (n00)--(n12)--(n23)--(n31)--(n40);
  \node[lbl,right] at (4.1,-0.0) {survivor};
  \draw[accent,thin] (4.6,-2.2)--(4.9,-2.2) node[right,lbl]{$u=0$};
  \draw[accent2,thin,dashed] (4.6,-2.8)--(4.9,-2.8) node[right,lbl]{$u=1$};
\end{tikzpicture}
\caption{Trellis of the $K=3$, $(7,5)_8$ code (4 states, MSB = most recent input). The Viterbi
algorithm keeps one survivor path per state and time; the green path is the decoded sequence
\code{1 1 0 0}.}
\label{fig:trellis}
\end{figure}

\subsection{Maximum-likelihood sequence decoding}
Given channel LLRs $L_{t,j}$ for the coded bits, the log-likelihood of a candidate codeword is,
up to a constant,
\begin{equation}
  \ln p(\vect y\,|\,\vect c)=\sum_{t,j}\frac{L_{t,j}}{2}(1-2c_{t,j})+\text{const},
\end{equation}
a \emph{correlation} metric. The Viterbi algorithm maximizes it over all paths of the trellis
by dynamic programming. At each time, for every state it keeps the better of the two incoming
paths (add--compare--select) and stores the decision. A traceback from state 0 (terminated
trellis) recovers the input bits. Path metrics are renormalized every step to avoid overflow.
\codeexcerpt{viterbi}{Viterbi decoding vectorized over states and frames (\code{Trellis.decode}).}
With hard decisions (bit inputs) the metric degenerates to the Hamming distance, losing about
2\,dB (\file{tests/test_physics.py::test_conv_coded_soft_vs_hard}). The asymptotic coding gain
of soft decoding is $10\log_{10}(R\,d_{\mathrm{free}})=7$\,dB for the $K=7$ code.

\subsection{Puncturing}
Deleting coded bits periodically raises the rate without a new decoder. The default pattern
\code{[1,1,1,0,0,1]} turns the rate-1/2 code into rate 3/4, as in IEEE 802.11. The
depuncturer reinserts the deleted positions as erasures ($L=0$, no information) before
decoding. When the original length is not available in the metadata, it is inferred as the
shortest mother length whose pattern keeps exactly the received number of bits.
\codeexcerpt{depuncture}{Erasure insertion (\blk{depuncturer}).}

% ------------------------------------------------------------------
\section{LDPC codes}\label{sec:ldpc}
\subsection{Structure}
A low-density parity-check code~\cite{gallager1962} is the null space of a sparse matrix $\mat H\in\{0,1\}^{m\times n}$:
$\mat H\vect c\T=\vect0$. Its Tanner graph has a variable node per code bit and a check node per
row (\cref{fig:tanner}). LinkStudio builds quasi-cyclic codes in the style of IEEE 802.11n and 5G
NR: $\mat H$ is an array of $z\times z$ circulant permutation matrices (or zero blocks), described
by a base matrix of shifts (\cref{fig:ldpc-base}).

\begin{figure}[htbp]
\centering
\begin{tikzpicture}[x=1.15cm,y=1cm]
  \foreach \i in {1,...,8} \node[circle,draw=accent,fill=soft,minimum size=6mm,inner sep=0pt,font=\scriptsize] (v\i) at (\i,0) {$c_\i$};
  \foreach \j in {1,...,4} \node[rectangle,draw=accent2,fill=softwarm,minimum size=6mm,font=\scriptsize] (c\j) at (2*\j-0.5,2) {$\oplus$};
  \foreach \j/\list in {1/{1,2,4,5},2/{2,3,6,7},3/{1,3,5,8},4/{4,6,7,8}}
    \foreach \i in \list \draw[black!50] (c\j.south)--(v\i.north);
  \draw[->,accent3,thick] (v5.north east) .. controls +(0.4,0.7) .. node[right,lbl]{$L_{v\to c}$} ($(c3.south)+(0.15,0)$);
  \node[lbl,right] at (8.5,0) {variable nodes};
  \node[lbl,right] at (8.5,2) {check nodes};
\end{tikzpicture}
\caption{A Tanner graph. Belief propagation exchanges messages along the edges: variables tell
checks what they believe, checks reply with what the other bits imply.}
\label{fig:tanner}
\end{figure}

\begin{figure}[htbp]
\centering
\begin{tikzpicture}[infoblk/.style={draw=accent,fill=soft,minimum size=5.6mm,inner sep=0pt},
  parblk/.style={draw=accent3,fill=softgreen,minimum size=5.6mm,inner sep=0pt},
  zeroblk/.style={draw=black!12,fill=white,minimum size=5.6mm,inner sep=0pt}]
  % Base matrix of the default QC-LDPC code (k_blocks=8, m_blocks=8, z=48), from QCLDPCCode.H
\node[zeroblk] at (0.00,0.00) {};
\node[zeroblk] at (0.62,0.00) {};
\node[infoblk] at (1.24,0.00) {\tiny 36};
\node[zeroblk] at (1.86,0.00) {};
\node[infoblk] at (2.48,0.00) {\tiny 12};
\node[infoblk] at (3.10,0.00) {\tiny 47};
\node[infoblk] at (3.72,0.00) {\tiny 13};
\node[zeroblk] at (4.34,0.00) {};
\node[parblk] at (4.96,0.00) {\tiny 0};
\node[zeroblk] at (5.58,0.00) {};
\node[zeroblk] at (6.20,0.00) {};
\node[zeroblk] at (6.82,0.00) {};
\node[zeroblk] at (7.44,0.00) {};
\node[zeroblk] at (8.06,0.00) {};
\node[zeroblk] at (8.68,0.00) {};
\node[zeroblk] at (9.30,0.00) {};
\node[zeroblk] at (0.00,-0.62) {};
\node[infoblk] at (0.62,-0.62) {\tiny 12};
\node[zeroblk] at (1.24,-0.62) {};
\node[zeroblk] at (1.86,-0.62) {};
\node[zeroblk] at (2.48,-0.62) {};
\node[zeroblk] at (3.10,-0.62) {};
\node[zeroblk] at (3.72,-0.62) {};
\node[infoblk] at (4.34,-0.62) {\tiny 24};
\node[parblk] at (4.96,-0.62) {\tiny 0};
\node[parblk] at (5.58,-0.62) {\tiny 0};
\node[zeroblk] at (6.20,-0.62) {};
\node[zeroblk] at (6.82,-0.62) {};
\node[zeroblk] at (7.44,-0.62) {};
\node[zeroblk] at (8.06,-0.62) {};
\node[zeroblk] at (8.68,-0.62) {};
\node[zeroblk] at (9.30,-0.62) {};
\node[infoblk] at (0.00,-1.24) {\tiny 39};
\node[infoblk] at (0.62,-1.24) {\tiny 19};
\node[zeroblk] at (1.24,-1.24) {};
\node[infoblk] at (1.86,-1.24) {\tiny 14};
\node[zeroblk] at (2.48,-1.24) {};
\node[zeroblk] at (3.10,-1.24) {};
\node[zeroblk] at (3.72,-1.24) {};
\node[infoblk] at (4.34,-1.24) {\tiny 5};
\node[zeroblk] at (4.96,-1.24) {};
\node[parblk] at (5.58,-1.24) {\tiny 0};
\node[parblk] at (6.20,-1.24) {\tiny 0};
\node[zeroblk] at (6.82,-1.24) {};
\node[zeroblk] at (7.44,-1.24) {};
\node[zeroblk] at (8.06,-1.24) {};
\node[zeroblk] at (8.68,-1.24) {};
\node[zeroblk] at (9.30,-1.24) {};
\node[infoblk] at (0.00,-1.86) {\tiny 6};
\node[zeroblk] at (0.62,-1.86) {};
\node[infoblk] at (1.24,-1.86) {\tiny 25};
\node[infoblk] at (1.86,-1.86) {\tiny 21};
\node[infoblk] at (2.48,-1.86) {\tiny 24};
\node[zeroblk] at (3.10,-1.86) {};
\node[zeroblk] at (3.72,-1.86) {};
\node[zeroblk] at (4.34,-1.86) {};
\node[zeroblk] at (4.96,-1.86) {};
\node[zeroblk] at (5.58,-1.86) {};
\node[parblk] at (6.20,-1.86) {\tiny 0};
\node[parblk] at (6.82,-1.86) {\tiny 0};
\node[zeroblk] at (7.44,-1.86) {};
\node[zeroblk] at (8.06,-1.86) {};
\node[zeroblk] at (8.68,-1.86) {};
\node[zeroblk] at (9.30,-1.86) {};
\node[zeroblk] at (0.00,-2.48) {};
\node[zeroblk] at (0.62,-2.48) {};
\node[zeroblk] at (1.24,-2.48) {};
\node[infoblk] at (1.86,-2.48) {\tiny 5};
\node[zeroblk] at (2.48,-2.48) {};
\node[zeroblk] at (3.10,-2.48) {};
\node[infoblk] at (3.72,-2.48) {\tiny 44};
\node[infoblk] at (4.34,-2.48) {\tiny 14};
\node[zeroblk] at (4.96,-2.48) {};
\node[zeroblk] at (5.58,-2.48) {};
\node[zeroblk] at (6.20,-2.48) {};
\node[parblk] at (6.82,-2.48) {\tiny 0};
\node[parblk] at (7.44,-2.48) {\tiny 0};
\node[zeroblk] at (8.06,-2.48) {};
\node[zeroblk] at (8.68,-2.48) {};
\node[zeroblk] at (9.30,-2.48) {};
\node[zeroblk] at (0.00,-3.10) {};
\node[zeroblk] at (0.62,-3.10) {};
\node[zeroblk] at (1.24,-3.10) {};
\node[zeroblk] at (1.86,-3.10) {};
\node[infoblk] at (2.48,-3.10) {\tiny 9};
\node[infoblk] at (3.10,-3.10) {\tiny 23};
\node[zeroblk] at (3.72,-3.10) {};
\node[zeroblk] at (4.34,-3.10) {};
\node[zeroblk] at (4.96,-3.10) {};
\node[zeroblk] at (5.58,-3.10) {};
\node[zeroblk] at (6.20,-3.10) {};
\node[zeroblk] at (6.82,-3.10) {};
\node[parblk] at (7.44,-3.10) {\tiny 0};
\node[parblk] at (8.06,-3.10) {\tiny 0};
\node[zeroblk] at (8.68,-3.10) {};
\node[zeroblk] at (9.30,-3.10) {};
\node[infoblk] at (0.00,-3.72) {\tiny 45};
\node[infoblk] at (0.62,-3.72) {\tiny 39};
\node[infoblk] at (1.24,-3.72) {\tiny 40};
\node[zeroblk] at (1.86,-3.72) {};
\node[zeroblk] at (2.48,-3.72) {};
\node[zeroblk] at (3.10,-3.72) {};
\node[infoblk] at (3.72,-3.72) {\tiny 25};
\node[zeroblk] at (4.34,-3.72) {};
\node[zeroblk] at (4.96,-3.72) {};
\node[zeroblk] at (5.58,-3.72) {};
\node[zeroblk] at (6.20,-3.72) {};
\node[zeroblk] at (6.82,-3.72) {};
\node[zeroblk] at (7.44,-3.72) {};
\node[parblk] at (8.06,-3.72) {\tiny 0};
\node[parblk] at (8.68,-3.72) {\tiny 0};
\node[zeroblk] at (9.30,-3.72) {};
\node[zeroblk] at (0.00,-4.34) {};
\node[zeroblk] at (0.62,-4.34) {};
\node[zeroblk] at (1.24,-4.34) {};
\node[zeroblk] at (1.86,-4.34) {};
\node[zeroblk] at (2.48,-4.34) {};
\node[infoblk] at (3.10,-4.34) {\tiny 5};
\node[zeroblk] at (3.72,-4.34) {};
\node[zeroblk] at (4.34,-4.34) {};
\node[zeroblk] at (4.96,-4.34) {};
\node[zeroblk] at (5.58,-4.34) {};
\node[zeroblk] at (6.20,-4.34) {};
\node[zeroblk] at (6.82,-4.34) {};
\node[zeroblk] at (7.44,-4.34) {};
\node[zeroblk] at (8.06,-4.34) {};
\node[parblk] at (8.68,-4.34) {\tiny 0};
\node[parblk] at (9.30,-4.34) {\tiny 0};

  \draw[decorate,decoration={brace,amplitude=5pt}] (-0.3,0.45) -- node[above=5pt,lbl]{information part $\mat H_i$ ($k_b=8$ block columns)} (4.64,0.45);
  \draw[decorate,decoration={brace,amplitude=5pt}] (4.66,0.45) -- node[above=5pt,lbl]{parity part $\mat H_p$ (dual diagonal)} (9.6,0.45);
\end{tikzpicture}
\caption{Base matrix of the default code ($k_b=m_b=8$, $z=48$, rate 1/2, $n=768$), read back from
\code{QCLDPCCode.H}. Each number is the cyclic shift of a $48\times48$ identity; empty cells are
zero blocks. Every information column has weight 3.}
\label{fig:ldpc-base}
\end{figure}

\subsection{Encoding by accumulation}
Split $\mat H=[\mat H_i\,|\,\mat H_p]$ and the codeword $\vect c=[\vect u\,|\,\vect p]$. The
parity part is block dual-diagonal: identities on the diagonal and on the sub-diagonal. The
parity equations $\mat H_i\vect u\T\oplus\mat H_p\vect p\T=\vect0$ then read, block by block,
\begin{equation}
  \vect p_0=\vect s_0,\qquad \vect p_j=\vect p_{j-1}\oplus\vect s_j,\qquad \vect s=\mat H_i\vect u\T ,
\end{equation}
so the parity is a cumulative XOR of the syndrome blocks. Encoding costs one sparse product and
needs no Gaussian elimination. The dual-diagonal structure also guarantees that $\mat H$ has
full rank.
\codeexcerpt{ldpc_encode}{Systematic QC-LDPC encoding.}

\subsection{Belief-propagation decoding}
With channel LLRs $L_v$, the flooding schedule alternates two updates.
\begin{align}
  \text{variable to check:}\quad & L_{v\to c}=L_v+\sum_{c'\in\mathcal N(v)\setminus c}L_{c'\to v},\\
  \text{check to variable:}\quad & L_{c\to v}=2\tanh^{-1}\!\prod_{v'\in\mathcal N(c)\setminus v}\tanh\frac{L_{v'\to c}}{2}
  \quad\text{(sum-product)},\\
  & L_{c\to v}\approx\alpha\prod_{v'}\sgn(L_{v'\to c})\cdot\min_{v'\in\mathcal N(c)\setminus v}|L_{v'\to c}|
  \quad\text{(normalized min-sum)}.
\end{align}
The min-sum rule only needs, for each check, the two smallest input magnitudes: the output to
the edge holding the minimum is the second minimum, and to all others the minimum. The factor
$\alpha\approx0.8$ compensates its overestimation of reliabilities. After every iteration the
hard decisions are tested against the syndrome. Frames that satisfy it stop early, and the
average number of iterations is reported as a metric.
\codeexcerpt{ldpc_minsum}{The normalized min-sum check-node update, vectorized over all checks of all frames.}
The message passing is vectorized with a padded table of the edges of every check node
(irregular row weights are padded with neutral entries) and sparse incidence matrices for the
variable-node sums.

% ------------------------------------------------------------------
\section{Polar codes}\label{sec:polar}
\subsection{Channel polarization}
Arıkan's transform~\cite{arikan2009} $\vect x=\vect u\mat G_N$, with
$\mat G_N=\begin{bsmallmatrix}1&0\\1&1\end{bsmallmatrix}^{\otimes n}$, $N=2^n$, combines $N$
uses of a channel $W$ into $N$ synthetic channels $W_N^{(i)}$ seen by the bits $u_i$ under
successive decoding. As $N$ grows they polarize: a fraction tends to noiseless channels, and the
rest to useless ones. One step splits a channel with Bhattacharyya parameter $Z$ into a worse
and a better one,
\begin{equation}
  Z(W^-)\le2Z-Z^2,\qquad Z(W^+)=Z^2 .
\end{equation}
Starting from $Z=e^{-E_s/N_0}$ (BPSK over AWGN at the \prm{design_snr_db}) and applying the
recursion $n$ times ranks the synthetic channels. The $K$ best carry information and the others
are \emph{frozen} to 0.
\codeexcerpt{polar_rel}{Reliability ranking by the Bhattacharyya recursion.}

\begin{figure}[htbp]
\centering
\begin{tikzpicture}
\begin{semilogyaxis}[width=11cm,height=5cm,xlabel={synthetic channel index $i$},ylabel={$Z(W_N^{(i)})$},
  xmin=-1,xmax=64,ymin=1e-12,ymax=2,legend pos=south west]
  \addplot[sim={accent2,mark=x,mark size=2pt}] table[x=i,y=z]{\datafile{polar_frozen}};
  \addplot[sim=accent3] table[x=i,y=z]{\datafile{polar_info}};
  \legend{frozen,information}
\end{semilogyaxis}
\end{tikzpicture}
\caption{Polarization for $N=64$, $K=32$ at a design SNR of 0\,dB (\code{reliability},
\code{frozen_mask}). Reliable channels concentrate at high indices, the frozen set at low
indices.}
\label{fig:polar}
\end{figure}

\subsection{Encoding and successive-cancellation decoding}
Without bit reversal, $\vect u=[\vect u_a,\vect u_b]$ encodes recursively as
$\vect x=[\mathrm{enc}(\vect u_a)\oplus\mathrm{enc}(\vect u_b),\ \mathrm{enc}(\vect u_b)]$,
computed in place with $\log_2N$ butterfly stages:
\codeexcerpt{polar_encode}{The polar transform in $\log_2 N$ vectorized stages.}
The SC decoder follows the same recursion on LLRs. For the halves $\vect L_1,\vect L_2$,
\begin{equation}
  f(a,b)=\sgn a\,\sgn b\,\min(|a|,|b|)\ \ (\text{left child}),\qquad
  g(a,b,\hat x)=b+(1-2\hat x)\,a\ \ (\text{right child}).
\end{equation}
The left half is decoded first; its re-encoded decisions $\hat{\vect x}_a$ enter the right half.
Subtrees made only of frozen bits return zeros immediately, and subtrees made only of
information bits (rate-1 nodes) are decided by a direct hard decision. Both shortcuts are exact.
\codeexcerpt{polar_sc}{Recursive SC decoding with rate-0 and rate-1 node shortcuts.}

% ------------------------------------------------------------------
\section{Comparison}
\begin{figure}[htbp]
\centering
\begin{tikzpicture}
\begin{semilogyaxis}[berplot,width=9.6cm,height=7.5cm,xmin=0,xmax=10,ymin=1e-6,ymax=0.5,
  legend style={at={(1.03,1)},anchor=north west}]
  \addplot[th=black!60] table[x=ebn0,y=uncoded]{\datafile{coding_theory}};
  \addplot[th=accent4,dashed] table[x=ebn0,y=hamming]{\datafile{coding_theory}};
  \addplot[sim=accent4] table[x=ebn0,y=hamming]{\datafile{coding_sim}};
  \addplot[simline={accent}] table[x=ebn0,y=conv]{\datafile{coding_sim}};
  \addplot[simline={accent3,mark=square*}] table[x=ebn0,y=polar]{\datafile{coding_sim}};
  \addplot[simline={accent2,mark=triangle*}] table[x=ebn0,y=ldpc]{\datafile{coding_sim}};
  \legend{uncoded (theory),{Hamming(7,4) hard (approx.)},{Hamming(7,4) simulated},
    conv. $K{=}7$ soft Viterbi,polar $N{=}256$ SC,QC-LDPC $n{=}768$ min-sum}
\end{semilogyaxis}
\end{tikzpicture}
\caption{QPSK over AWGN with the codes of this chapter (rate $\approx1/2$ except Hamming 4/7;
max-log LLRs; simulated with the real blocks, up to 200 errors or 20 batches per point). At
$10^{-4}$ the Hamming code gains well under 1\,dB with hard decisions, while the convolutional and
LDPC codes gain about 5 and 5.5\,dB.}
\label{fig:coding}
\end{figure}

\begin{keypoint}
Soft decisions are worth about 2\,dB, and iterative decoding of long sparse codes approaches
capacity. The rate-1/2 limit for binary-input AWGN is 0.19\,dB. The 768-bit LDPC code is about
3\,dB away at $10^{-5}$; longer codes get closer.
\end{keypoint}

\chapter{Wireline communications}\label{ch:wireline}

Copper links (Ethernet over twisted pair, SerDes over backplanes, DSL over telephone loops)
are real-valued baseband systems. Their channel is a lossy, dispersive transmission line, and
their noise is dominated by crosstalk, echo and jitter rather than thermal noise.
\file{phy/wireline.py} models these effects; all its blocks keep real signals real.

\begin{figure}[htbp]
\centering
\begin{tikzpicture}[node distance=3mm and 4mm,every node/.style={font=\sffamily\scriptsize}]
  \node[lbl,anchor=east] at (-0.2,0) {SerDes};
  \node[blockbox,minimum width=11mm] (a1) at (0.6,0) {PAM-4};
  \node[blockbox,right=of a1,minimum width=11mm] (a2) {pulse};
  \node[warmbox,right=of a2,minimum width=13mm] (a3) {backplane};
  \node[warmbox,right=of a3,minimum width=11mm] (a4) {jitter};
  \node[greenbox,right=of a4,minimum width=11mm] (a5) {CTLE};
  \node[greenbox,right=of a5,minimum width=11mm] (a6) {FFE+DFE};
  \foreach \x/\y in {a1/a2,a2/a3,a3/a4,a4/a5,a5/a6} \draw[sig] (\x)--(\y);
  \node[lbl,anchor=east] at (-0.2,-1.2) {Ethernet};
  \node[blockbox,minimum width=11mm] (b1) at (0.6,-1.2) {4B/5B};
  \node[blockbox,right=of b1,minimum width=11mm] (b2) {MLT-3};
  \node[warmbox,right=of b2,minimum width=13mm] (b3) {Cat5e};
  \node[warmbox,right=of b3,minimum width=11mm] (b4) {echo};
  \node[greenbox,right=of b4,minimum width=11mm] (b5) {echo canc.};
  \node[greenbox,right=of b5,minimum width=11mm] (b6) {DFE};
  \foreach \x/\y in {b1/b2,b2/b3,b3/b4,b4/b5,b5/b6} \draw[sig] (\x)--(\y);
  \node[lbl,anchor=east] at (-0.2,-2.4) {DSL};
  \node[blockbox,minimum width=11mm] (c1) at (0.6,-2.4) {DMT};
  \node[warmbox,right=of c1,minimum width=11mm] (c2) {loop};
  \node[warmbox,right=of c2,minimum width=13mm] (c3) {NEXT/FEXT};
  \node[warmbox,right=of c3,minimum width=11mm] (c4) {impulse};
  \node[greenbox,right=of c4,minimum width=11mm] (c5) {FFT + FEQ};
  \node[greenbox,right=of c5,minimum width=11mm] (c6) {demap};
  \foreach \x/\y in {c1/c2,c2/c3,c3/c4,c4/c5,c5/c6} \draw[sig] (\x)--(\y);
\end{tikzpicture}
\caption{Simplified chains of the wireline scenarios (10GBASE-KR / PAM-4 SerDes, 100BASE-TX,
ADSL/VDSL).}
\end{figure}

% ------------------------------------------------------------------
\section{Line codes}\label{sec:line-codes}
A line code shapes the spectrum of the transmitted symbols: it removes DC (transformer
coupling), guarantees transitions (clock recovery) or concentrates energy at low frequency
(cable loss, emissions). \blk{line_encoder} implements four codes:
\begin{center}\small
\begin{tabular}{@{}llll@{}}
\toprule
code & rule & levels & PSD with rectangular pulses ($T$ = bit time)\\
\midrule
NRZ & $1\to+1$, $0\to-1$ & 2 & $T\sinc^2(fT)$\\
Manchester & $0\to(+1,-1)$, $1\to(-1,+1)$ & 2 & $T\sinc^2(fT/2)\sin^2(\pi fT/2)$\\
AMI & $0\to0$, $1\to\pm1$ alternately & 3 & $T\sinc^2(fT)\sin^2(\pi fT)$\\
MLT-3 & each 1 steps through $0,+1,0,-1$ & 3 & concentrated below $1/(4T)$\\
\bottomrule
\end{tabular}
\end{center}
Manchester and AMI have a spectral null at DC. MLT-3, used by 100BASE-TX, completes a full
cycle every four ones, so its fundamental is at most a quarter of the bit rate.

\begin{figure}[htbp]
\centering
\begin{tikzpicture}
\begin{axis}[width=11cm,height=5.5cm,xlabel={frequency / bit rate},ylabel={PSD [dB]},xmin=0,xmax=3,
  ymin=-40,ymax=3,legend pos=north east]
  \addplot[th=accent] table[x=f,y=nrz]{\datafile{linecode_psd}};
  \addplot[th=accent2] table[x=f,y=manchester]{\datafile{linecode_psd}};
  \addplot[th=accent3] table[x=f,y=ami]{\datafile{linecode_psd}};
  \addplot[th=accent4] table[x=f,y=mlt3]{\datafile{linecode_psd}};
  \legend{NRZ,Manchester,AMI,MLT-3}
\end{axis}
\end{tikzpicture}
\caption{Welch spectra of random data through \blk{line_encoder} and a rectangular
\blk{pulse_shaper}, normalized to the NRZ peak and clipped at $-40$\,dB. Note the DC nulls of
Manchester and AMI, and Manchester's doubled bandwidth.}
\label{fig:linecode-psd}
\end{figure}

\paragraph{Block line codes.}
\blk{line_block_encoder} implements two block codes.
\begin{itemize}
  \item \textbf{4B/5B} maps each nibble to a 5-bit word with at most one leading and two trailing
  zeros. Long zero runs are thus excluded, at rate $0.8$.
  \item \textbf{64B/66B} prepends the sync header \code{01} to each 64-bit block and scrambles
  the payload with the self-synchronizing polynomial $x^{58}+x^{39}+1$,
  $s[n]=d[n]\oplus s[n-39]\oplus s[n-58]$. The descrambler needs no synchronization, but each
  channel error produces three errors after descrambling (error multiplication, tested in
  \file{tests/test_wireline_optical.py}).
\end{itemize}

% ------------------------------------------------------------------
\section{The cable as a transmission line}\label{sec:cable}
A uniform line of length $\ell$, terminated in its characteristic impedance, has transfer
function $H(f)=e^{-\gamma(f)\ell}$ with propagation constant
$\gamma=\sqrt{(R+j\omega L)(G+j\omega C)}=\alpha+j\beta$. Two loss mechanisms dominate:
\begin{itemize}
  \item the \emph{skin effect}: current crowds near the conductor surface, so $R\propto\sqrt f$
  and the attenuation $\alpha_s\propto\sqrt f$;
  \item \emph{dielectric loss}: $G\propto f$ gives $\alpha_d\propto f$.
\end{itemize}
For the fitted presets, in dB per km,
\begin{equation}
  \alpha_{\mathrm{dB}}(f)=a\sqrt{f_{\mathrm{MHz}}}+b\,f_{\mathrm{MHz}}\;\Big(+\;c/\sqrt{f_{\mathrm{MHz}}}\Big) .
\end{equation}
The last term belongs to the TIA-568 category formulas (Cat\,5e: $c=0.050$ per 100\,m, Cat\,6A: 0.25)
and is held constant below 1\,MHz, where the formula is not specified. It is applied as an
attenuation only; it is a fraction of a dB in the Ethernet band.

\begin{keypoint}[Causality fixes the phase]
A causal, minimum-phase channel cannot have an arbitrary attenuation: the log-magnitude and the
phase form a Hilbert-transform pair. For the skin effect the pair is exact and elementary. The
function $e^{-k\sqrt{j\omega}}$ is analytic in the right half $s$-plane (it is the
Laplace-domain response of a diffusion line), and
\[
  e^{-k\sqrt{j\omega}}=e^{-k\sqrt{\omega}\,(1+j)/\sqrt2}:
\]
the $\sqrt f$ loss comes with an \emph{equal} $\sqrt f$ phase lag. LinkStudio therefore uses
\begin{equation}
  H(f)=\exp\!\Big(-\alpha_s(f)\,\ell\,\big(1+j\,\sgn f\big)-\alpha_d(f)\,\ell\Big),
\end{equation}
with $\alpha$ in nepers. Dropping the phase would produce a symmetric, non-causal pulse;
keeping it produces the long causal tail characteristic of copper (\cref{fig:cable}, right),
which is what makes DSL and Ethernet equalizers work hard.

The dielectric term needs the same care. A loss exactly proportional to $|f|$ at all frequencies
violates the Paley--Wiener condition, so it has no exact causal partner; its Hilbert transform
diverges logarithmically. Physical dielectrics have a nearly constant loss tangent over a wide but
finite band. The Djordjevic--Sarkar permittivity~\cite{djordjevic2001} captures this:
\[
  \varepsilon(\omega)=\varepsilon_\infty+\frac{\Delta\varepsilon}{m_2-m_1}\log_{10}\frac{\omega_2+j\omega}{\omega_1+j\omega},
  \qquad \omega_{1,2}=10^{m_{1,2}}\ \mathrm{rad/s}.
\]
It is a continuum of Debye relaxations, hence causal, and its loss tangent is nearly flat between
$\omega_1$ and $\omega_2$. LinkStudio uses $m_1=4$, $m_2=12$, sets $\tan\delta$ so that the
attenuation is $b\,f$ (within 1\,\% from 10\,MHz to 1\,GHz) and $\varepsilon'(1\,\mathrm{GHz})=\varepsilon_r$
(4.0 for FR-4, 2.3 for polyethylene, 1.5 for foamed coax). It then applies
$e^{-j\omega\ell(\sqrt{\varepsilon(\omega)}-\sqrt{\varepsilon_\infty})/c}$; subtracting
$\sqrt{\varepsilon_\infty}$ removes only the front delay. On the 1\,m backplane the energy of the
impulse response before $t=0$ drops from 40\,\% (zero-phase dielectric term) to $5\times10^{-7}$.
\end{keypoint}
\codeexcerpt{cable}{Cable response: skin effect with its phase, TIA term and causal dielectric (\file{wireline.py}).}
\codeexcerpt{dielectric}{Djordjevic--Sarkar dielectric response.}

\subsection{Telephone loops from primary constants}
DSL standards specify loops through their primary constants, not through fitted losses.
ITU-T G.996.1 Annex~A~\cite{g9961} tabulates $R$, $L$, $G$ and $C$ of polyethylene-insulated
26\,AWG (0.4\,mm) and 24\,AWG (0.5\,mm) cable from 1\,Hz to 5\,MHz. The \code{pic_26awg} and
\code{pic_24awg} presets interpolate these tables at 70\,\textdegree F in $\log f$. Above 5\,MHz
they extend them with $R\propto\sqrt f$, $G\propto f$ and constant $L$, $C$. The loop is a
two-port with ABCD matrix
\begin{equation}
  \begin{pmatrix}A&B\\C&D\end{pmatrix}=\begin{pmatrix}\cosh\gamma\ell & Z_0\sinh\gamma\ell\\ \sinh(\gamma\ell)/Z_0 & \cosh\gamma\ell\end{pmatrix},\qquad
  Z_0=\sqrt{\frac{R+j\omega L}{G+j\omega C}},
\end{equation}
between source and load resistances $Z_S=Z_L=Z$ (\prm{termination_ohm}, 100\,$\Omega$ as in
G.996.1). The insertion transfer function relative to a direct connection is
\begin{equation}
  H(f)=\frac{2Z}{AZ+B+CZ^2+DZ} .
\end{equation}
Unlike the $\sqrt f$ fit, this model has the right low-frequency behaviour: below about 20\,kHz the
loss is set by the DC resistance $R_0\ell$ and the impedance mismatch, not by $\sqrt f$. G.996.1
also publishes the insertion loss of its test loops. \Cref{fig:pic} compares the model with
Table~2 for the Mid-C loop (1829\,m of 26\,AWG): the deviation is below 0.3\,dB from 20\,kHz to
1.1\,MHz, the rounding of the table. The \emph{ADSL-like DMT} scenario now runs over this loop.
\codeexcerpt{pic}{ABCD response of a PIC loop.}

\begin{figure}[htbp]
\centering
\begin{tikzpicture}
\begin{semilogxaxis}[width=10cm,height=5.6cm,xlabel={frequency [kHz]},ylabel={insertion loss [dB]},
  xmin=10,xmax=2000,ymin=0,ymax=70,legend pos=north west]
  \addplot[th=accent] table[x=f_khz,y=model]{\datafile{pic_loss}};
  \addplot[only marks,mark=square*,mark size=1.6pt,accent2] table[x=f_khz,y=table]{\datafile{pic_g9961}};
  \addplot[th=black!50,dashed] table[x=f_khz,y=fit]{\datafile{pic_loss}};
  \legend{\code{pic_26awg} (RLCG + ABCD),{G.996.1 Table 2, Mid-C},\code{dsl_0.4mm} ($\sqrt f$ fit)}
\end{semilogxaxis}
\end{tikzpicture}
\caption{Insertion loss of 1829\,m of 26\,AWG cable between 100\,$\Omega$ terminations: RLCG model
from the G.996.1 primary constants against the values published in G.996.1 Table~2, and the
earlier $\sqrt f$ fit for comparison.}
\label{fig:pic}
\end{figure}

\begin{figure}[htbp]
\centering
\begin{tikzpicture}
\begin{groupplot}[group style={group size=2 by 1,horizontal sep=17mm},width=6.9cm,height=5.3cm]
\nextgroupplot[xmode=log,xlabel={frequency [MHz]},ylabel={insertion loss [dB]},xmin=0.01,xmax=30,
  ymin=0,ymax=120,legend pos=north west,title={loss}]
  \addplot[th=accent] table[x=f_mhz,y=dsl1km]{\datafile{cable_loss}};
  \addplot[th=accent3] table[x=f_mhz,y=dsl2km]{\datafile{cable_loss}};
  \addplot[th=accent2] table[x=f_mhz,y=dsl3km]{\datafile{cable_loss}};
  \addplot[th=accent4,dashed] table[x=f_mhz,y=cat5e100m]{\datafile{cable_loss}};
  \legend{{0.4\,mm, 1\,km},{0.4\,mm, 2\,km},{0.4\,mm, 3\,km},{Cat\,5e, 100\,m}}
\nextgroupplot[xlabel={time [\si{\micro s}]},ylabel={normalized $h(t)$},xmin=0,xmax=40,ymin=-0.25,ymax=1.05,
  legend pos=north east,title={pulse response, 0.4\,mm, 2\,km}]
  \addplot[th=accent] table[x=t_us,y=h]{\datafile{cable_impulse_min}};
  \addplot[th=black!50,dashed] table[x=t_us,y=h]{\datafile{cable_impulse_zero}};
  \legend{with skin phase (causal),magnitude only (shifted)}
\end{groupplot}
\end{tikzpicture}
\caption{Left: insertion loss of the presets (\code{cable_response}). Right: impulse response at
17.664\,MSa/s with the minimum-phase term (used by \blk{copper_line}) and, for comparison, the
zero-phase response with the same magnitude, shifted by 20\,\si{\micro s}.}
\label{fig:cable}
\end{figure}

\begin{incode}
\blk{copper_line} applies $H(f)$ in the frequency domain, with zero padding to avoid circular
wrap-around. It publishes the sampled impulse response, truncated at 99.95\,\% of its energy,
on \code{csi} (\meta{csi_type}\code{="cir"}), for \blk{linear_equalizer} or the DMT FEQ.
Presets: Cat\,5e, Cat\,6A, 0.4 and 0.5\,mm telephone pairs (fit and G.996.1 RLCG), RG-6 coax
and an FR-4 backplane trace. The RLCG presets remove only the front delay $\ell\sqrt{LC}$; the
remaining dispersive delay is part of the channel.
\end{incode}

% ------------------------------------------------------------------
\section{Crosstalk}
Signals on neighbouring pairs of a binder couple electromagnetically. The ANSI T1.417~\cite{ansi417} 99\,\%
worst-case models for $N$ disturbers give the coupling power transfer functions
\begin{align}
  |H_{\mathrm{NEXT}}(f)|^2&=8.536\times10^{-15}\,N^{0.6}f^{1.5},\\
  |H_{\mathrm{FEXT}}(f)|^2&=7.744\times10^{-21}\,N^{0.6}\,\ell_{\mathrm{ft}}\,f^2\,|H_{\mathrm{line}}(f)|^2 .
\end{align}
NEXT (near end) does not depend on the length and grows with frequency. FEXT (far end) travels
the whole line and is attenuated by it. \blk{crosstalk} builds the disturbance from the
spectrum of the neighbour's transmit signal (input \code{tx}, or the signal itself) with random
phases, shapes it with the model and adds it. The crosstalk-to-signal ratio is reported.
\codeexcerpt{xtalk}{ANSI crosstalk models.}

% ------------------------------------------------------------------
\section{Full duplex: hybrid echo and echo cancellation}
Full-duplex transceivers (1000BASE-T, DSL) send and receive on the same pair. The hybrid
leaks part of the local transmit signal into the receiver as \emph{echo}, often much stronger
than the far-end signal. \blk{hybrid_echo} adds the transmit signal filtered by a short echo
path; \blk{echo_canceller} identifies that path with an NLMS filter driven by the known local
transmit signal and subtracts its estimate.

NLMS converges at a speed proportional to $\mu$, but leaves an excess mean-square error, the
misadjustment $\mathcal M\approx\mu/(2-\mu)$ for NLMS. That excess error is noise added to the
far-end signal. The block therefore uses a \emph{gear shift}: a large step for the first
\prm{train_samples} samples (fast acquisition, as at link start-up), then a small one (low
residual). The power reduction after convergence is reported as a metric.
\codeexcerpt{echo_nlms}{NLMS echo cancellation with gear shifting.}

% ------------------------------------------------------------------
\section{SerDes equalization}\label{sec:dfe}
\subsection{Continuous-time linear equalizer}
A CTLE is an analog high-pass with one zero and two poles,
\begin{equation}
  H_{\mathrm{CTLE}}(s)=A\,\frac{1+s/\omega_z}{(1+s/\omega_{p1})(1+s/\omega_{p2})},
\end{equation}
providing a peaking of about $20\log_{10}(\omega_{p1}/\omega_z)$\,dB that partially inverts the
channel loss (\cref{fig:ctle}). It amplifies high-frequency noise and crosstalk as well, so it
is followed by a decision-feedback equalizer.

\begin{figure}[htbp]
\centering
\begin{tikzpicture}
\begin{axis}[width=10cm,height=5.2cm,xlabel={frequency [GHz]},ylabel={magnitude [dB]},xmin=0,xmax=20,
  ymin=-25,ymax=15,legend pos=south west]
  \addplot[th=accent2] table[x=f_ghz,y=channel_db]{\datafile{ctle}};
  \addplot[th=accent3] table[x=f_ghz,y=ctle_db]{\datafile{ctle}};
  \addplot[th=accent] table[x=f_ghz,y=total_db]{\datafile{ctle}};
  \legend{backplane 0.8\,m,CTLE ($f_z{=}3$, $f_{p1}{=}20$, $f_{p2}{=}30$\,GHz),combined}
  \draw[black!50,dashed] (axis cs:14,-25) -- (axis cs:14,15) node[pos=0.92,left,font=\scriptsize]{Nyquist, 28\,GBd};
\end{axis}
\end{tikzpicture}
\caption{The CTLE keeps the combined response within about 2\,dB up to the Nyquist frequency of
a 28\,GBd link (\code{cable_response} and \code{CTLE.response}). Beyond it the second pole and
the growing loss take over.}
\label{fig:ctle}
\end{figure}

\subsection{Feed-forward plus decision-feedback equalizer}
The DFE output is
\begin{equation}
  z[n]=\sum_{i=0}^{F-1}w_i\,y[n+D-i]-\sum_{j=1}^{K}b_j\,\hat a[n-j],
\end{equation}
where the feedback taps cancel the post-cursor ISI of past decisions $\hat a$. Because past
decisions are noise-free when correct, the DFE removes ISI \emph{without} noise enhancement,
unlike a linear equalizer. The price is error propagation: a wrong decision feeds wrong ISI back.
Both filters adapt jointly with NLMS on the stacked regressor $[\vect u;\hat{\vect a}]$, trained on
\prm{n_train} known symbols and then decision-directed.
\codeexcerpt{dfe_loop}{Joint NLMS adaptation of the FFE and the DFE (\blk{dfe_equalizer}).}

\subsection{Jitter}
Sampling-clock jitter perturbs the sampling instants, $t_n=nT_s+\Delta(n)$, with a random
(Gaussian, RJ) and a sinusoidal (SJ) component expressed in unit intervals.
\blk{jitter} resamples the waveform at the perturbed instants with the cubic interpolator; RJ is
drawn once per symbol and interpolated between symbols. To first order, the error is
$x'(t)\,\Delta$, which is why steep edges, that is high bandwidth, make jitter costly.
\codeexcerpt{jitter}{Perturbed sampling instants.}

% ------------------------------------------------------------------
\section{Discrete multitone (DSL)}\label{sec:dmt}
\subsection{Real OFDM}
DMT is OFDM for real baseband channels. Imposing Hermitian symmetry, $X_{N-k}=X_k\conj$, makes
the IFFT output real. Tones $k=1,\dots,N/2-1$ carry independent QAM symbols; DC and Nyquist are
unused. The cyclic prefix turns the loop into $N/2-1$ parallel flat channels, as in
\cref{ch:ofdm}. The demodulator equalizes each tone with one division (FEQ) using the channel's
impulse response.
\codeexcerpt{dmt_hermitian}{Per-tone mapping and Hermitian symmetry (\blk{dmt_modulator}).}

\subsection{Bit loading and the SNR gap}
For uncoded QAM at a target error rate, the achievable bits per tone are well approximated by
the \emph{gap} formula
\begin{equation}
  b_k=\log_2\!\Big(1+\frac{\mathrm{SNR}_k}{\Gamma\,\gamma_m}\Big),
  \qquad\Gamma\approx9.8\,\mathrm{dB}\ (P_s=10^{-7}),
\end{equation}
where $\gamma_m$ is a noise margin~\cite{starr1999}. The block computes the expected
$\mathrm{SNR}_k$ from the cable model and the AWGN level. With the time-domain power
$P=2N_{\mathrm{tones}}/N$ and $\sigma^2=P/\mathrm{SNR}$,
$\mathrm{SNR}_k=|H_k|^2\,\mathrm{SNR}\,N/(2N_{\mathrm{tones}})$. It then rounds $b_k$ down to the
sizes the mapper supports: 0, 1 (BPSK) or an even number (square QAM).
\codeexcerpt{waterfill}{Gap-approximation bit loading.}

\begin{pitfall}[Bit loading versus water-filling]
The transmit PSD is flat; only the bits adapt to the channel. True water-filling would also
redistribute power, pouring more where the SNR is high. With a flat PSD the loss is small at
high SNR, but it is a modelling choice to keep in mind.
\end{pitfall}

\begin{figure}[htbp]
\centering
\begin{tikzpicture}
\begin{axis}[width=11cm,height=5.4cm,xlabel={tone index $k$ (4.3125\,kHz spacing)},ylabel={bits per tone},
  xmin=30,xmax=256,ymin=0,ymax=15,legend pos=north east,const plot]
  \addplot[accent,thick] table[x=tone,y=bits1km]{\datafile{dmt_bits}};
  \addplot[accent3,thick] table[x=tone,y=bits2km]{\datafile{dmt_bits}};
  \addplot[accent2,thick] table[x=tone,y=bits3km]{\datafile{dmt_bits}};
  \legend{1\,km (\dmtrateA\,Mb/s),2\,km (\dmtrateB\,Mb/s),3\,km (\dmtrateC\,Mb/s)}
\end{axis}
\end{tikzpicture}
\caption{Bit loading of \blk{dmt_modulator} on a 0.4\,mm loop (ADSL downstream tones 33--255,
50\,dB SNR, 9.8\,dB gap, 3\,dB margin). Data rates are those reported by the block.}
\label{fig:dmt}
\end{figure}

\chapter{Optical communications}\label{ch:optical}

\file{phy/optical.py} models fiber-optic and free-space optical links, from the laser to the
coherent receiver's digital signal processing. \Cref{fig:optical-chains} shows the two families
of systems it supports.

\begin{figure}[htbp]
\centering
\begin{tikzpicture}[node distance=3mm and 3.6mm,every node/.style={font=\sffamily\scriptsize}]
  \node[lbl,anchor=east] at (-0.15,0) {IM/DD};
  \node[blockbox,minimum width=10mm] (a1) at (0.55,0) {PAM/NRZ};
  \node[blockbox,right=of a1,minimum width=10mm] (a2) {laser /\\EML};
  \node[warmbox,right=of a2,minimum width=10mm] (a3) {fiber};
  \node[warmbox,right=of a3,minimum width=10mm] (a4) {EDFA};
  \node[greenbox,right=of a4,minimum width=10mm] (a5) {PIN /\\APD};
  \node[greenbox,right=of a5,minimum width=10mm] (a6) {AGC};
  \node[greenbox,right=of a6,minimum width=10mm] (a7) {decision};
  \foreach \x/\y in {a1/a2,a2/a3,a3/a4,a4/a5,a5/a6,a6/a7} \draw[sig] (\x)--(\y);
  \node[lbl,anchor=east] at (-0.15,-1.35) {coherent};
  \node[blockbox,minimum width=10mm] (b1) at (0.55,-1.35) {DP-QAM};
  \node[blockbox,right=of b1,minimum width=10mm] (b2) {IQ-MZM};
  \node[warmbox,right=of b2,minimum width=12mm] (b3) {$N\times$(span\\+EDFA)};
  \node[greenbox,right=of b3,minimum width=10mm] (b4) {coherent\\front end};
  \node[greenbox,right=of b4,minimum width=8mm] (b5) {CD};
  \node[greenbox,right=of b5,minimum width=8mm] (b6) {FO};
  \node[greenbox,right=of b6,minimum width=10mm] (b7) {$2{\times}2$\\MIMO};
  \node[greenbox,right=of b7,minimum width=8mm] (b8) {CPR};
  \foreach \x/\y in {b1/b2,b2/b3,b3/b4,b4/b5,b5/b6,b6/b7,b7/b8} \draw[sig] (\x)--(\y);
\end{tikzpicture}
\caption{Intensity modulation with direct detection, and dual-polarization coherent transmission.}
\label{fig:optical-chains}
\end{figure}

\begin{incode}
Optical signals are complex field envelopes around the optical carrier, scaled so that
$|E|^2$ is the instantaneous optical power in watts. Dual polarization uses the shape
$(B,2,n)$. Optical blocks maintain the following metadata:
\begin{itemize}
  \item \meta{wavelength_nm};
  \item \meta{power}, the average launch or received power in W;
  \item \meta{cd_ps_nm}, the accumulated chromatic dispersion $D\ell$;
  \item \meta{ase_psd}, the accumulated ASE density per polarization in W/Hz.
\end{itemize}
The OSNR is computed from the last two at any point of the link.
\end{incode}

% ==================================================================
\section{Transmitters}
\subsection{Directly modulated lasers}
An intensity modulator maps the electrical drive $x\in[-1,1]$ to the optical power
\begin{equation}
  P(t)=P_{\mathrm{avg}}\big(1+m\,x(t)\big),\qquad m=\frac{\mathrm{ER}-1}{\mathrm{ER}+1},
\end{equation}
since the extinction ratio is $\mathrm{ER}=P_1/P_0=(1+m)/(1-m)$. Three non-idealities are
modelled:
\begin{itemize}
  \item \textbf{Relative intensity noise}: multiplicative Gaussian noise of variance
  $\mathrm{RIN}\cdot f_s/2$ on $P$.
  \item \textbf{Phase noise}: a Wiener process with the laser linewidth (\cref{eq:wiener}).
  \item \textbf{Transient chirp} of a directly modulated laser. The instantaneous frequency
  follows the power changes, $\Delta\nu(t)=\frac{\alpha}{4\pi}\frac{d}{dt}\ln P(t)$, with $\alpha$
  the Henry linewidth-enhancement factor. Integrating gives the phase
  \begin{equation}
    \phi_{\mathrm{chirp}}(t)=\frac{\alpha}{2}\ln\frac{P(t)}{P_{\mathrm{avg}}} .
  \end{equation}
\end{itemize}
The field is $E=\sqrt P\,e^{j\phi}$.
\codeexcerpt{im_mod}{Laser intensity modulation with RIN, phase noise and chirp (\blk{intensity_modulator}).}

\begin{pitfall}[Not a rate-equation model]
The laser is modelled behaviourally (power, chirp, linewidth). Relaxation oscillations, the
adiabatic chirp term and thermal effects, which require the laser rate equations, are not
included.
\end{pitfall}

\subsection{Mach--Zehnder and IQ modulators}
A Mach--Zehnder modulator in push--pull configuration, biased at null, has the field transfer
function $E_{\mathrm{out}}/E_{\mathrm{in}}=\sin\big(\frac{\pi}{2}V/V_\pi\big)$. An IQ modulator
nests two of them, with a $90^\circ$ phase shift between the arms:
\begin{equation}
  \frac{E_{\mathrm{out}}}{E_{\mathrm{in}}}\propto\sin\Big(\frac\pi2\frac{V_I}{V_\pi}\Big)+j\sin\Big(\frac\pi2\frac{V_Q}{V_\pi}\Big).
\end{equation}
With the peak drive set to \prm{drive_swing}$\times V_\pi$, small swings are nearly linear,
while large swings compress the outer constellation points (\cref{fig:mzm}). The \prm{linear}
model maps the drive directly onto the field. The laser power is split equally between the
polarizations.

\begin{figure}[htbp]
\centering
\begin{tikzpicture}
\begin{axis}[width=9cm,height=5cm,xlabel={drive $V/V_\pi$},ylabel={normalized output},xmin=-1,xmax=1,
  ymin=-1.1,ymax=1.1,legend pos=north west,samples=101]
  \addplot[th=accent,domain=-1:1] {sin(90*x)};
  \addplot[th=accent2,domain=-1:1] {sin(90*x)^2};
  \addplot[th=black!50,dashed,domain=-0.6:0.6] {pi/2*x};
  \draw[accent3,fill=accent3!15,opacity=0.5] (axis cs:-0.5,-1.1) rectangle (axis cs:0.5,1.1);
  \legend{field $\sin(\frac\pi2 V/V_\pi)$,power,linear approx.}
\end{axis}
\end{tikzpicture}
\caption{MZM transfer function at null bias. The shaded region is the default drive swing
($0.5\,V_\pi$).}
\label{fig:mzm}
\end{figure}
\codeexcerpt{iq_mzm}{Sinusoidal MZM transfer and power normalization (\blk{iq_modulator}).}

% ==================================================================
\section{Propagation in single-mode fiber}\label{sec:nlse}
\subsection{The nonlinear Schrödinger equation}
The slowly varying envelope $A(z,t)$ of a single-mode fiber, in a frame moving with the group
velocity, obeys
\begin{equation}
  \frac{\partial A}{\partial z}=
  \underbrace{-\frac{\alpha}{2}A-j\frac{\beta_2}{2}\frac{\partial^2A}{\partial t^2}}_{\hat D A}
  \;+\;\underbrace{j\gamma|A|^2A}_{\hat N A},
  \label{eq:nlse}
\end{equation}
with three ingredients~\cite{agrawal2019}:
\begin{itemize}
  \item the power attenuation $\alpha$, given in dB/km as \prm{attenuation_db_km};
  \item the group-velocity dispersion $\beta_2$, related to the dispersion parameter $D$
  (ps/nm/km) by
  \begin{equation}
    \beta_2=-\frac{D\lambda^2}{2\pi c}\qquad(\beta_2\approx\SI{-21.7}{ps^2/km}\ \text{for }D=17,\ \lambda=1550\,\mathrm{nm});
  \end{equation}
  \item the Kerr coefficient $\gamma$, given in 1/W/km as \prm{gamma_w_km}.
\end{itemize}
With the Fourier convention of \cref{eq:fourier-convention}, $\partial_t\to j\omega$, so the
linear part has the exact solution
\begin{equation}
  \tilde A(z,\omega)=\tilde A(0,\omega)\exp\!\Big(j\frac{\beta_2}{2}\omega^2z-\frac{\alpha}{2}z\Big),
  \label{eq:dispersion-operator}
\end{equation}
an all-pass filter with quadratic phase. The dispersion therefore depends only on the
accumulated $D\ell$ (ps/nm), the quantity tracked in \meta{cd_ps_nm}.
\codeexcerpt{beta2}{From $D$ to $\beta_2$.}
\codeexcerpt{disp_op}{The dispersion operator of \cref{eq:dispersion-operator}.}

\paragraph{Dispersion alone.}
An unchirped Gaussian pulse $A(0,t)=e^{-t^2/(2T_0^2)}$ stays Gaussian and broadens as
\begin{equation}
  \frac{T_1(z)}{T_0}=\sqrt{1+\Big(\frac{z}{L_D}\Big)^2},\qquad L_D=\frac{T_0^2}{|\beta_2|}.
\end{equation}
For $T_0=5$\,ps and $|\beta_2|=\fiberbeta$\,ps$^2$/km, $L_D=\fiberLD$\,km.
\Cref{fig:dispersion} compares the \code{propagate} solver with this result. The simulated RMS
widths agree with the formula to all printed digits, as expected from an exact spectral method.

\begin{figure}[htbp]
\centering
\begin{tikzpicture}
\begin{groupplot}[group style={group size=2 by 1,horizontal sep=16mm},width=6.9cm,height=5.2cm]
\nextgroupplot[xlabel={$t$ [ps]},ylabel={$|A|^2/P_0$},xmin=-40,xmax=40,ymin=0,ymax=1.05,
  legend pos=north east,title={pulse shape}]
  \addplot[th=black!70] table[x=t_ps,y=z0]{\datafile{fiber_pulses}};
  \addplot[th=accent] table[x=t_ps,y=z1]{\datafile{fiber_pulses}};
  \addplot[th=accent3] table[x=t_ps,y=z2]{\datafile{fiber_pulses}};
  \addplot[th=accent2] table[x=t_ps,y=z4]{\datafile{fiber_pulses}};
  \legend{$z=0$,$L_D$,$2L_D$,$4L_D$}
\nextgroupplot[xlabel={$z/L_D$},ylabel={$T_1/T_0$},xmin=0,xmax=5,ymin=0.8,ymax=5.4,legend pos=north west,
  title={RMS broadening}]
  \addplot[th=accent] table[x=z,y=theory]{\datafile{fiber_broadening}};
  \addplot[sim=accent2] table[x=z,y=sim]{\datafile{fiber_broadening}};
  \legend{$\sqrt{1+(z/L_D)^2}$,\code{propagate}}
\end{groupplot}
\end{tikzpicture}
\caption{Chromatic dispersion of a 5-ps Gaussian pulse in standard single-mode fiber
($D=17$\,ps/nm/km, no loss).}
\label{fig:dispersion}
\end{figure}

\paragraph{Kerr nonlinearity alone.}
Without dispersion \cref{eq:nlse} integrates exactly. The power profile is unchanged and the phase
grows in proportion to it, $A(L,t)=A(0,t)\,e^{-\alpha L/2}e^{j\gamma|A(0,t)|^2L_{\mathrm{eff}}}$ with
$L_{\mathrm{eff}}=(1-e^{-\alpha L})/\alpha$. This self-phase modulation (SPM) chirps the pulse and
broadens its spectrum, with a number of spectral peaks that grows with the peak phase
$\phi_{\max}=\gamma P_0L_{\mathrm{eff}}$ (\cref{fig:spm}).

\begin{figure}[htbp]
\centering
\begin{tikzpicture}
\begin{groupplot}[group style={group size=3 by 1,horizontal sep=4mm,y descriptions at=edge left},
  width=5.1cm,height=4.4cm,xmin=-1.5,xmax=1.5,ymin=0,ymax=1.05,xlabel={$(\nu-\nu_0)T_0$}]
\nextgroupplot[ylabel={spectral density},title={$\phi_{\max}=0$}]
  \addplot[th=black!70] table[x=f,y=p0]{\datafile{spm_spectra}};
\nextgroupplot[title={$\phi_{\max}=1.5\pi$}]
  \addplot[th=accent] table[x=f,y=p15]{\datafile{spm_spectra}};
\nextgroupplot[title={$\phi_{\max}=3.5\pi$}]
  \addplot[th=accent2] table[x=f,y=p35]{\datafile{spm_spectra}};
\end{groupplot}
\end{tikzpicture}
\caption{SPM-broadened spectra of a Gaussian pulse computed by \code{propagate} with $D=0$,
the classic figure of Agrawal~\cite{agrawal2019}: $\phi_{\max}=(M-\frac12)\pi$ produces $M$ peaks.}
\label{fig:spm}
\end{figure}

\subsection{The split-step Fourier method}
With both terms present, \cref{eq:nlse} has no general closed form. Formally
$A(z+h)=e^{h(\hat D+\hat N)}A(z)$. The two operators do not commute, but each is trivial to
apply on its own: $\hat D$ is a multiplication in the frequency domain, and $\hat N$ a
multiplication in the time domain. The symmetric (Strang) splitting
\begin{equation}
  A(z+h)=e^{\frac h2\hat D}\,e^{h\hat N}\,e^{\frac h2\hat D}A(z)+\mathcal O(h^3)
  \label{eq:strang}
\end{equation}
is exact to second order, by the Baker--Campbell--Hausdorff formula: the first commutator
$[\hat D,\hat N]$ cancels by symmetry. Over a fixed length the $L/h$ local errors add up to a
global error $\mathcal O(h^2)$.

Attenuation is included in the half steps. The nonlinear phase is applied at mid-step to the
field whose power has been reduced by the first half step. The effective length that makes
this exact for a constant envelope is
\begin{equation}
  L_{\mathrm{eff}}(h)=\int_{-h/2}^{h/2}e^{-\alpha z'}dz'=\frac{2\sinh(\alpha h/2)}{\alpha} .
\end{equation}
For dual polarization the Manakov equation replaces $|A|^2$ by $|A_x|^2+|A_y|^2$ and $\gamma$ by
$\frac89\gamma$, the average over fast random polarization evolution.
\codeexcerpt{ssfm}{The split-step Fourier solver of the NLSE / Manakov equation (\code{propagate}).}

\begin{figure}[htbp]
\centering
\begin{tikzpicture}
\begin{groupplot}[group style={group size=2 by 1,horizontal sep=17mm},width=6.9cm,height=5.4cm]
\nextgroupplot[xmode=log,ymode=log,xlabel={step $h$ [km]},ylabel={relative error after $4L_D$},
  xmin=0.005,xmax=4,ymin=1e-6,ymax=1,legend pos=north west,title={convergence}]
  \addplot[sim=accent2,mark size=2pt] table[x=h,y=err]{\datafile{soliton_convergence}};
  \addplot[th=black!50,dashed,domain=0.005:3] {0.09*x^2};
  \legend{split-step,slope 2 ($\propto h^2$)}
  \draw[accent,densely dotted] (axis cs:2,1e-6) -- (axis cs:2,1) node[pos=0.5,left,font=\scriptsize,rotate=90,anchor=south]{default 0.5 step/km};
\nextgroupplot[xlabel={$t$ [ps]},ylabel={$|A|^2/P_0$},xmin=-30,xmax=30,ymin=0,ymax=1.05,
  legend pos=north east,title={after $4L_D$}]
  \addplot[th=black!70] table[x=t_ps,y=input]{\datafile{soliton_profiles}};
  \addplot[th=accent,dashed] table[x=t_ps,y=linear]{\datafile{soliton_profiles}};
  \addplot[sim={accent2,mark size=0.9pt}] table[x=t_ps,y=nlse]{\datafile{soliton_profiles}};
  \legend{input / exact,dispersion only,NLSE (split-step)}
\end{groupplot}
\end{tikzpicture}
\caption{Validation of the solver on the fundamental soliton, an exact solution of
\cref{eq:nlse} with $\alpha=0$. With $T_0=5$\,ps, $P_0=|\beta_2|/(\gamma T_0^2)=\solitonPzero$\,W
and $L=\solitonL$\,km $=4L_D$, the pulse must return unchanged. Left: the relative $L^2$ error
falls as $h^2$, confirming the order of \cref{eq:strang}. Right: dispersion alone would broaden
the pulse, while the balance with SPM preserves it.}
\label{fig:soliton}
\end{figure}

\begin{pitfall}[Choosing the step]
The default \prm{steps_per_km}\code{=0.5} (2-km steps) is adequate for typical launch powers
around 0\,dBm per channel, where the nonlinear length $1/(\gamma P)$ is hundreds of km. The
soliton of \cref{fig:soliton} carries 28\,dBm, a nonlinear length of about 1\,km. There the
default step leaves a 12\,\% error, and 16 steps per km are needed for $10^{-3}$. For high powers,
dense WDM or solitons, reduce the step until the result stops changing. Without \prm{nonlinear} the solver applies one exact
FFT for the whole span.
\end{pitfall}

\subsection{Polarization-mode dispersion}
Random birefringence rotates the state of polarization and delays the two principal states by
the differential group delay (DGD). The first-order model is
\begin{equation}
  \tilde{\vect A}_{\mathrm{out}}(\omega)=\mat U_2\begin{pmatrix}e^{j\omega\tau/2}&0\\0&e^{-j\omega\tau/2}\end{pmatrix}\mat U_1\tilde{\vect A}_{\mathrm{in}}(\omega),
  \qquad\tau=D_{\mathrm{PMD}}\sqrt{\ell},
\end{equation}
with $\mat U_{1,2}$ Haar-random unitary matrices. These are generated by a QR decomposition of
a complex Gaussian matrix with the phases of the diagonal of $\mat R$ removed.
\codeexcerpt{pmd}{Random polarization rotation and first-order PMD.}
Here $\tau$ is set to its mean value; a full statistical model would draw it from a Maxwellian
distribution.

% ==================================================================
\section{Optical amplification and OSNR}
\subsection{ASE noise}
An erbium-doped fiber amplifier of gain $G$ adds amplified spontaneous emission with power
spectral density, per polarization,
\begin{equation}
  S_{\mathrm{ASE}}=n_{\mathrm{sp}}\,h\nu\,(G-1),\qquad \mathrm{NF}=\frac{2n_{\mathrm{sp}}(G-1)}{G}\ \Rightarrow\ n_{\mathrm{sp}}=\frac{\mathrm{NF}\cdot G}{2(G-1)},
\end{equation}
where the noise figure NF is at least 2 (3\,dB) for $n_{\mathrm{sp}}\ge1$. In the simulation band
$f_s$, the ASE is circular Gaussian noise of variance $S_{\mathrm{ASE}}f_s$ per polarization.
\codeexcerpt{ase}{ASE density from gain and noise figure.}
\codeexcerpt{add_ase}{ASE noise in the simulated bandwidth.}

\subsection{Optical signal-to-noise ratio}
The OSNR is measured in a reference bandwidth $B_{\mathrm{ref}}=12.5$\,GHz (0.1\,nm at 1550\,nm),
counting the noise in both polarizations:
\begin{equation}
  \mathrm{OSNR}=\frac{P}{2S_{\mathrm{ASE}}B_{\mathrm{ref}}} .
\end{equation}
In a chain of $N$ identical spans, each followed by an amplifier compensating its loss
($G=e^{\alpha L_{\mathrm{span}}}$), the ASE adds up while the signal power is restored. With the
NF definition above, $2S_{\mathrm{ASE}}\approx\mathrm{NF}\,G\,h\nu$, so
\begin{equation}
  \mathrm{OSNR}_{\mathrm{dB}}\approx\underbrace{-10\log_{10}\frac{h\nu B_{\mathrm{ref}}}{1\,\mathrm{mW}}}_{58.0\ \mathrm{dB}}
  +P_{\mathrm{dBm}}-G_{\mathrm{dB}}-\mathrm{NF}_{\mathrm{dB}}-10\log_{10}N,
  \label{eq:osnr58}
\end{equation}
the ``58\,dB rule'' of link design.
\codeexcerpt{link_loop}{The multi-span link: propagate, amplify, accumulate the ASE density.}

\begin{figure}[htbp]
\centering
\begin{tikzpicture}
\begin{semilogxaxis}[width=10cm,height=5.4cm,xlabel={number of spans $N$},ylabel={OSNR [dB]},
  xmin=0.8,xmax=50,ymin=18,ymax=40,legend pos=north east]
  \addplot[th=accent] table[x=spans,y=rule58]{\datafile{osnr_spans}};
  \addplot[sim=accent2] table[x=spans,y=sim]{\datafile{osnr_spans}};
  \legend{58\,dB rule \eqref{eq:osnr58},\blk{optical_link} metric}
\end{semilogxaxis}
\end{tikzpicture}
\caption{OSNR after $N\times80$\,km (0.2\,dB/km, NF\,=\,5\,dB, 0\,dBm per channel). The block's
metric and the rule differ by less than 0.05\,dB, the effect of the $G-1\approx G$ approximation
and of the exact photon energy.}
\label{fig:osnr}
\end{figure}

\subsection{From OSNR to $E_s/N_0$}
For a coherent receiver whose noise is the ASE, the noise density per polarization is
$N_0=S_{\mathrm{ASE}}$ and the signal power per polarization is $P/p$ for $p$ polarizations.
Hence
\begin{equation}
  \frac{E_s}{N_0}=\frac{P/p}{S_{\mathrm{ASE}}R_s}=\mathrm{OSNR}\cdot\frac{2B_{\mathrm{ref}}}{p\,R_s},
\end{equation}
implemented as \code{osnr_to_esn0_db}. The test-suite uses it to check that the full coherent
DP-QPSK link (\cref{sec:coherent-dsp}) performs within about 0.3\,dB of theory.

% ==================================================================
\section{Wavelength-division multiplexing}
\blk{wdm_multiplexer} surrounds the channel under test with $N_{\mathrm{ch}}-1$ neighbours at
multiples of the channel spacing. The neighbours carry independent symbols with the same
constellation, pulse shape and power. Propagating the composite field through a nonlinear
fiber produces cross-phase modulation (the power of the neighbours modulates the phase of the
channel under test) and four-wave mixing. \blk{wdm_demultiplexer} selects a channel with a
second-order super-Gaussian optical filter, $|H(f)|=\exp\!\big(-\tfrac{\ln2}{2}(2f/B)^4\big)$, which
is 3\,dB down at $\pm B/2$. The sample rate must cover all channels. The \emph{WDM with fiber
nonlinearity} scenario shows the MER of the centre channel dropping by several dB when the
neighbours are switched on.
\codeexcerpt{wdm_demux}{Channel selection by frequency shift and optical filtering.}

% ==================================================================
\section{Free-space optics}
Atmospheric turbulence makes the received intensity fluctuate (scintillation). \blk{fso_turbulence}
multiplies the field by $\sqrt h$, where the normalized irradiance $h$ ($\E h=1$) is constant
over a frame and follows one of two models.
\begin{itemize}
  \item \textbf{Log-normal} (weak turbulence): $\ln h\sim\N(-\sigma^2/2,\sigma^2)$.
  \item \textbf{Gamma-gamma} (moderate to strong turbulence): $h$ is the product of two
  independent unit-mean gamma variables with parameters $\alpha$ (large eddies) and $\beta$
  (small eddies). Its density is
  \begin{equation}
    f(h)=\frac{2(\alpha\beta)^{\frac{\alpha+\beta}{2}}}{\Gamma(\alpha)\Gamma(\beta)}\,h^{\frac{\alpha+\beta}{2}-1}K_{\alpha-\beta}\big(2\sqrt{\alpha\beta h}\big).
  \end{equation}
\end{itemize}
\codeexcerpt{fso}{Scintillation models.}

\begin{figure}[htbp]
\centering
\begin{tikzpicture}
\begin{axis}[width=10cm,height=5.2cm,xlabel={normalized irradiance $h$},ylabel={pdf},xmin=0,xmax=3,ymin=0,
  legend pos=north east]
  \addplot[th=accent] table[x=h,y=ln_pdf]{\datafile{fso_pdf}};
  \addplot[sim=accent,mark size=1.2pt] table[x=h,y=ln_hist]{\datafile{fso_pdf}};
  \addplot[th=accent2] table[x=h,y=gg_pdf]{\datafile{fso_pdf}};
  \addplot[sim={accent2,mark=square*},mark size=1.1pt] table[x=h,y=gg_hist]{\datafile{fso_pdf}};
  \legend{log-normal $\sigma^2{=}0.1$,histogram,{gamma-gamma $\alpha{=}4$, $\beta{=}2$},histogram}
\end{axis}
\end{tikzpicture}
\caption{Histograms of $h$ from \blk{fso_turbulence} (60\,000 frames) against the closed-form densities.}
\label{fig:fso}
\end{figure}

% ==================================================================
\section{Direct detection}
\subsection{PIN and APD receivers}
A photodiode converts power into current, $I=\mathcal RP$, with responsivity $\mathcal R$
(A/W). Over the noise bandwidth $B=f_s/2$ the current carries three contributions:
\begin{itemize}
  \item \textbf{shot noise}, $\sigma_{\mathrm{sh}}^2=2qM^2F\,(\mathcal RP+I_d)B$;
  \item \textbf{thermal noise}, $\sigma_{\mathrm{th}}^2=i_n^2B$, with $i_n$ the input-referred
  current density (\prm{thermal_noise_pa_rthz});
  \item \textbf{avalanche gain} (APD): the multiplication $M$ comes with McIntyre's excess noise
  factor $F=kM+(1-k)(2-1/M)$, with $k$ the ionization ratio.
\end{itemize}
The APD gain raises the signal above thermal noise but amplifies shot noise by $M^2F$, so there
is an optimum $M$.
\codeexcerpt{photodiode}{Square-law detection with shot, thermal and APD excess noise.}

\subsection{OOK sensitivity and the Q factor}
For on--off keying with levels $I_{1,0}=\mathcal RP_{1,0}$ and Gaussian noise of standard
deviations $\sigma_{1,0}$, the optimum threshold gives
\begin{equation}
  \mathrm{BER}=\frac12\erfc\Big(\frac{Q}{\sqrt2}\Big),\qquad Q=\frac{I_1-I_0}{\sigma_1+\sigma_0},
  \qquad P_{1,0}=\frac{2P_{\mathrm{avg}}}{1+\mathrm{ER}^{\mp1}} .
\end{equation}
Thermal noise is signal independent, so $\sigma_1=\sigma_0$ and $Q\propto P_{\mathrm{avg}}$: the BER
falls by many decades within a few dB (\cref{fig:ook}). The receiver sensitivity is
conventionally quoted at $10^{-9}$ ($Q=6$) or at a FEC threshold.

\begin{figure}[htbp]
\centering
\begin{tikzpicture}
\begin{semilogyaxis}[width=10cm,height=6cm,xlabel={average received power [dBm]},ylabel={BER},
  xmin=-28,xmax=-18,ymin=1e-7,ymax=0.5,legend pos=south west]
  \addplot[th=accent] table[x=p,y=ber]{\datafile{ook_theory}};
  \addplot[sim=accent2] table[x=p,y=ber]{\datafile{ook_sim}};
  \legend{Q-factor theory,simulated}
\end{semilogyaxis}
\end{tikzpicture}
\caption{10\,Gb/s NRZ-OOK into a thermal-noise-limited PIN receiver ($\mathcal R=0.8$\,A/W,
20\,pA/$\sqrt{\mathrm{Hz}}$, ER\,=\,10\,dB): chain \blk{intensity_modulator}$\to$\blk{photodiode}$\to$\blk{agc}$\to$\blk{line_decoder}
against the Q-factor formula.}
\label{fig:ook}
\end{figure}

\begin{incode}
\blk{agc} removes the mean (the DC photocurrent) and scales each frame to the nominal symbol
power of the upstream line code or constellation. It subtracts the known noise variance when
estimating the signal power, so the decision thresholds stay correct at low SNR.
\end{incode}

% ==================================================================
\section{Coherent detection}
\subsection{The coherent front end}
A coherent receiver mixes the signal with a local oscillator (LO) in a $90^\circ$ hybrid and
detects the four outputs with two balanced photodiode pairs per polarization
(\cref{fig:coherent-fe}). Balanced detection cancels the direct-detection terms $|E|^2$ and
$|E_{\mathrm{LO}}|^2$ and leaves the beat term:
\begin{equation}
  i_I+j\,i_Q=\mathcal R\sqrt{P_{\mathrm{LO}}}\;E(t)\,e^{-j(2\pi\Delta f\,t+\phi_{\mathrm{LO}}(t))}+n(t).
\end{equation}
The field is thus translated linearly to the electrical domain, including its phase and
polarization. With a strong LO, the LO shot noise $2q\mathcal RP_{\mathrm{LO}}B$ dominates the
thermal noise.

\begin{figure}[htbp]
\centering
\begin{tikzpicture}[node distance=5mm and 8mm]
  \node[blockbox] (pbs) {PBS};
  \node[blockbox,right=of pbs,yshift=10mm,minimum height=11mm] (hx) {$90^\circ$\\hybrid X};
  \node[blockbox,right=of pbs,yshift=-10mm,minimum height=11mm] (hy) {$90^\circ$\\hybrid Y};
  \node[warmbox,above=4mm of hx] (lox) {LO ($\Delta f$, $\Delta\nu$)};
  \node[warmbox,below=4mm of hy] (loy) {LO ($\Delta f$, $\Delta\nu$)};
  \node[greenbox,right=of hx,minimum width=13mm] (bx) {2 balanced\\PDs};
  \node[greenbox,right=of hy,minimum width=13mm] (by) {2 balanced\\PDs};
  \node[blockbox,right=of bx] (ax) {ADC};
  \node[blockbox,right=of by] (ay) {ADC};
  \node[left=6mm of pbs,lbl] (in) {$E_x,E_y$};
  \draw[sig] (in)--(pbs); \draw[sig] (pbs.east)--(hx.west); \draw[sig] (pbs.east)--(hy.west);
  \draw[sig] (lox)--(hx); \draw[sig] (loy)--(hy);
  \draw[sig] (hx)--(bx); \draw[sig] (hy)--(by); \draw[sig] (bx)--node[above,lbl]{$I_x,Q_x$}(ax); \draw[sig] (by)--node[above,lbl]{$I_y,Q_y$}(ay);
\end{tikzpicture}
\caption{Polarization-diverse coherent front end, modelled by \blk{coherent_receiver}. One LO laser
is split between the two hybrids.}
\label{fig:coherent-fe}
\end{figure}

\begin{incode}
The simulator processes each frame with circular FFTs (fiber, dispersion compensation). An LO
phase that does not return to its initial value at the end of the frame would create a
discontinuity at the frame boundary, spread over the frame by the dispersion compensator. The
block therefore makes the LO phase periodic:
\begin{itemize}
  \item the frequency offset is rounded to the frame's frequency grid $f_s/n$;
  \item the Wiener phase is turned into a Brownian bridge, $\phi_k-\frac kn\phi_n$, which has the
  same local increments but returns to zero.
\end{itemize}
\end{incode}
\codeexcerpt{coherent_rx}{Coherent detection with a periodic LO phase (\blk{coherent_receiver}).}

\subsection{Coherent digital signal processing}\label{sec:coherent-dsp}
The chain follows the architecture described by Savory~\cite{savory2010}.
\Cref{fig:coherent-const} follows a 128\,Gb/s DP-QPSK signal through the DSP chain after
$20\times80$\,km, with random polarization, a 1\,GHz LO offset and 100\,kHz lasers.

\begin{figure}[htbp]
\centering
\begin{tikzpicture}
\begin{groupplot}[group style={group size=4 by 1,horizontal sep=4mm,y descriptions at=edge left},
  constaxis,width=4.1cm,height=4.1cm,xmin=-2.2,xmax=2.2,ymin=-2.2,ymax=2.2,xtick={-2,0,2},ytick={-2,0,2}]
\nextgroupplot[title={front-end output}]
  \addplot[scatterpts=black!50] table[x=i,y=q]{\datafile{coh_rx}};
\nextgroupplot[title={after CD + FO comp.}]
  \addplot[scatterpts=accent4] table[x=i,y=q]{\datafile{coh_fo}};
\nextgroupplot[title={after $2{\times}2$ equalizer}]
  \addplot[scatterpts=accent] table[x=i,y=q]{\datafile{coh_eq}};
\nextgroupplot[title={after phase recovery}]
  \addplot[scatterpts=accent3] table[x=i,y=q]{\datafile{coh_cpr}};
\end{groupplot}
\end{tikzpicture}
\caption{Polarization X of the \emph{Coherent DP-QPSK long haul} scenario at four points of
the DSP chain (normalized to unit power).}
\label{fig:coherent-const}
\end{figure}

\paragraph{Chromatic dispersion compensation.}
The accumulated dispersion is known from the metadata (or given), and its all-pass response
\cref{eq:dispersion-operator} is inverted exactly by multiplying with its conjugate in the
frequency domain.
\codeexcerpt{cd_comp}{Frequency-domain CD compensation.}

\paragraph{Frequency-offset estimation.}
Raising QPSK to the fourth power removes the modulation: $(s\,e^{j2\pi\Delta f t})^4=s^4e^{j2\pi4\Delta ft}$
with $s^4=-1$ for every QPSK symbol, and approximately constant on average for square QAM. The
spectrum of $r^4$ therefore has a line at $4\Delta f$. The estimate is the position of its peak
divided by 4, refined by zero-padding and computed on a low-pass filtered copy to reject noise.
In dual polarization a rotated mixture $as_1+bs_2$ has a fourth-power line proportional to
$a^4+b^4$, which can vanish. The estimator therefore also uses the combinations
$(x_0+x_1)/\sqrt2$ and $(x_0+jx_1)/\sqrt2$.

Compensating the offset \emph{after} the dispersion compensator leaves a residual: the CD
filter was applied to a spectrum shifted by $\Delta\omega$, so the error is
$\exp\!\big(j\frac{\beta_2L}{2}(2\omega\Delta\omega+\Delta\omega^2)\big)$, a bulk delay
$\beta_2L\Delta\omega$ plus a constant phase. The block removes it.
\codeexcerpt{fo_est}{Fourth-power FO estimation over several polarization combinations.}

\paragraph{Adaptive $2\times2$ equalization.}
Polarization rotation, PMD and residual dispersion form a $2\times2$ MIMO channel, inverted by
a fractionally spaced butterfly of four FIR filters (\cref{fig:butterfly}):
\begin{equation}
  y_p[k]=\sum_{q\in\{x,y\}}\vect w_{pq}\T\,\vect u_q[k],\qquad
  \vect w_{pq}\leftarrow\vect w_{pq}+\mu\,e_p[k]\,\vect u_q\conj[k] .
\end{equation}
The error $e_p$ is either:
\begin{itemize}
  \item the constant-modulus error $y_p(R-|y_p|^2)$ with $R=\E|s|^4/\E|s|^2$ (Godard's
  CMA~\cite{godard1980}), which is blind and insensitive to the carrier phase; or
  \item the LMS error $d_p-y_p$, trained on known symbols and then decision-directed, with a
  larger step during training.
\end{itemize}

\begin{figure}[htbp]
\centering
\begin{tikzpicture}[x=1cm,y=1cm]
  % input x feeds w_xx (to x out) and w_yx (to y out); input y feeds w_xy and w_yy
  \node[lbl,anchor=east] (xi) at (0,0) {$x$ in};
  \node[lbl,anchor=east] (yi) at (0,-3.3) {$y$ in};
  \node[blockbox,minimum width=11mm] (wxx) at (2.2,0.2) {$\vect w_{xx}$};
  \node[blockbox,minimum width=11mm] (wyx) at (2.2,-1.0) {$\vect w_{yx}$};
  \node[blockbox,minimum width=11mm] (wxy) at (2.2,-2.3) {$\vect w_{xy}$};
  \node[blockbox,minimum width=11mm] (wyy) at (2.2,-3.5) {$\vect w_{yy}$};
  \node[sumnode] (sx) at (4.4,0) {$+$};
  \node[sumnode] (sy) at (4.4,-3.3) {$+$};
  \node[lbl,anchor=west] (xo) at (5.2,0) {$x$ out};
  \node[lbl,anchor=west] (yo) at (5.2,-3.3) {$y$ out};
  \coordinate (bx) at (0.7,0); \coordinate (by) at (0.7,-3.3);
  \draw[thick] (xi.east) -- (bx); \draw[thick] (yi.east) -- (by);
  \draw[sig] (bx) |- (wxx.west); \draw[sig] (bx) |- (wyx.west);
  \draw[sig] (by) |- (wxy.west); \draw[sig] (by) |- (wyy.west);
  \fill (bx) circle (1.2pt); \fill (by) circle (1.2pt);
  \draw[sig] (wxx.east) -- (sx); \draw[sig] (wxy.east) -- (sx.south west);
  \draw[sig] (wyy.east) -- (sy); \draw[sig] (wyx.east) -- (sy.north west);
  \draw[sig] (sx) -- (xo); \draw[sig] (sy) -- (yo);
\end{tikzpicture}
\caption{The butterfly equalizer (2 samples per symbol, \prm{n_taps} taps per filter).}
\label{fig:butterfly}
\end{figure}
\codeexcerpt{butterfly}{Butterfly update with CMA or trained/decision-directed LMS.}

\paragraph{Carrier phase recovery.}
Laser phase noise of $2\times100$\,kHz at 32\,GBd is slow, but large enough to need
feed-forward tracking. Two estimators are available.
\begin{itemize}
  \item The Viterbi--Viterbi estimator removes the QPSK modulation with the fourth power:
  $\hat\phi=\frac14\big(\arg\sum_{|i|\le W}y_{k+i}^4-\pi\big)$.
  \item Blind phase search (BPS)~\cite{pfau2009} works for any square QAM. It rotates the
  symbols by $B$ test phases in $[-\pi/4,\pi/4)$, sums over a window the squared distances to
  the nearest constellation points, and keeps the phase with the smallest sum.
\end{itemize}
Both estimates are defined modulo $\pi/2$; they are unwrapped (on $4\hat\phi$) to follow the phase
continuously. The remaining quadrant ambiguity is resolved once per frame with a short block of
known symbols, taken after the equalizer has converged.
\codeexcerpt{bps}{Viterbi--Viterbi and blind phase search.}

\chapter{Validation}\label{ch:validation}

A simulator is only useful if its numbers can be trusted. LinkStudio is validated at three levels:
\begin{itemize}
  \item unit tests of the engine and the DSP helpers;
  \item \emph{physics tests}, which compare complete chains with closed-form results;
  \item a test that loads and runs every shipped scenario.
\end{itemize}
The demo scripts repeat the physics checks in readable form, and this manual's figures are a
third, visual check of the same models.

\section{Running the checks}
\begin{lstlisting}[style=shell]
cd backend
python -m pytest                 # all tests (about 220, a few minutes)
python -m pytest tests/test_physics.py -q
python scripts/main.py           # demos: simulated vs theory, PASS/FAIL per check
python scripts/main.py --quick   # fewer errors per point, faster
\end{lstlisting}
Physics tests stop after 300--500 errors and accept a relative deviation of 20--25\,\%. By
\cref{sec:mc} this is several standard deviations, so false alarms are negligible, while a
systematic error of 1\,dB (a factor of 2--3 in BER on steep curves) is caught.

\section{Physics checks}
\begin{center}
\small
\begin{longtable}{@{}p{3.3cm}p{5.6cm}p{5.3cm}@{}}
\toprule
\textbf{Model} & \textbf{Reference} & \textbf{Checked in}\\
\midrule
\endhead
AWGN, PSK/QAM/PAM & exact Gray BER (\cref{sec:ber-awgn}) & \file{test_physics.py}, \cref{fig:ber-awgn}\\
Soft demapping & max-log decisions = hard decisions & \file{test_physics.py}\\
Pulse shaping & RRC link BER = symbol-rate theory; noise metadata after MF & \file{test_physics.py}\\
OFDM over AWGN & per-subcarrier $E_s/N_0$ theory & \file{test_physics.py}\\
OFDM multipath & LS and perfect CSI equalization & \file{test_physics.py}, \cref{fig:ofdm-chest}\\
Rayleigh + ZF & \cref{eq:rayleigh-ber} & \file{test_physics.py}, \cref{fig:diversity}\\
ISI equalizers & ZF/MMSE/LMS remove ISI & \file{test_physics.py}\\
Soft vs hard Viterbi & soft better by $\approx2$\,dB & \file{test_physics.py}\\
Synchronization & receiver close to theory with delay and CFO; fails without sync & \file{test_advanced.py}\\
Carrier recovery & tracks phase noise & \file{test_advanced.py}\\
Timing recovery & estimates the fractional delay & \file{test_advanced.py}\\
LDPC & codewords satisfy $\mat H\vect c\T=0$; corrects noise; coding gain & \file{test_advanced.py}\\
Polar & noiseless decoding; frozen set; coding gain & \file{test_advanced.py}\\
Puncturing & rate 3/4; decodable & \file{test_advanced.py}\\
SIMO MRC & \cref{eq:mrc}, $L=2$ & \file{test_advanced.py}, \cref{fig:diversity}\\
Spatial multiplexing & ZF $2\times2$ = SISO Rayleigh; ML $<$ MMSE $<$ ZF & \file{test_advanced.py}, \cref{fig:smux}\\
Alamouti & MRC $L=2$ shifted by 3\,dB & \file{test_advanced.py}, \cref{fig:diversity}\\
Time-varying multipath & OFDM with Jakes taps & \file{test_advanced.py}\\
Line codes & round trips; 4B/5B and 64B/66B rates; error multiplication & \file{test_wireline_optical.py}\\
Real NRZ over AWGN & BPSK theory & \file{test_wireline_optical.py}\\
Cable & insertion loss equals the preset formula & \file{test_wireline_optical.py}, \cref{fig:cable}\\
SerDes & CTLE peaking; CTLE + DFE open the eye & \file{test_wireline_optical.py}, \cref{fig:ctle}\\
DMT & water-filling bit loading & \file{test_wireline_optical.py}, \cref{fig:dmt}\\
Crosstalk, echo & coupling level; canceller reduces echo & \file{test_wireline_optical.py}\\
Jitter & zero jitter is identity; RJ degrades & \file{test_wireline_optical.py}\\
Fiber loss, metadata & $P\,e^{-\alpha L}$, $D\ell$ & \file{test_wireline_optical.py}\\
Dispersion & Gaussian broadening; exact compensation & \file{test_wireline_optical.py}, \cref{fig:dispersion}\\
SPM & CW phase $\gamma PL_{\mathrm{eff}}$ & \file{test_wireline_optical.py}, \cref{fig:spm}\\
Split-step solver & second-order convergence on the soliton & \cref{fig:soliton}\\
EDFA, multi-span link & OSNR formulas ($n_{sp}=F/2$, $GA\ne1$, course exercise), saturation & \file{test_optical_models.py}, \cref{fig:osnr-spans,fig:osnr-span-length}\\
Raman & on/off gain, lumped-limit ASE, $F_{\mathrm{eq}}<3$\,dB & \file{test_optical_models.py}, \cref{fig:raman}\\
GN model & long-span form, 3\,dB rule, link metrics; split-step comparison & \file{test_optical_models.py}, \cref{fig:gn-launch}\\
Lasers & threshold and slope, settling, linearized response, adiabatic chirp, VCSEL overlaps, thermal rollover, hole-burning turn-off & \file{test_lasers.py}, \cref{fig:vcsel-li,fig:vcsel-transient,fig:laser-response}\\
Modulators, detectors & MZM sine law and ER, PAM levels, TIA output, APD power law, ISNR & \file{test_optical_models.py}, \cref{fig:snr-receivers}\\
PM-QAM & course BER formulas vs exact; back-to-back simulation & \file{test_optical_models.py}, \cref{fig:coherent-ber}\\
IM/DD OOK, PAM-4 & Q-factor BER; PAM-M formula of the course; dispersion limit & \file{test_wireline_optical.py}, \file{test_optical_models.py}, \cref{fig:pam-imdd,fig:cd-limit}\\
IM/DD PAM-4 & dispersion-induced power fading & \file{test_wireline_optical.py}\\
Coherent DP-QPSK & BER vs OSNR theory within $\approx0.3$\,dB; FO estimate & \file{test_wireline_optical.py}\\
FSO & mean and variance of the scintillation & \file{test_wireline_optical.py}, \cref{fig:fso}\\
WDM & XPM lowers the MER of the centre channel & \file{test_wireline_optical.py}\\
Multimode fiber & $\alpha_{\mathrm{opt}}$, $\sigma_{\min}$, WKB delay, continuum limit, class bandwidths, launch closed forms, variance additivity & \file{test_multimode.py}, \cref{fig:mmf-profile,fig:mmf-launch}\\
Data-center lane & Bessel--Thomson noise bandwidth; reach by fiber class & \file{test_multimode.py}, \cref{fig:mmf-eyes}\\
Fading generator & Zheng--Xiao: $R_h=J_0$, uncorrelated I/Q & \file{test_channels.py}, \cref{fig:jakes}\\
NR TDL profiles & unit normalized delay spread; TDL-D $K_1=13.3$\,dB; fractional delay & \file{test_channels.py}\\
Path loss & P.525; TR~38.901 breakpoint continuity, NLOS $\ge$ LOS, shadowing, LOS probability & \file{test_propagation.py}, \cref{fig:pathloss}\\
Thermal noise & $k_BTF f_s$; free-space QPSK link budget BER & \file{test_propagation.py}\\
Copper loops & G.996.1 Mid-C insertion loss within 0.35\,dB; causal dielectric & \file{test_wireline_optical.py}, \cref{fig:pic}\\
Dispersion slope & group delay $D\ell\Delta\lambda+\frac{S\ell}2\Delta\lambda^2$; CD+slope compensation & \file{test_wireline_optical.py}\\
PMD & Maxwellian DGD (mean, $\langle\tau^2\rangle/\langle\tau\rangle^2$) & \file{test_wireline_optical.py}, \cref{fig:pmd}\\
Adaptive split step & soliton error $<10^{-3}$ with default criterion & \file{test_wireline_optical.py}\\
DML, coherent RX & adiabatic chirp frequency step; LO split per polarization & \file{test_wireline_optical.py}\\
FSO channel & $\alpha,\beta$ from $\sigma_R^2$; Kim, geometric and pointing losses & \file{test_wireline_optical.py}\\
CP-based OFDM sync & recovers delay and CFO & \file{test_wireline_optical.py}\\
Confidence interval & Wilson bounds contain the BER & \file{test_wireline_optical.py}\\
All scenarios & load, validate, run; BER below a threshold & \file{test_scenarios.py}\\
\bottomrule
\end{longtable}
\end{center}

\section{Known limitations}
Every model in this manual states its assumptions; the most important simplifications are
collected here.
\begin{itemize}
  \item Laser dynamics are small-signal (no large-signal rate equations). The multimode fiber
  omits mode-partition and modal noise.
  \item FSO fading is quasi-static (no time correlation, aperture averaging or beam wander).
  \item WDM simulates the full multi-channel field but demodulates only one channel.
  \item The OFDM CP synchronizer estimates the fractional CFO only.
  \item DMT adapts bits to a flat transmit PSD (no power water-filling).
  \item MIMO channels are flat; frequency-selective MIMO-OFDM is not yet a block.
  \item Subsystems are edited in their YAML files, not in the GUI.
  \item Error counts assume independent errors; for coded and block-fading links, also judge
  the accuracy by the number of frame errors.
\end{itemize}


\appendix
\part{Reference}
\chapter{Block reference}\label{app:blocks}

This appendix lists the \numblocks{} registered blocks, grouped by functional class in signal-chain order as in the GUI palette. It is generated by
\file{manual/tools/generate.py} from \code{block_catalogue()}, the same source as the GUI
palette and \code{python -m linkstudio describe}. Each entry gives:
\begin{itemize}
  \item the first paragraph of the block documentation;
  \item the input and output ports with their accepted kinds;
  \item the parameters with their defaults. \emph{auto} means the value is taken from upstream
  metadata or measured; \emph{required} means the parameter has no default.
\end{itemize}

\makeatletter
\newcommand{\blockcategory}[1]{%
  \section{\texorpdfstring{\MakeUppercase{#1}}{#1}}}
\newcommand{\blockentry}[5]{%
  \par\noindent
  \begin{minipage}{\linewidth}\raggedright
    {\sffamily\bfseries\color{accent}\texttt{#1}}\quad{\small #2}\\[1pt]
    {\footnotesize\textsf{in:} \texttt{#3}\qquad\textsf{out:} \texttt{#4}}\\[1pt]
    {\footnotesize\color{commentcolor}#5}
  \end{minipage}\par\vspace{5pt}}
\makeatother

% Generated by manual/tools/generate.py from linkstudio.core.registry.block_catalogue()
\blockcategory{source}
\blockentry{prbs}{Pseudo-random binary sequence (PRBS-7/9/11/15/20/23/31), continuous across frames.}{--}{out:BITS}{\texttt{order}=9; \texttt{n\_bits}=1024}
\blockentry{random\_bits}{Independent Bernoulli bits, a new random frame for every batch row.}{--}{out:BITS}{\texttt{n\_bits}=1024; \texttt{p\_one}=0.5}
\blockcategory{modulation}
\blockentry{demapper}{Hard or soft (LLR) demapping.}{in:SYMBOLS}{out:BITS|LLR}{\texttt{method}=maxlog; \texttt{scheme}=auto; \texttt{order}=4; \texttt{phase\_offset\_deg}=auto; \texttt{points}=auto; \texttt{noise\_var}=auto}
\blockentry{mapper}{Map groups of bits onto Gray-labelled constellation points (unit average energy).}{in:BITS}{out:SYMBOLS}{\texttt{scheme}=qam; \texttt{order}=4; \texttt{phase\_offset\_deg}=auto; \texttt{points}=auto; \texttt{symbol\_rate}=1e+06}
\blockcategory{waveform}
\blockentry{downsampler}{Keep one sample out of \texttt{factor} starting at \texttt{offset}.}{in:SYMBOLS|SAMPLES}{out:SYMBOLS|SAMPLES}{\texttt{factor}=2; \texttt{offset}=0; \texttt{output\_kind}=samples}
\blockentry{fir\_filter}{Generic FIR filter with user-supplied taps (causal, output length preserved).}{in:SYMBOLS|SAMPLES}{out:same as 'in'}{\texttt{taps}=[1.0]}
\blockentry{matched\_filter}{Receive filter + symbol-rate sampling.}{in:SAMPLES}{out:SYMBOLS|SAMPLES}{\texttt{shape}=matched; \texttt{sps}=auto; \texttt{span}=10; \texttt{rolloff}=0.25; \texttt{bt}=0.3; \texttt{normalize}=True; \texttt{decimate}=True}
\blockentry{ofdm\_demodulator}{Remove the cyclic prefix and FFT every OFDM symbol into a resource grid.}{in:SAMPLES}{out:GRID}{\texttt{n\_fft}=auto; \texttt{cp\_len}=16; \texttt{n\_active}=52; \texttt{pilot\_spacing}=6; \texttt{fft\_offset}=0}
\blockentry{ofdm\_modulator}{Map data symbols and pilots on subcarriers, IFFT and add the cyclic prefix.}{in:SYMBOLS}{out:SAMPLES}{\texttt{n\_fft}=64; \texttt{cp\_len}=16; \texttt{n\_active}=52; \texttt{pilot\_spacing}=6; \texttt{subcarrier\_spacing}=15000; \texttt{tail\_zeros}=0}
\blockentry{pulse\_shaper}{Upsample symbols by \texttt{sps} and filter them with a unit-energy pulse.}{in:SYMBOLS}{out:SAMPLES}{\texttt{shape}=rrc; \texttt{sps}=8; \texttt{span}=10; \texttt{rolloff}=0.25; \texttt{bt}=0.3}
\blockentry{upsampler}{Insert \texttt{factor - 1} zeros between samples.}{in:SYMBOLS|SAMPLES}{out:SAMPLES}{\texttt{factor}=2}
\blockcategory{channel}
\blockentry{awgn}{Add circularly-symmetric complex Gaussian noise.}{in:SYMBOLS|SAMPLES}{out:same as 'in'}{\texttt{snr\_db}=10; \texttt{mode}=ebn0; \texttt{power\_reference}=nominal; \texttt{signal\_power}=auto; \texttt{bits\_per\_symbol}=auto; \texttt{code\_rate}=auto; \texttt{noise\_figure\_db}=5; \texttt{temperature\_k}=290}
\blockentry{carrier\_offset}{Carrier frequency offset and phase rotation: y = x * exp(j(2 pi cfo t + phi)).}{in:SYMBOLS|SAMPLES}{out:same as 'in'}{\texttt{cfo\_hz}=0; \texttt{phase\_deg}=0; \texttt{random\_phase}=False}
\blockentry{flat\_fading}{Frequency-flat Rayleigh or Rician fading, E|h|\textasciicircum{}2 = 1.}{in:SYMBOLS|SAMPLES}{out:same as 'in', csi:CSI}{\texttt{model}=rayleigh; \texttt{k\_factor\_db}=0; \texttt{variation}=block; \texttt{doppler\_hz}=100; \texttt{n\_sinusoids}=16}
\blockentry{gain}{Scale the signal by \texttt{gain\_db} (e.g. path loss). Noise metadata is scaled too.}{in:SYMBOLS|SAMPLES}{out:same as 'in'}{\texttt{gain\_db}=0}
\blockentry{iq\_imbalance}{Receiver IQ imbalance and DC offset: y = K1 x + K2 conj(x) + dc.}{in:SYMBOLS|SAMPLES}{out:same as 'in'}{\texttt{amplitude\_db}=0; \texttt{phase\_deg}=0; \texttt{dc\_i}=0; \texttt{dc\_q}=0}
\blockentry{mimo\_channel}{Flat Rayleigh MIMO channel y = H x, H of size Nr x Nt, normalized by 1/sqrt(Nt).}{in:SYMBOLS|SAMPLES}{out:same as 'in', csi:CSI}{\texttt{n\_rx}=2; \texttt{variation}=block; \texttt{tx\_correlation}=0; \texttt{rx\_correlation}=0}
\blockentry{multipath}{Tapped-delay-line channel (frequency selective).}{in:SYMBOLS|SAMPLES}{out:same as 'in', csi:CSI}{\texttt{preset}=auto; \texttt{delay\_spread\_ns}=100; \texttt{fractional\_delays}=True; \texttt{delays}=[0.0, 1.0, 2.0]; \texttt{powers\_db}=[0.0, -3.0, -6.0]; \texttt{delay\_unit}=samples; \texttt{fading}=True; \texttt{variation}=block; \texttt{doppler\_hz}=50; \texttt{normalize}=True}
\blockentry{path\_loss}{Distance-dependent path loss and log-normal shadow fading (one value per frame).}{in:SYMBOLS|SAMPLES}{out:same as 'in'}{\texttt{model}=free\_space; \texttt{frequency\_ghz}=3.5; \texttt{distance\_m}=100; \texttt{h\_bs\_m}=auto; \texttt{h\_ut\_m}=auto; \texttt{condition}=probabilistic; \texttt{shadow\_fading}=True}
\blockentry{phase\_noise}{Oscillator phase noise modelled as a Wiener process with 3-dB \texttt{linewidth\_hz}.}{in:SYMBOLS|SAMPLES}{out:same as 'in'}{\texttt{linewidth\_hz}=100}
\blockentry{power\_amplifier}{Memoryless PA nonlinearity (Rapp, Saleh or ideal clipper).}{in:SYMBOLS|SAMPLES}{out:same as 'in'}{\texttt{model}=rapp; \texttt{input\_backoff\_db}=6; \texttt{smoothness}=2; \texttt{saleh\_beta\_a}=1.1517; \texttt{saleh\_alpha\_p}=4.0033; \texttt{saleh\_beta\_p}=9.104; \texttt{normalize\_output}=True}
\blockentry{quantizer}{Uniform mid-rise ADC on I and Q with clipping.}{in:SYMBOLS|SAMPLES}{out:same as 'in'}{\texttt{bits}=8; \texttt{full\_scale}=auto; \texttt{loading\_factor}=4}
\blockentry{timing\_offset}{Delay the signal by a (fractional) number of samples, keeping its length.}{in:SYMBOLS|SAMPLES}{out:same as 'in'}{\texttt{delay\_samples}=0}
\blockcategory{sync}
\blockentry{carrier\_recovery}{Decision-directed second-order PLL tracking residual phase and frequency offsets.}{in:SYMBOLS}{out:SYMBOLS}{\texttt{loop\_bandwidth}=0.005; \texttt{damping}=0.707}
\blockentry{fo\_compensation}{Blind frequency-offset estimation from the spectral peak of the M-th power (M = 4 for QPSK / square QAM), joint over polarizations, and correction.}{in:SYMBOLS|SAMPLES}{out:same as 'in'}{\texttt{power}=4; \texttt{estimation\_bandwidth}=0.75}
\blockentry{frame\_sync}{Preamble-based frame synchronization.}{in:SYMBOLS}{out:SYMBOLS}{\texttt{search}=32; \texttt{correct\_cfo}=True}
\blockentry{ofdm\_sync}{Blind OFDM timing and fractional CFO synchronization from the cyclic prefix (van de Beek).}{in:SAMPLES}{out:SAMPLES}{\texttt{max\_delay}=auto; \texttt{backoff}=auto; \texttt{correct\_cfo}=True}
\blockentry{phase\_recovery}{Feed-forward carrier phase recovery for large laser linewidths.}{in:SYMBOLS, ref:SYMBOLS}{out:SYMBOLS}{\texttt{method}=bps; \texttt{n\_test\_phases}=32; \texttt{half\_window}=16; \texttt{n\_ref}=64; \texttt{ref\_start}=600}
\blockentry{preamble}{Prepend a known preamble and append guard symbols to every frame.}{in:SYMBOLS}{out:SYMBOLS}{\texttt{length}=64; \texttt{guard}=16; \texttt{seed}=2024}
\blockentry{timing\_recovery}{Symbol timing recovery: Gardner detector, cubic interpolator, second-order loop.}{in:SAMPLES}{out:SYMBOLS}{\texttt{loop\_bandwidth}=0.01; \texttt{damping}=0.707; \texttt{detector\_gain}=1; \texttt{initial\_estimate}=True; \texttt{estimate\_symbols}=128}
\blockcategory{receiver}
\blockentry{adaptive\_mimo\_equalizer}{Fractionally spaced 2x2 (NxN) butterfly FIR equalizer for coherent receivers.}{in:SAMPLES, ref:SYMBOLS}{out:SYMBOLS}{\texttt{n\_taps}=15; \texttt{step\_size}=0.001; \texttt{train\_step\_size}=0.01; \texttt{algorithm}=lms; \texttt{n\_train}=512; \texttt{sps}=2}
\blockentry{agc}{Automatic gain control: remove the mean (optional) and scale every frame and stream to the nominal symbol power (\texttt{power} metadata of the upstream mapper / line code).}{in:SYMBOLS|SAMPLES}{out:same as 'in'}{\texttt{remove\_dc}=True; \texttt{target\_power}=auto}
\blockentry{linear\_equalizer}{FIR equalizer designed from the known channel impulse response (\texttt{csi} input).}{in:SYMBOLS, csi:CSI}{out:SYMBOLS}{\texttt{n\_taps}=15; \texttt{method}=mmse; \texttt{delay}=auto}
\blockentry{lms\_equalizer}{Adaptive (N)LMS FIR equalizer: trained on known symbols, then decision directed.}{in:SYMBOLS, ref:SYMBOLS}{out:SYMBOLS}{\texttt{n\_taps}=11; \texttt{step\_size}=0.1; \texttt{normalized}=True; \texttt{n\_train}=200; \texttt{decision\_directed}=True}
\blockentry{mimo\_detector}{Linear (ZF, MMSE) or maximum-likelihood MIMO detection with known channel.}{in:SYMBOLS, csi:CSI}{out:SYMBOLS}{\texttt{method}=mmse}
\blockentry{ofdm\_channel\_estimator}{Estimate the channel frequency response on every subcarrier.}{in:GRID, csi:CSI}{out:CSI}{\texttt{method}=ls; \texttt{average\_symbols}=False}
\blockentry{ofdm\_equalizer}{Extract the data subcarriers and equalize them (one-tap ZF per subcarrier).}{in:GRID, csi:CSI}{out:SYMBOLS}{--}
\blockentry{one\_tap\_equalizer}{Coherent equalization of flat fading with known gains (\texttt{csi} input).}{in:SYMBOLS, csi:CSI}{out:SYMBOLS}{\texttt{method}=zf}
\blockcategory{mimo}
\blockentry{alamouti\_decoder}{Linear Alamouti combining with known (block-fading) channel and any number of RX antennas.}{in:SYMBOLS, csi:CSI}{out:SYMBOLS}{--}
\blockentry{alamouti\_encoder}{Alamouti space-time block code: 2 transmit antennas, rate 1, full diversity.}{in:SYMBOLS}{out:SYMBOLS}{--}
\blockentry{stream\_demux}{Spatial multiplexing: distribute symbols round-robin over \texttt{n\_streams} antennas.}{in:SYMBOLS}{out:SYMBOLS}{\texttt{n\_streams}=2}
\blockentry{stream\_mux}{Inverse of \texttt{stream\_demux}: interleave the streams back into one sequence.}{in:SYMBOLS}{out:SYMBOLS}{--}
\blockcategory{coding}
\blockentry{conv\_encoder}{Feed-forward convolutional encoder of rate 1/n.}{in:BITS}{out:BITS}{\texttt{generators}=[171, 133]; \texttt{constraint\_length}=7; \texttt{terminate}=True}
\blockentry{crc\_decoder}{Check and strip the CRC; reports the frame error rate detected by the CRC.}{in:BITS|LLR}{out:BITS, valid:FLAGS}{\texttt{polynomial}=crc16}
\blockentry{crc\_encoder}{Append CRC parity bits to every frame.}{in:BITS}{out:BITS}{\texttt{polynomial}=crc16}
\blockentry{depuncturer}{Re-insert punctured positions as erasures (LLR = 0) before decoding.}{in:BITS|LLR}{out:LLR}{\texttt{pattern}=[1, 1, 1, 0, 0, 1]}
\blockentry{hamming\_decoder}{Syndrome decoding (hard decision; LLR inputs are sliced first).}{in:BITS|LLR}{out:BITS}{\texttt{m}=3}
\blockentry{hamming\_encoder}{Systematic Hamming encoder (frame length must be a multiple of k).}{in:BITS}{out:BITS}{\texttt{m}=3}
\blockentry{interleaver}{Permute the bits (or LLRs) of every frame.}{in:BITS|LLR|SYMBOLS}{out:same as 'in'}{\texttt{kind}=random; \texttt{rows}=16; \texttt{seed}=1234; \texttt{inverse}=False}
\blockentry{ldpc\_decoder}{Belief-propagation LDPC decoder (normalized min-sum or sum-product) with early stopping.}{in:BITS|LLR}{out:BITS}{\texttt{k\_blocks}=8; \texttt{m\_blocks}=8; \texttt{z}=48; \texttt{col\_weight}=3; \texttt{seed}=1; \texttt{max\_iter}=30; \texttt{method}=min\_sum; \texttt{alpha}=0.8}
\blockentry{ldpc\_encoder}{Systematic quasi-cyclic LDPC encoder (frame length must be a multiple of k).}{in:BITS}{out:BITS}{\texttt{k\_blocks}=8; \texttt{m\_blocks}=8; \texttt{z}=48; \texttt{col\_weight}=3; \texttt{seed}=1}
\blockentry{polar\_decoder}{Successive-cancellation polar decoder (soft LLR or hard input).}{in:BITS|LLR}{out:BITS}{\texttt{block\_length}=256; \texttt{k}=128; \texttt{design\_snr\_db}=0}
\blockentry{polar\_encoder}{Polar encoder (frame length must be a multiple of k).}{in:BITS}{out:BITS}{\texttt{block\_length}=256; \texttt{k}=128; \texttt{design\_snr\_db}=0}
\blockentry{puncturer}{Remove coded bits according to a periodic pattern (raises the code rate).}{in:BITS}{out:BITS}{\texttt{pattern}=[1, 1, 1, 0, 0, 1]}
\blockentry{repetition\_decoder}{Soft (LLR sum) or hard (majority vote) repetition decoding.}{in:BITS|LLR}{out:BITS}{\texttt{n}=3}
\blockentry{repetition\_encoder}{Repeat every bit \texttt{n} times (rate 1/n).}{in:BITS}{out:BITS}{\texttt{n}=3}
\blockentry{scrambler}{XOR every frame with the same LFSR sequence.}{in:BITS|LLR}{out:same as 'in'}{\texttt{taps}=[7, 4]; \texttt{seed}=93}
\blockentry{viterbi\_decoder}{Maximum-likelihood sequence decoder; soft decisions with LLR input, hard with bits.}{in:BITS|LLR}{out:BITS}{\texttt{generators}=[171, 133]; \texttt{constraint\_length}=7; \texttt{terminate}=True}
\blockcategory{wireline}
\blockentry{copper\_line}{Cable channel with frequency-dependent loss (twisted pair, coax, PCB trace).}{in:SYMBOLS|SAMPLES}{out:same as 'in', csi:CSI}{\texttt{preset}=cat5e; \texttt{length\_m}=100; \texttt{termination\_ohm}=100; \texttt{a}=auto; \texttt{b}=auto; \texttt{max\_taps}=256}
\blockentry{crosstalk}{Near-end (NEXT) or far-end (FEXT) crosstalk from \texttt{n\_disturbers} similar pairs.}{in:SYMBOLS|SAMPLES, tx:SYMBOLS|SAMPLES}{out:same as 'in'}{\texttt{model}=next; \texttt{n\_disturbers}=10; \texttt{preset}=dsl\_0.4mm; \texttt{length\_m}=1000}
\blockentry{ctle}{Continuous-time linear equalizer: H(s) = A (1 + s/wz) / ((1 + s/wp1)(1 + s/wp2)).}{in:SYMBOLS|SAMPLES}{out:same as 'in'}{\texttt{zero\_hz}=1e+09; \texttt{pole1\_hz}=6e+09; \texttt{pole2\_hz}=2.5e+10; \texttt{dc\_gain\_db}=0}
\blockentry{dfe\_equalizer}{Adaptive feed-forward + decision-feedback equalizer (NLMS), symbol spaced.}{in:SYMBOLS, ref:SYMBOLS}{out:SYMBOLS}{\texttt{ffe\_taps}=7; \texttt{dfe\_taps}=4; \texttt{step\_size}=0.05; \texttt{n\_train}=300}
\blockentry{dmt\_demodulator}{Remove the cyclic prefix, FFT, one-tap frequency-domain equalization and per-tone demapping.}{in:SAMPLES, csi:CSI}{out:BITS}{--}
\blockentry{dmt\_modulator}{Discrete multitone (ADSL/VDSL): real IFFT with Hermitian symmetry, per-tone bit loading.}{in:BITS}{out:SAMPLES}{\texttt{n\_fft}=512; \texttt{cp\_len}=40; \texttt{first\_tone}=33; \texttt{last\_tone}=255; \texttt{sample\_rate}=2.208e+06; \texttt{loading}=waterfill; \texttt{uniform\_bits}=4; \texttt{max\_bits}=14; \texttt{gap\_db}=9.8; \texttt{margin\_db}=3; \texttt{snr\_db}=60; \texttt{line\_preset}=dsl\_0.4mm; \texttt{line\_length\_m}=2000}
\blockentry{echo\_canceller}{NLMS adaptive filter estimating the echo from the local \texttt{tx} signal and subtracting it.}{in:SYMBOLS|SAMPLES, tx:SYMBOLS|SAMPLES}{out:same as 'in'}{\texttt{n\_taps}=16; \texttt{train\_step\_size}=0.05; \texttt{step\_size}=0.002; \texttt{train\_samples}=1000}
\blockentry{hybrid\_echo}{Leakage of the local transmitter into the receiver through the hybrid (full duplex).}{in:SYMBOLS|SAMPLES, tx:SYMBOLS|SAMPLES}{out:same as 'in'}{\texttt{echo\_db}=-10; \texttt{taps}=[1.0, 0.6, 0.3, 0.1]}
\blockentry{impulse\_noise}{Bursts of strong Gaussian noise (electrical switching, appliances).}{in:SYMBOLS|SAMPLES}{out:same as 'in'}{\texttt{burst\_probability}=0.001; \texttt{burst\_samples}=32; \texttt{power\_db}=10}
\blockentry{jitter}{Sampling-clock jitter: random (RJ, per symbol) and sinusoidal (SJ), in unit intervals.}{in:SAMPLES}{out:same as 'in'}{\texttt{rj\_rms\_ui}=0.02; \texttt{sj\_amplitude\_ui}=0; \texttt{sj\_frequency\_hz}=1e+06}
\blockentry{line\_block\_decoder}{Inverse of \texttt{line\_block\_encoder}; counts invalid code words / sync headers.}{in:BITS|LLR}{out:BITS}{\texttt{code}=64b66b}
\blockentry{line\_block\_encoder}{4B/5B (FDDI, 100BASE-TX) or 64B/66B (10G/100G Ethernet) block line code.}{in:BITS}{out:BITS}{\texttt{code}=64b66b}
\blockentry{line\_decoder}{Hard decisions for NRZ / Manchester / AMI / MLT-3 symbols (the real part is used).}{in:SYMBOLS|SAMPLES}{out:BITS}{\texttt{code}=nrz; \texttt{threshold}=0.5}
\blockentry{line\_encoder}{Baseband line code: bits -> real symbols.}{in:BITS}{out:SYMBOLS}{\texttt{code}=nrz; \texttt{symbol\_rate}=1e+09}
\blockcategory{optical}
\blockentry{cd\_compensation}{Frequency-domain chromatic dispersion compensation (inverse all-pass filter).}{in:SAMPLES}{out:SAMPLES}{\texttt{dispersion\_ps\_nm}=auto; \texttt{slope\_ps\_nm2}=auto; \texttt{wavelength\_nm}=auto}
\blockentry{coherent\_receiver}{Coherent front-end: local oscillator, 90-degree hybrid and balanced photodiodes.}{in:SAMPLES}{out:SAMPLES}{\texttt{responsivity}=0.8; \texttt{thermal\_noise\_pa\_rthz}=20; \texttt{shot\_noise}=True; \texttt{lo\_power\_dbm}=13; \texttt{lo\_linewidth\_hz}=100000; \texttt{frequency\_offset\_hz}=0}
\blockentry{fiber}{Single-mode fiber span: attenuation, chromatic dispersion, optional Kerr nonlinearity and PMD.}{in:SAMPLES}{out:SAMPLES}{\texttt{attenuation\_db\_km}=0.2; \texttt{dispersion\_ps\_nm\_km}=17; \texttt{nonlinear}=False; \texttt{dispersion\_slope\_ps\_nm2\_km}=0; \texttt{gamma\_w\_km}=1.3; \texttt{steps\_per\_km}=0.5; \texttt{nonlinear\_phase\_rad}=0.05; \texttt{pmd\_ps\_sqrt\_km}=0; \texttt{pmd\_sections}=10; \texttt{polarization\_rotation}=False; \texttt{length\_km}=80}
\blockentry{fso\_channel}{Terrestrial free-space optical link from geometry and weather.}{in:SAMPLES}{out:SAMPLES, csi:CSI}{\texttt{distance\_m}=1000; \texttt{visibility\_km}=auto; \texttt{divergence\_mrad}=2; \texttt{rx\_aperture\_m}=0.08; \texttt{cn2}=1e-14; \texttt{turbulence}=gamma\_gamma; \texttt{pointing\_jitter\_m}=0}
\blockentry{fso\_turbulence}{Free-space optical scintillation: log-normal (weak) or gamma-gamma (moderate/strong).}{in:SAMPLES}{out:SAMPLES, csi:CSI}{\texttt{model}=lognormal; \texttt{log\_intensity\_variance}=0.1; \texttt{alpha}=4; \texttt{beta}=2; \texttt{rytov\_variance}=auto}
\blockentry{intensity\_modulator}{Laser with intensity modulation (directly modulated laser / EML) for IM/DD links.}{in:SYMBOLS|SAMPLES}{out:SAMPLES}{\texttt{power\_dbm}=0; \texttt{extinction\_ratio\_db}=10; \texttt{rin\_db\_hz}=auto; \texttt{chirp\_alpha}=0; \texttt{adiabatic\_chirp\_ghz\_mw}=0; \texttt{resonance\_ghz}=auto; \texttt{damping\_per\_ns}=50; \texttt{parasitic\_ghz}=auto; \texttt{linewidth\_hz}=0; \texttt{wavelength\_nm}=1550}
\blockentry{iq\_modulator}{Laser + IQ Mach-Zehnder modulator for coherent transmission (single or dual polarization).}{in:SYMBOLS|SAMPLES}{out:SAMPLES}{\texttt{power\_dbm}=0; \texttt{model}=linear; \texttt{drive\_swing}=0.5; \texttt{linewidth\_hz}=100000; \texttt{wavelength\_nm}=1550}
\blockentry{multimode\_fiber}{Graded-index multimode fiber (OM1-OM5) for short-reach IM/DD links.}{in:SAMPLES}{out:SAMPLES}{\texttt{fiber\_type}=OM4; \texttt{length\_m}=100; \texttt{model}=bandwidth; \texttt{modal\_bandwidth\_mhz\_km}=auto; \texttt{attenuation\_db\_km}=auto; \texttt{alpha}=auto; \texttt{profile\_dispersion}=0; \texttt{launch}=overfilled; \texttt{launch\_waist\_um}=4; \texttt{launch\_offset\_um}=0; \texttt{group\_weights}=[]; \texttt{core\_index}=1.467; \texttt{group\_index}=1.482; \texttt{zero\_dispersion\_nm}=1320; \texttt{dispersion\_slope}=0.101; \texttt{source\_rms\_width\_nm}=auto}
\blockentry{optical\_amplifier}{Erbium-doped fiber amplifier with amplified spontaneous emission (ASE) noise.}{in:SAMPLES}{out:SAMPLES}{\texttt{mode}=power; \texttt{gain\_db}=20; \texttt{target\_power\_dbm}=0; \texttt{noise\_figure\_db}=5}
\blockentry{optical\_link}{Multi-span amplified link: \texttt{n\_spans} x (fiber span + EDFA compensating the span loss).}{in:SAMPLES}{out:SAMPLES}{\texttt{attenuation\_db\_km}=0.2; \texttt{dispersion\_ps\_nm\_km}=17; \texttt{nonlinear}=False; \texttt{dispersion\_slope\_ps\_nm2\_km}=0; \texttt{gamma\_w\_km}=1.3; \texttt{steps\_per\_km}=0.5; \texttt{nonlinear\_phase\_rad}=0.05; \texttt{pmd\_ps\_sqrt\_km}=0; \texttt{pmd\_sections}=10; \texttt{polarization\_rotation}=False; \texttt{n\_spans}=10; \texttt{span\_km}=80; \texttt{noise\_figure\_db}=5; \texttt{launch\_power\_dbm}=auto}
\blockentry{photodiode}{Direct detection (PIN or APD): square-law current with shot and thermal noise.}{in:SAMPLES}{out:SAMPLES}{\texttt{responsivity}=0.8; \texttt{thermal\_noise\_pa\_rthz}=20; \texttt{shot\_noise}=True; \texttt{dark\_current\_na}=5; \texttt{apd\_gain}=1; \texttt{apd\_k}=0.2; \texttt{bandwidth\_ghz}=auto}
\blockentry{vcsel}{Directly modulated 850 nm VCSEL for multimode data-center links.}{in:SYMBOLS|SAMPLES}{out:SAMPLES}{\texttt{power\_dbm}=0; \texttt{extinction\_ratio\_db}=4; \texttt{resonance\_ghz}=20; \texttt{damping\_per\_ns}=100; \texttt{parasitic\_ghz}=20; \texttt{rin\_db\_hz}=-140; \texttt{wavelength\_nm}=850; \texttt{spectral\_width\_nm}=0.4}
\blockentry{wdm\_demultiplexer}{Select one WDM channel with an optical band-pass filter (rectangular or 2nd-order super-Gaussian).}{in:SAMPLES}{out:SAMPLES}{\texttt{channel}=0; \texttt{bandwidth\_ghz}=auto; \texttt{shape}=supergaussian}
\blockentry{wdm\_multiplexer}{Wavelength-division multiplexing: add neighbouring channels around the channel under test.}{in:SAMPLES}{out:SAMPLES}{\texttt{n\_channels}=5; \texttt{spacing\_ghz}=50; \texttt{interferer\_power\_db}=0}
\blockcategory{metrics}
\blockentry{ber\_meter}{Count bit, symbol and frame errors between transmitted and received bits.}{tx:BITS, rx:BITS|LLR}{--}{\texttt{bits\_per\_symbol}=auto; \texttt{skip\_bits}=0}
\blockentry{evm\_meter}{Error vector magnitude and modulation error ratio between reference and received symbols.}{ref:SYMBOLS, in:SYMBOLS}{--}{\texttt{skip\_symbols}=0}
\blockentry{optical\_power\_meter}{Average optical power (dBm) and OSNR (0.1 nm) of an optical signal.}{in:SAMPLES}{--}{--}
\blockentry{power\_meter}{Average power and peak-to-average power ratio of a waveform.}{in:SYMBOLS|SAMPLES}{--}{--}


\chapter{Metadata keys}\label{app:metadata}

Metadata travels with every signal (\cref{sec:metadata}). The table lists the keys used by the
built-in library, the blocks that set them and the blocks that read them.

\begin{center}
\small
\begin{longtable}{@{}p{3.7cm}p{4.8cm}p{5.5cm}@{}}
\toprule
\textbf{Key} & \textbf{Meaning} & \textbf{Set by / read by}\\
\midrule
\endhead
\meta{bits_per_symbol} & bits per constellation symbol $k$ & mapper, line encoder / awgn ($E_b/N_0$), BER meter\\
\meta{constellation} & the \code{Constellation} object & mapper / demapper, equalizers (decisions), PLL, BPS, AGC, WDM\\
\meta{code_rate} & cumulative code rate $R$ (with tails, CRC, puncturing) & encoders, CRC, puncturer, block line codes / awgn\\
\meta{power} & nominal mean power per sample & mapper, pulse shaper, OFDM/DMT, optical blocks / awgn (SNR reference), AGC, OSNR\\
\meta{sps} & samples per symbol $N_{\mathrm{sps}}$ & pulse shaper, rate changers, OFDM ($N/N_{\mathrm{act}}$) / awgn, matched filter, timing recovery, equalizers\\
\meta{delay} & known processing delay in samples & pulse shaper, matched filter, FIR / matched filter, timing recovery, MIMO equalizer\\
\meta{pulse_taps} & transmit pulse & pulse shaper / matched filter, WDM multiplexer\\
\meta{n_symbols} & symbols per frame & mapper, pulse shaper, preamble, MIMO / matched filter, equalizers\\
\meta{noise_var} & noise variance per complex sample or symbol (scalar or array) & awgn, equalizers, detectors, photodiode, coherent receiver / soft demapper, AGC, MIMO detector\\
\meta{spatial_rate} & streams per channel use $N_s$ & stream demux / awgn\\
\meta{csi_type} & form of CSI: \code{gain}, \code{cir}, \code{cir_tv}, \code{freq}, \code{mimo} & channels, estimators / equalizers, detectors\\
\meta{ofdm} & \code{OFDMLayout} & OFDM modulator / demodulator, estimator, equalizer, sync\\
\meta{n_data_symbols}, \meta{data_power} & payload size and data-subcarrier power & OFDM modulator / OFDM equalizer\\
\meta{preamble}, \meta{n_payload_symbols} & preamble sequence and payload size & preamble / frame sync\\
\meta{punctured_length} & mother-code length before puncturing & puncturer / depuncturer\\
\meta{line_code} & line code name & line encoder / AGC (target power)\\
\meta{dmt}, \meta{n_bits} & DMT modulator and payload size & DMT modulator / DMT demodulator\\
\meta{optical} & the signal is an optical field & optical transmitters / detectors\\
\meta{wavelength_nm} & carrier wavelength & optical transmitters / fiber, EDFA, CD and FO compensation\\
\meta{cd_ps_nm} & accumulated dispersion $D\ell$ & fiber, link / CD compensation\\
\meta{cd_compensated_ps_nm} & dispersion removed so far & CD compensation / FO compensation (bulk-delay fix)\\
\meta{ase_psd} & accumulated ASE density per polarization & EDFA, link, fiber / OSNR metrics\\
\meta{wdm} & channel count and spacing & WDM multiplexer / WDM demultiplexer\\
\bottomrule
\end{longtable}
\end{center}

\chapter{Scenario library}\label{app:scenarios}

The \numscenarios{} scenario files shipped in \file{backend/scenarios/} (subsystem files
excluded) are listed below. The table is generated from the files themselves. Open them from
the GUI scenario browser (Ctrl+O) or run them with \code{python -m linkstudio run <file>}. Each
scenario's \code{description} field explains what to observe and which parameters to sweep.

\makeatletter\let\primitiveinput\@@input\makeatother
\begin{center}
\footnotesize
\begin{longtable}{@{}>{\raggedright\arraybackslash}p{1.9cm}>{\raggedright\arraybackslash}p{5.9cm}>{\raggedright\arraybackslash}p{3.9cm}cl@{}}
\toprule
\textbf{Category} & \textbf{File} & \textbf{Name} & \textbf{Blocks} & \textbf{Mode}\\
\midrule
\endhead
\primitiveinput generated/scenarios.tex 
\bottomrule
\end{longtable}
\end{center}

Two reusable subsystems, \file{subsystems/coded_qam_tx.yaml} and
\file{subsystems/coded_qam_rx.yaml}, are used by \file{coded_link_subsystems.yaml}.


\backmatter
\begin{thebibliography}{99}
\addcontentsline{toc}{chapter}{Bibliography}

\bibitem{proakis2008} J.~G. Proakis and M.~Salehi, \emph{Digital Communications}, 5th ed. McGraw-Hill, 2008.
\bibitem{agrawal2019} G.~P. Agrawal, \emph{Nonlinear Fiber Optics}, 6th ed. Academic Press, 2019.
\bibitem{mengali1997} U.~Mengali and A.~N. D'Andrea, \emph{Synchronization Techniques for Digital Receivers}. Plenum Press, 1997.
\bibitem{cho2002} K.~Cho and D.~Yoon, ``On the general BER expression of one- and two-dimensional amplitude modulations,'' \emph{IEEE Trans. Commun.}, vol.~50, no.~7, pp.~1074--1080, 2002.
\bibitem{craig1991} J.~W. Craig, ``A new, simple and exact result for calculating the probability of error for two-dimensional signal constellations,'' in \emph{Proc. IEEE MILCOM}, 1991.
\bibitem{saleh1981} A.~A.~M. Saleh, ``Frequency-independent and frequency-dependent nonlinear models of TWT amplifiers,'' \emph{IEEE Trans. Commun.}, vol.~29, no.~11, pp.~1715--1720, 1981.
\bibitem{gardner1986} F.~M. Gardner, ``A BPSK/QPSK timing-error detector for sampled receivers,'' \emph{IEEE Trans. Commun.}, vol.~34, no.~5, pp.~423--429, 1986.
\bibitem{oerder1988} M.~Oerder and H.~Meyr, ``Digital filter and square timing recovery,'' \emph{IEEE Trans. Commun.}, vol.~36, no.~5, pp.~605--612, 1988.
\bibitem{schmidl1997} T.~M. Schmidl and D.~C. Cox, ``Robust frequency and timing synchronization for OFDM,'' \emph{IEEE Trans. Commun.}, vol.~45, no.~12, pp.~1613--1621, 1997.
\bibitem{vandebeek1997} J.-J. van de Beek, M.~Sandell, and P.~O. Börjesson, ``ML estimation of time and frequency offset in OFDM systems,'' \emph{IEEE Trans. Signal Process.}, vol.~45, no.~7, pp.~1800--1805, 1997.
\bibitem{alamouti1998} S.~M. Alamouti, ``A simple transmit diversity technique for wireless communications,'' \emph{IEEE J. Sel. Areas Commun.}, vol.~16, no.~8, pp.~1451--1458, 1998.
\bibitem{arikan2009} E.~Arıkan, ``Channel polarization: a method for constructing capacity-achieving codes for symmetric binary-input memoryless channels,'' \emph{IEEE Trans. Inf. Theory}, vol.~55, no.~7, pp.~3051--3073, 2009.
\bibitem{gallager1962} R.~G. Gallager, ``Low-density parity-check codes,'' \emph{IRE Trans. Inf. Theory}, vol.~8, no.~1, pp.~21--28, 1962.
\bibitem{starr1999} T.~Starr, J.~M. Cioffi, and P.~J. Silverman, \emph{Understanding Digital Subscriber Line Technology}. Prentice Hall, 1999.
\bibitem{godard1980} D.~N. Godard, ``Self-recovering equalization and carrier tracking in two-dimensional data communication systems,'' \emph{IEEE Trans. Commun.}, vol.~28, no.~11, pp.~1867--1875, 1980.
\bibitem{pfau2009} T.~Pfau, S.~Hoffmann, and R.~Noé, ``Hardware-efficient coherent digital receiver concept with feedforward carrier recovery for M-QAM constellations,'' \emph{J. Lightw. Technol.}, vol.~27, no.~8, pp.~989--999, 2009.
\bibitem{savory2010} S.~J. Savory, ``Digital coherent optical receivers: algorithms and subsystems,'' \emph{IEEE J. Sel. Topics Quantum Electron.}, vol.~16, no.~5, pp.~1164--1179, 2010.
\bibitem{wilson1927} E.~B. Wilson, ``Probable inference, the law of succession, and statistical inference,'' \emph{J. Amer. Statist. Assoc.}, vol.~22, no.~158, pp.~209--212, 1927.
\bibitem{clarke1968} R.~H. Clarke, ``A statistical theory of mobile-radio reception,'' \emph{Bell Syst. Tech. J.}, vol.~47, no.~6, pp.~957--1000, 1968.
\bibitem{3gpp36104} 3GPP TS 36.104, \emph{Evolved Universal Terrestrial Radio Access (E-UTRA); Base Station (BS) radio transmission and reception}, Annex B.
\bibitem{ansi417} ANSI T1.417-2001, \emph{Spectrum Management for Loop Transmission Systems}.

\bibitem{gloge1973} D.~Gloge and E.~A.~J. Marcatili, ``Multimode theory of graded-core fibers,'' \emph{Bell Syst. Tech. J.}, vol.~52, no.~9, pp.~1563--1578, 1973.
\bibitem{olshansky1976} R.~Olshansky and D.~B. Keck, ``Pulse broadening in graded-index optical fibers,'' \emph{Appl. Opt.}, vol.~15, no.~2, pp.~483--491, 1976.
\bibitem{cunningham2002} D.~Cunningham and P.~Dawe, \emph{Review of the 10 Gigabit Ethernet Link Model}, white paper, Agilent Technologies, 2002.
\bibitem{golowich2003} S.~E. Golowich, P.~Pepeljugoski \emph{et al.}, ``Modeling and simulation of next-generation multimode fiber links,'' \emph{J. Lightw. Technol.}, vol.~21, no.~5, 2003.
\bibitem{coldren2012} L.~A. Coldren, S.~W. Corzine, and M.~L. Ma\v{s}anovi\'c, \emph{Diode Lasers and Photonic Integrated Circuits}, 2nd ed. Wiley, 2012.

\bibitem{zheng2003} Y.~R. Zheng and C.~Xiao, ``Simulation models with correct statistical properties for Rayleigh fading channels,'' \emph{IEEE Trans. Commun.}, vol.~51, no.~6, pp.~920--928, 2003.
\bibitem{tr38901} 3GPP TR 38.901 v19.1.0, \emph{Study on channel model for frequencies from 0.5 to 100 GHz}, 2025.
\bibitem{p525} ITU-R Recommendation P.525-5, \emph{Calculation of free-space attenuation}, 2024.
\bibitem{g9961} ITU-T Recommendation G.996.1, \emph{Test procedures for digital subscriber line (DSL) transceivers}, 2001.
\bibitem{djordjevic2001} A.~R. Djordjevic, R.~M. Biljic, V.~D. Likar-Smiljanic, and T.~K. Sarkar, ``Wideband frequency-domain characterization of FR-4 and time-domain causality,'' \emph{IEEE Trans. Electromagn. Compat.}, vol.~43, no.~4, pp.~662--667, 2001.
\bibitem{g652} ITU-T Recommendation G.652, \emph{Characteristics of a single-mode optical fibre and cable}, 2024.
\bibitem{sinkin2003} O.~V. Sinkin, R.~Holzlöhner, J.~Zweck, and C.~R. Menyuk, ``Optimization of the split-step Fourier method in modeling optical-fiber communications systems,'' \emph{J. Lightw. Technol.}, vol.~21, no.~1, pp.~61--68, 2003.
\bibitem{marcuse1997} D.~Marcuse, C.~R. Menyuk, and P.~K.~A. Wai, ``Application of the Manakov-PMD equation to studies of signal propagation in optical fibers with randomly varying birefringence,'' \emph{J. Lightw. Technol.}, vol.~15, no.~9, pp.~1735--1746, 1997.
\bibitem{koch1984} T.~L. Koch and J.~E. Bowers, ``Nature of wavelength chirping in directly modulated semiconductor lasers,'' \emph{Electron. Lett.}, vol.~20, no.~25, pp.~1038--1040, 1984.
\bibitem{alhabash2001} M.~A. Al-Habash, L.~C. Andrews, and R.~L. Phillips, ``Mathematical model for the irradiance probability density function of a laser beam propagating through turbulent media,'' \emph{Opt. Eng.}, vol.~40, no.~8, pp.~1554--1562, 2001.
\bibitem{p1814} ITU-R Recommendation P.1814-1, \emph{Prediction methods required for the design of terrestrial free-space optical links}, 2025.
\bibitem{farid2007} A.~A. Farid and S.~Hranilovic, ``Outage capacity optimization for free-space optical links with pointing errors,'' \emph{J. Lightw. Technol.}, vol.~25, no.~7, pp.~1702--1710, 2007.

\end{thebibliography}

\end{document}
